% La diversité physique du Cameroun — Géographie Tle A — Cours signé OnBuch+
% Programme officiel MINESEC. Compiler avec : tectonic N01-diversite-physique-cameroun.tex
\newcommand{\DOCMATIERE}{Géographie}
\newcommand{\DOCNIVEAU}{Tle A}
\newcommand{\DOCMODULE}{Module 1 : Le Cameroun}
\newcommand{\DOCLECON}{N01}
\newcommand{\DOCTITRE}{La diversité physique du Cameroun}
% OnBuch+ — préambule « cours signés » ENRICHI (compilé avec tectonic / XeLaTeX)
% Système visuel « Pop » d'OnBuch+ : fond crème (jamais blanc), bordures encre
% épaisses, ombres « dures » décalées, titres Archivo Black en capitales.
% Version MATHÉMATIQUES : graphes (pgfplots/tikz), boîtes pédagogiques
% (définition, propriété, méthode, attention, le savais-tu, exercice résolu,
% exercice, TP, bilan), page de garde et signature OnBuch+ ancrée en bas.
%
% Macros attendues dans main.tex AVANT \input{preamble} :
%   \DOCMATIERE  \DOCNIVEAU  \DOCMODULE  \DOCLECON (numéro)  \DOCTITRE
\documentclass[11pt]{article}

\usepackage{fontspec}
\usepackage[french]{babel}
\usepackage[a4paper,margin=2cm,top=3.4cm,bottom=2.8cm,headheight=1.4cm,footskip=1.3cm]{geometry}
\usepackage{amsmath,amssymb,amsfonts}
\usepackage{mathtools} % \xmapsto, \coloneqq... souvent utilisés par les modèles
\usepackage{cancel} % simplifications barrées (bilans de forces)
% \abs{z} (module, valeur absolue) et \norm{u} (norme), utilisés par les modèles
\providecommand{\abs}{}\let\abs\relax\DeclarePairedDelimiter{\abs}{\lvert}{\rvert}
\providecommand{\norm}{}\let\norm\relax\DeclarePairedDelimiter{\norm}{\lVert}{\rVert}
\usepackage[table]{xcolor} % [table] : fournit aussi \rowcolors (alternance de lignes)
\usepackage{pagecolor}
\usepackage{tikz}
\usetikzlibrary{arrows.meta,positioning,calc,shapes.geometric,shapes.misc,shapes.arrows,decorations.pathmorphing,decorations.pathreplacing,decorations.markings,patterns,shadows,mindmap,trees,fit,backgrounds,matrix,intersections,angles,quotes}
\usepackage{pgfplots}
\usepackage[version=4]{mhchem} % équations nucléaires \ce{^{235}_{92}U}
\usepackage[european,siunitx]{circuitikz} % schémas électriques
\pgfplotsset{compat=1.18}
\usepgfplotslibrary{fillbetween,groupplots} % aires entre courbes (calcul intégral)
\usepackage{enumitem}
\usepackage{fancyhdr}
% [most] charge toutes les bibliothèques tcolorbox (listings, minted, raster,
% documentation...) et gonfle inutilement la mémoire TeX sur un document long
% (~300 boîtes) : on ne charge que ce qui est réellement utilisé.
\usepackage[skins,breakable]{tcolorbox}
\usepackage{graphicx}
\graphicspath{{/home/user/t-t/geo-onbuch/images/}}
\usepackage{colortbl}
\usepackage{booktabs}
\usepackage{tabularx}
\usepackage{multirow}
\usepackage{diagbox} % case d'en-tête diagonale des tableaux à double entrée
% Colonnes extensibles centrée / gauche / droite (C, L, R), souvent utilisées par les modèles
\newcolumntype{C}{>{\centering\arraybackslash}X}
\newcolumntype{L}{>{\raggedright\arraybackslash}X}
\newcolumntype{R}{>{\raggedleft\arraybackslash}X}
\usepackage{array}
\usepackage{multicol}
\usepackage{siunitx}
% parse-numbers=false : accepte aussi \SI{2,5 \times 10^{-3}}{...}
\sisetup{locale=FR,output-decimal-marker={,},group-minimum-digits=4,parse-numbers=false}
\DeclareSIUnit\ohm{\ensuremath{\symup{\Omega}}} % U+2126 absent de la police
\DeclareSIUnit\division{div} % réglages d'oscilloscope (ms/div, V/div)
\DeclareSIUnit\tr{tr} % tours (vitesse de rotation, ex. \SI{1500}{\tr\per\minute})
\DeclareSIUnit\molar{M} % concentration molaire (ex. \SI{3}{\molar}, \SI{1}{\micro\molar})
\DeclareSIUnit\an{an} % année (le modèle confond parfois avec \year, un compteur TeX) — ex. \SI{5}{\kilo\meter\per\an}
\DeclareSIUnit\FCFA{FCFA} % franc CFA (ex. \SI{500}{\FCFA\per\kilo\gram})
\DeclareSIUnit\degreeBrix{\textdegree Brix} % teneur en sucre d'un sirop/jus (réfractométrie)
\DeclareSIUnit\month{mois} % le modèle utilise parfois \month, un compteur TeX, comme unité de durée
\DeclareSIUnit\microliter{\micro\liter} % le modèle écrit parfois \microliter en un seul token
\DeclareSIUnit\million{millions} % dénombrement (ex. \SI{15}{\million\per\mL} spermatozoïdes)
\usepackage{qrcode}
% \degree : commande fréquemment utilisée par les modèles pour un angle ou une
% température en dehors de \SI{}{} (non fournie par les paquets chargés).
\providecommand{\degree}{^{\circ}}
% \qed (fin de démonstration), utilisé par les modèles sans amsthm
\providecommand{\qed}{\hfill$\square$}
% \barre : double barre des valeurs interdites dans un tableau de variations (tkz-tab)
\providecommand{\barre}{\ensuremath{\|}}
% environnement proof (amsthm non chargé) utilisé par les modèles
\ifdefined\proof\else\newenvironment{proof}[1][Démonstration]{\par\noindent\textit{#1.}\ }{\qed\par}\fi
\usepackage{lastpage}
\usepackage{hyperref}
\hypersetup{hidelinks}

% ============================================================
% Palette « Pop » OnBuch+ — mêmes hex que la classe P (plus_ui.dart)
% ============================================================
\definecolor{poporange}{HTML}{F59321}
\definecolor{popdark}{HTML}{E07A0C}
\definecolor{popcream}{HTML}{FFF6E8}
\definecolor{popink}{HTML}{1C1714}
\definecolor{poppurple}{HTML}{7A5AE0}
\definecolor{popgreen}{HTML}{2E9E6A}
\definecolor{poppink}{HTML}{E0457B}
\definecolor{popblue}{HTML}{2F7DE1}
\definecolor{popgold}{HTML}{C98A12}
\definecolor{popmuted}{HTML}{8A7F76}
\definecolor{popline}{HTML}{EADFD2}
\definecolor{popgray}{HTML}{8A7F76}
\colorlet{popgrey}{popgray}
% Alias tolérants (le modèle nomme parfois les couleurs différemment)
\colorlet{popwhite}{white}
\colorlet{popblack}{popink}
\colorlet{popred}{poppink}
\colorlet{popyellow}{popgold}
% Teintes claires (fonds de tableaux/encadrés) — le modèle nomme parfois
% les couleurs avec un suffixe L (ex. popblueL) sans les avoir définies.
\colorlet{poporangeL}{poporange!20}
\colorlet{popdarkL}{popdark!20}
\colorlet{poppurpleL}{poppurple!20}
\colorlet{popgreenL}{popgreen!20}
\colorlet{poppinkL}{poppink!20}
\colorlet{popblueL}{popblue!20}
\colorlet{popgoldL}{popgold!20}
\colorlet{popredL}{popred!20}
\colorlet{popgrayL}{popgray!20}
% Teintes claires pour fonds de tableaux / zones
\colorlet{popblueL}{popblue!10!white}
\colorlet{poporangeL}{poporange!14!white}
\colorlet{popgreenL}{popgreen!12!white}
\colorlet{poppurpleL}{poppurple!12!white}
\colorlet{poppinkL}{poppink!12!white}
\colorlet{popgoldL}{popgold!14!white}

\pagecolor{popcream}

% ============================================================
% Typographie
% ============================================================
\defaultfontfeatures{Ligatures=TeX}
\setmainfont{PlusJakartaSans-Medium.ttf}[Path=../../../fonts/,BoldFont=PlusJakartaSans-Bold.ttf]
% Maths en sans-serif (Fira Math) avec chiffres et lettres droites de Plus
% Jakarta, pour que les formules se fondent dans le texte.
\usepackage[math-style=ISO,bold-style=ISO]{unicode-math}
\setmathfont{FiraMath-Regular.otf}[Path=../../../fonts/]
\setmathfont{PlusJakartaSans-Medium.ttf}[Path=../../../fonts/,range={up/num,up/{Latin,latin},\mathup}]
\usepackage{icomma} % virgule décimale sans espace en mode math
% Caractères mathématiques Unicode absents des polices de texte (Plus Jakarta,
% Archivo Black) : rendus par leur équivalent mathématique, en texte comme en
% formule — sinon ils s'affichent en carré vide (ex. « Arithmétique dans ℤ »).
\usepackage{newunicodechar}
\newunicodechar{ℝ}{\ensuremath{\mathbb{R}}}
\newunicodechar{ℤ}{\ensuremath{\mathbb{Z}}}
\newunicodechar{ℂ}{\ensuremath{\mathbb{C}}}
\newunicodechar{ℕ}{\ensuremath{\mathbb{N}}}
\newunicodechar{ℚ}{\ensuremath{\mathbb{Q}}}
\newunicodechar{𝔻}{\ensuremath{\mathbb{D}}}
\newunicodechar{α}{\ensuremath{\alpha}}
\newunicodechar{β}{\ensuremath{\beta}}
\newunicodechar{γ}{\ensuremath{\gamma}}
\newunicodechar{δ}{\ensuremath{\delta}}
\newunicodechar{ε}{\ensuremath{\varepsilon}}
\newunicodechar{θ}{\ensuremath{\theta}}
\newunicodechar{λ}{\ensuremath{\lambda}}
\newunicodechar{μ}{\ensuremath{\mu}}
\newunicodechar{σ}{\ensuremath{\sigma}}
\newunicodechar{τ}{\ensuremath{\tau}}
\newunicodechar{φ}{\ensuremath{\varphi}}
\newunicodechar{ψ}{\ensuremath{\psi}}
\newunicodechar{ω}{\ensuremath{\omega}}
\newunicodechar{Σ}{\ensuremath{\Sigma}}
% majuscules produites par \MakeUppercase dans les titres (α-aminés -> Α)
\newunicodechar{Α}{\ensuremath{\alpha}}
\newunicodechar{Β}{\ensuremath{\beta}}
\newunicodechar{Γ}{\ensuremath{\Gamma}}
\newunicodechar{Δ}{\ensuremath{\Delta}}
\newunicodechar{Ε}{\ensuremath{\varepsilon}}
\newunicodechar{Θ}{\ensuremath{\Theta}}
\newunicodechar{Λ}{\ensuremath{\Lambda}}
\newunicodechar{Μ}{\ensuremath{\mu}}
\newunicodechar{Τ}{\ensuremath{\tau}}
\newunicodechar{Φ}{\ensuremath{\Phi}}
\newunicodechar{Ψ}{\ensuremath{\Psi}}
\newunicodechar{Ω}{\ensuremath{\Omega}}
\newunicodechar{↦}{\ensuremath{\mapsto}}
\newunicodechar{∈}{\ensuremath{\in}}
\newunicodechar{∉}{\ensuremath{\notin}}
\newunicodechar{∧}{\ensuremath{\wedge}}
\newunicodechar{⇔}{\ensuremath{\Leftrightarrow}}
\newunicodechar{⟺}{\ensuremath{\Longleftrightarrow}}
\newunicodechar{⇒}{\ensuremath{\Rightarrow}}
\newunicodechar{⟹}{\ensuremath{\Longrightarrow}}
\newunicodechar{∘}{\ensuremath{\circ}}
\newunicodechar{∪}{\ensuremath{\cup}}
\newunicodechar{∩}{\ensuremath{\cap}}
\newunicodechar{≡}{\ensuremath{\equiv}}
\newunicodechar{⊂}{\ensuremath{\subset}}
\newunicodechar{⊥}{\ensuremath{\perp}}
\newunicodechar{∃}{\ensuremath{\exists}}
\newunicodechar{∀}{\ensuremath{\forall}}
\newunicodechar{∛}{\ensuremath{\sqrt[3]{\,}}}
\newunicodechar{∜}{\ensuremath{\sqrt[4]{\,}}}
\newunicodechar{⃗}{\ensuremath{{}^{\to}}}
\newunicodechar{ⁿ}{\ifmmode^{n}\else\textsuperscript{\ensuremath{n}}\fi}
\newunicodechar{ˣ}{\ifmmode^{x}\else\textsuperscript{\ensuremath{x}}\fi}
\newunicodechar{ᵇ}{\ifmmode^{b}\else\textsuperscript{\ensuremath{b}}\fi}
\newunicodechar{ʳ}{\ifmmode^{r}\else\textsuperscript{\ensuremath{r}}\fi}
\newunicodechar{ᵘ}{\ifmmode^{u}\else\textsuperscript{\ensuremath{u}}\fi}
\newunicodechar{ʸ}{\ifmmode^{y}\else\textsuperscript{\ensuremath{y}}\fi}
\newunicodechar{ᵃ}{\ifmmode^{a}\else\textsuperscript{\ensuremath{a}}\fi}
\newunicodechar{ᵉ}{\ifmmode^{e}\else\textsuperscript{\ensuremath{e}}\fi}
\newunicodechar{ᵏ}{\ifmmode^{k}\else\textsuperscript{\ensuremath{k}}\fi}
\newunicodechar{ᵖ}{\ifmmode^{p}\else\textsuperscript{\ensuremath{p}}\fi}
\newunicodechar{ᴺ}{\ifmmode^{N}\else\textsuperscript{\ensuremath{N}}\fi}
\newunicodechar{ᶜ}{\ifmmode^{c}\else\textsuperscript{\ensuremath{c}}\fi}
\newunicodechar{ʰ}{\ifmmode^{h}\else\textsuperscript{\ensuremath{h}}\fi}
\newunicodechar{ᵠ}{\ifmmode^{q}\else\textsuperscript{\ensuremath{q}}\fi}
\newunicodechar{ₙ}{\ifmmode_{n}\else\textsubscript{\ensuremath{n}}\fi}
\newunicodechar{ₐ}{\ifmmode_{a}\else\textsubscript{\ensuremath{a}}\fi}
\newunicodechar{ᵢ}{\ifmmode_{i}\else\textsubscript{\ensuremath{i}}\fi}
\newunicodechar{ᵣ}{\ifmmode_{r}\else\textsubscript{\ensuremath{r}}\fi}
\newunicodechar{ₖ}{\ifmmode_{k}\else\textsubscript{\ensuremath{k}}\fi}
\newunicodechar{ᵦ}{\ifmmode_{\beta}\else\textsubscript{\ensuremath{\beta}}\fi}
\newunicodechar{ₓ}{\ifmmode_{x}\else\textsubscript{\ensuremath{x}}\fi}
\newfontfamily\popdisplay{ArchivoBlack-Regular.ttf}[Path=../../../fonts/]
\newfontfamily\poplogo{SpaceGrotesk-Bold.ttf}[Path=../../../fonts/]
\newfontfamily\popsemi{PlusJakartaSans-SemiBold.ttf}[Path=../../../fonts/]
\newfontfamily\popxbold{PlusJakartaSans-ExtraBold.ttf}[Path=../../../fonts/]
\newfontfamily\popxboldls{PlusJakartaSans-ExtraBold.ttf}[Path=../../../fonts/,LetterSpace=3.0]

\setlist[enumerate]{leftmargin=1.7em,itemsep=5pt,topsep=4pt}
\setlist[itemize]{leftmargin=1.7em,itemsep=4pt,topsep=4pt,label={\color{poporange}\textbullet}}
\setlength{\parindent}{0pt}
\setlength{\parskip}{5pt}
\renewcommand{\arraystretch}{1.25}

% pgfplots : style Pop par défaut
\pgfplotsset{
  every axis/.append style={axis line style={popink,line width=1pt},tick style={popink},
    grid style={popline,line width=0.5pt},label style={font=\small},tick label style={font=\footnotesize}},
  popcurve/.style={poporange,line width=1.8pt,no markers,smooth},
}
% Même style disponible dans un \draw TikZ simple (hors pgfplots)
\tikzset{popcurve/.style={poporange,line width=1.8pt}}
\tikzset{popgrid/.style={popline,very thin}}
\colorlet{popcurve}{poporange} % le nom du style est parfois utilisé comme couleur

% ============================================================
% Pill / squiggle
% ============================================================
\newcommand{\poppill}[3]{%
  \tikz[baseline=-0.55ex]\node[draw=popink,line width=1.2pt,rounded corners=2.6pt,%
    fill=#2,inner xsep=6.5pt,inner ysep=3pt]{%
    \popxboldls\fontsize{7}{7}\selectfont\color{#3}\MakeUppercase{#1}};%
}
\newcommand{\popsquiggle}[1]{%
  \tikz[baseline=-0.4ex]{
    \draw[#1,line width=2.6pt,line cap=round]
      plot [smooth] coordinates {(0,0.09) (0.34,0.01) (0.68,0.21) (1.02,0.01) (1.36,0.10)};
  }%
}
% Mot-clé surligné (marqueur orange sous le texte)
\newcommand{\cle}[1]{{\popxbold\color{popdark}#1}}

% ============================================================
% En-tête / pied
% ============================================================
\pagestyle{fancy}
\fancyhf{}
\renewcommand{\headrulewidth}{0pt}
\renewcommand{\footrulewidth}{0pt}
\lhead{\raisebox{-0.15ex}{\poplogo\fontsize{13}{13}\selectfont\color{popink}OnBuch\color{poporange}+}}
\rhead{\poppill{\DOCMATIERE\ \textbullet\ \DOCNIVEAU\ \textbullet\ Leçon \DOCLECON}{white}{popink}}
\lfoot{\popxbold\fontsize{7.6}{7.6}\selectfont\color{popmuted}\MakeUppercase{Cours signé OnBuch+}\ \textbullet\ {\popsemi\DOCTITRE}}
\rfoot{\tikz[baseline=-0.6ex]\node[rounded corners=8.5pt,fill=popink,minimum height=17pt,minimum width=17pt,inner xsep=5pt,inner sep=0]{\popxbold\fontsize{8}{8}\selectfont\color{popcream}\thepage\,/\,\pageref*{LastPage}};}

% ============================================================
% Titres
% ============================================================
\newcommand{\chaptitre}[2]{%
  \begin{tcolorbox}[enhanced,colback=poporange,colframe=popink,boxrule=2.6pt,arc=11pt,
      drop shadow=popink,left=15pt,right=15pt,top=12pt,bottom=11pt,before skip=3pt,after skip=18pt]
    \poppill{#2}{popink}{popcream}\par\vspace{9pt}
    {\raggedright\hyphenpenalty=10000\exhyphenpenalty=10000\popdisplay\fontsize{22}{24}\selectfont\color{popcream}\MakeUppercase{#1}\par}
  \end{tcolorbox}
}
% Section numérotée : pastille orange avec le numéro + titre Archivo
\newcounter{popsec}
\newcommand{\coursec}[1]{%
  \stepcounter{popsec}\setcounter{popsub}{0}%
  \par\vspace{16pt}\noindent
  \begin{minipage}{\linewidth}
  \tikz[baseline=-0.7ex]\node[draw=popink,line width=1.6pt,fill=poporange,rounded corners=4pt,minimum size=22pt,inner sep=0]{\popdisplay\fontsize{12}{12}\selectfont\color{popcream}\thepopsec};%
  \hspace{8pt}{\popdisplay\fontsize{15}{17}\selectfont\color{popink}\MakeUppercase{#1}}\par
  \vspace{4pt}\noindent{\color{poporange}\rule{36pt}{3.6pt}}
  \end{minipage}\par\nopagebreak\vspace{8pt}
}
\newcounter{popsub}
\newcommand{\courssub}[1]{%
  \stepcounter{popsub}%
  \par\vspace{9pt}\noindent{\popxbold\fontsize{11.5}{13}\selectfont\color{popink}{\color{poporange}\thepopsec.\thepopsub}\hspace{6pt}#1}\par\nopagebreak\vspace{4pt}
}

% ============================================================
% Boîtes pédagogiques — même recette : fond blanc, bordure épaisse colorée,
% ombre dure encre, badge plein. Titre optionnel [#1].
% ============================================================
\tcbset{popbase/.style={breakable,enhanced,colback=white,boxrule=2.3pt,arc=9pt,
  shadow={3pt}{-3pt}{0pt}{popink},left=14pt,right=14pt,top=11pt,bottom=11pt,before skip=11pt,after skip=15pt,
  fontupper=\popsemi\fontsize{10.6}{14.5}\selectfont\color{popink},
  fontlower=\popsemi\fontsize{10.6}{14.5}\selectfont\color{popink}}}
% Coche dessinée (le glyphe ✓ n'existe pas dans les polices du document)
\newcommand{\popcheck}[1]{\tikz[baseline=-0.5ex]\draw[#1,line width=1.6pt,line cap=round,line join=round](-0.13,0.0)--(-0.04,-0.09)--(0.13,0.1);}
\providecommand{\coche}{\popcheck{popgreen}}
\providecommand{\resultat}[1]{\par\smallskip\noindent\fcolorbox{popgreen}{popgreenL}{\parbox{\dimexpr\linewidth-2\fboxsep-2\fboxrule\relax}{\textbf{Résultat.} #1}}\par\smallskip}
\newcommand{\popboxhead}[3]{\poppill{#1}{#2}{white}\ifx\relax#3\relax\else\hspace{6pt}{\popxbold\fontsize{10.6}{12}\selectfont\color{popink}#3}\fi\par\vspace{7pt}}

\newtcolorbox{aretenir}[1][]{popbase,colframe=popgreen,before upper={\popboxhead{À retenir}{popgreen}{#1}}}
\newtcolorbox{exemplebox}[1][]{popbase,colframe=poporange,before upper={\popboxhead{Exemple}{poporange}{#1}}}
\newtcolorbox{definition}[1][]{popbase,colframe=popblue,before upper={\popboxhead{Définition}{popblue}{#1}}}
\newtcolorbox{propriete}[1][]{popbase,colframe=poppurple,before upper={\popboxhead{Propriété}{poppurple}{#1}}}
\newtcolorbox{methode}[1][]{popbase,colframe=popgold,before upper={\popboxhead{Méthode}{popgold}{#1}}}
\newtcolorbox{attention}[1][]{popbase,colframe=poppink,before upper={\popboxhead{Attention}{poppink}{#1}}}
\newtcolorbox{experience}[1][]{popbase,colframe=popblue,colback=popblueL,before upper={\popboxhead{Expérience}{popblue}{#1}}}
% « Le savais-tu ? » : carte violette pleine (contexte, culture, Cameroun)
\newtcolorbox{savaistu}[1][]{popbase,colback=poppurple,colframe=popink,
  fontupper=\popsemi\fontsize{10.4}{14.2}\selectfont\color{white},
  before upper={\poppill{Le savais-tu\,?}{white}{poppurple}\ifx\relax#1\relax\else\hspace{6pt}{\popxbold\color{white}#1}\fi\par\vspace{7pt}}}
% Exercice résolu : énoncé en haut, \tcblower puis la solution sur fond crème
\newcounter{popexo}\newcounter{popexr}
\newtcolorbox{exoresolu}[1][]{popbase,colframe=poporange,colbacklower=popcream,
  segmentation style={popink,line width=1.2pt,dashed},
  before upper={\stepcounter{popexr}\popboxhead{Exercice résolu \thepopexr}{poporange}{#1}},
  before lower={\poppill{Solution}{popink}{popcream}\par\vspace{6pt}}}
% Exercice (non résolu en place) — la correction est dans la partie Corrigés
\newtcolorbox{exercice}[1][]{popbase,colframe=popink,
  before upper={\stepcounter{popexo}\popboxhead{Exercice \thepopexo}{popink}{#1}}}
% Corrigé
\newtcolorbox{corrige}[1][]{popbase,colframe=popgreen,colback=popgreenL,
  before upper={\popboxhead{Corrigé}{popgreen}{#1}}}
% Objectifs / prérequis / bilan
\newtcolorbox{objectifs}{popbase,colframe=popink,colback=poporangeL,
  before upper={\popboxhead{Objectifs de la leçon}{popink}{}}}
\newtcolorbox{prerequis}{popbase,colframe=popink,colback=popblueL,
  before upper={\popboxhead{Prérequis}{popblue}{}}}
\newtcolorbox{bilan}{popbase,colframe=popink,colback=white,boxrule=2.8pt,
  before upper={\popboxhead{Fiche bilan}{poporange}{L'essentiel à connaître pour l'examen}}}
% Figure encadrée avec légende
\newtcolorbox{popfigure}[1][]{enhanced,breakable=false,colback=white,colframe=popline,boxrule=1.4pt,arc=8pt,
  left=10pt,right=10pt,top=10pt,bottom=8pt,before skip=10pt,after skip=12pt,halign=center,
  lower separated=false,halign lower=center,
  fontlower=\popsemi\fontsize{9}{11.5}\selectfont\color{popmuted},
  #1}
\newcommand{\legende}[1]{\par\vspace{6pt}{\popsemi\fontsize{9}{11.5}\selectfont\color{popmuted}#1}}

% ============================================================
% Signature OnBuch+
% ============================================================
\newcommand{\poplogoinline}[1]{{\poplogo\fontsize{#1}{#1}\selectfont\color{popink}OnBuch\color{poporange}+}}
\newcommand{\signatureonbuch}{%
  \begin{tcolorbox}[enhanced,colback=popink,colframe=popink,boxrule=2.2pt,arc=10pt,
      drop shadow=poporange,left=14pt,right=14pt,top=10pt,bottom=10pt,before skip=0pt,after skip=0pt]
    \tikz[baseline=-0.7ex]\node[circle,fill=popgreen,draw=popcream,line width=1.3pt,minimum size=22pt,inner sep=0]{\popcheck{white}};%
    \hspace{9pt}\begin{minipage}[c]{0.62\linewidth}
      {\popxbold\fontsize{12}{13}\selectfont\color{popcream}Signé\ \textbullet\ {\poplogo OnBuch\color{poporange}+}}\par\vspace{2pt}
      {\raggedright\popsemi\fontsize{8}{10.5}\selectfont\color{popline}\MakeUppercase{Équipe pédagogique OnBuch+}\\\MakeUppercase{Conforme au programme officiel MINESEC}\par}
    \end{minipage}\hfill
    {\poplogo\fontsize{22}{22}\selectfont\color{popcream}OnBuch\color{poporange}+}
  \end{tcolorbox}%
}

% ============================================================
% Page de garde
% #1 titre · #2 matière · #3 niveau · #4 module · #5 n° leçon · #6 durée ·
% #7 description / objectifs (texte ou itemize)
% ============================================================
\newcommand{\pagegarde}[7]{%
  \thispagestyle{empty}
  \newgeometry{a4paper,margin=1.7cm,top=1.6cm,bottom=1.4cm}%
  \begin{tikzpicture}[remember picture,overlay]
    \fill[poporange!18] ([xshift=-3.2cm,yshift=-3.6cm]current page.north east) circle (4.6cm);
    \fill[poppurple!14] ([xshift=2.4cm,yshift=7.6cm]current page.south west) circle (3.8cm);
  \end{tikzpicture}%
  {\poplogo\fontsize{30}{30}\selectfont\color{popink}OnBuch\color{poporange}+}\hfill
  \raisebox{0.6ex}{\poppill{Cours signé \textbullet\ Premium}{popink}{popcream}}\par
  \vspace{1pt}\popsquiggle{poppurple}\par
  \vspace{8pt}
  \begin{tcolorbox}[enhanced,colback=poporange,colframe=popink,boxrule=2.8pt,arc=13pt,
      drop shadow=popink,left=18pt,right=18pt,top=12pt,bottom=13pt,after skip=10pt]
    \poppill{#2 \textbullet\ #3}{popink}{popcream}\hspace{5pt}\poppill{#4}{white}{popink}\par\vspace{8pt}
    {\popdisplay\fontsize{14}{14}\selectfont\color{popink}LEÇON #5}\par\vspace{4pt}
    {\raggedright\hyphenpenalty=10000\exhyphenpenalty=10000\popdisplay\fontsize{25}{27}\selectfont\color{popcream}\MakeUppercase{#1}\par}
    \vspace{8pt}
    {\popxbold\fontsize{9.5}{9.5}\selectfont\color{popink}\MakeUppercase{Durée conseillée : #6}}
  \end{tcolorbox}
  \begin{tcolorbox}[enhanced,colback=white,colframe=popink,boxrule=2.2pt,arc=10pt,
      drop shadow=popink,left=16pt,right=16pt,top=10pt,bottom=10pt,after skip=8pt]
    \poppill{Dans cette leçon}{poppurple}{white}\par\vspace{5pt}
    \popsemi\fontsize{9.3}{12.6}\selectfont\color{popink}#7
  \end{tcolorbox}
  \noindent\begin{minipage}[t]{0.487\linewidth}
    \begin{tcolorbox}[enhanced,colback=popgreen,colframe=popink,boxrule=2.2pt,arc=10pt,
        drop shadow=popink,left=11pt,right=11pt,top=10pt,bottom=10pt,height=3.1cm,valign=center]
      \tcbox[colback=white,colframe=popink,boxrule=1.3pt,arc=5pt,boxsep=0pt,nobeforeafter,
        left=3pt,right=3pt,top=3pt,bottom=3pt,valign=center]{\qrcode[height=1.9cm]{https://play.google.com/store/apps/details?id=com.luvvix.onbuch&hl=fr}}%
      \hspace{8pt}\begin{minipage}[c]{0.5\linewidth}
        {\popdisplay\fontsize{11}{12.5}\selectfont\color{white}\MakeUppercase{Télécharge OnBuch}}\par\vspace{4pt}
        {\raggedright\popsemi\fontsize{8.4}{11}\selectfont\color{white}Gratuit \textbullet\ Google Play\\Tuteur IA, annales, concours, résultats\par}
      \end{minipage}
    \end{tcolorbox}
  \end{minipage}\hfill
  \begin{minipage}[t]{0.487\linewidth}
    \begin{tcolorbox}[enhanced,colback=poppurple,colframe=popink,boxrule=2.2pt,arc=10pt,
        drop shadow=popink,left=11pt,right=11pt,top=10pt,bottom=10pt,height=3.1cm,valign=center]
      \tikz[baseline=-0.7ex]\node[circle,fill=white,draw=popink,line width=1.3pt,minimum size=1.6cm,inner sep=0]{\popdisplay\fontsize{24}{24}\selectfont\color{poppurple}+};%
      \hspace{8pt}\begin{minipage}[c]{0.6\linewidth}
        {\popdisplay\fontsize{11}{12.5}\selectfont\color{white}\MakeUppercase{Passe à OnBuch+}}\par\vspace{4pt}
        {\raggedright\popsemi\fontsize{8.4}{11}\selectfont\color{white}Tous les cours signés, Tuteur IA illimité, fiches PDF — onglet OnBuch+ de l'app\par}
      \end{minipage}
    \end{tcolorbox}
  \end{minipage}
  \vspace{8pt}
  \popcontenu
  \vfill
  \signatureonbuch
  \restoregeometry
  \newpage
}

% Rangée « Ce cours contient » (page de garde)
\newcommand{\poptile}[3]{% #1 couleur #2 gros texte #3 légende
  \begin{tcolorbox}[enhanced,colback=#1,colframe=popink,boxrule=2pt,arc=9pt,shadow={2.5pt}{-2.5pt}{0pt}{popink},
      left=6pt,right=6pt,top=9pt,bottom=9pt,halign=center,valign=center,height=1.75cm,width=\linewidth,nobeforeafter]
    {\popdisplay\fontsize{12.5}{13}\selectfont\color{white}\MakeUppercase{#2}}\par\vspace{3pt}
    {\centering\popsemi\fontsize{7.8}{9.5}\selectfont\color{white}#3\par}
  \end{tcolorbox}}
\newcommand{\popcontenu}{%
  \par\noindent{\popxboldls\fontsize{7.5}{7.5}\selectfont\color{popmuted}CE COURS SIGNÉ CONTIENT}\par\vspace{7pt}
  \noindent\begin{minipage}[t]{0.235\linewidth}\poptile{popblue}{Cours}{complet, illustré, pas à pas}\end{minipage}\hfill
  \begin{minipage}[t]{0.235\linewidth}\poptile{poporange}{Méthodes}{et exercices résolus}\end{minipage}\hfill
  \begin{minipage}[t]{0.235\linewidth}\poptile{poppink}{Exercices}{type examen + corrigés}\end{minipage}\hfill
  \begin{minipage}[t]{0.235\linewidth}\poptile{popgreen}{Fiche bilan}{l'essentiel à retenir}\end{minipage}\par}

% Sommaire
\newtcolorbox{sommaire}{enhanced,colback=white,colframe=popink,boxrule=2.2pt,arc=10pt,
  drop shadow=popink,left=18pt,right=18pt,top=13pt,bottom=13pt,before skip=6pt,after skip=20pt,
  before upper={\poppill{Sommaire}{poppurple}{white}\par\vspace{9pt}\popsemi\fontsize{10.6}{14.5}\selectfont\color{popink}}}

% Signature de fin de cours — poussée tout en bas de la dernière page
\newcommand{\findecours}{%
  \par\vfill\nopagebreak
  \begin{center}\popsquiggle{poporange}\end{center}\vspace{4pt}
  \signatureonbuch
}
\AtBeginDocument{\def\llbracket{\mathopen{[\mkern-3.5mu[}}\def\rrbracket{\mathclose{]\mkern-3.5mu]}}} % intervalles d'entiers (glyphe ⟦ absent de la police)
% Glyphes absents de Fira Math : redéfinis à partir de glyphes disponibles
\AtBeginDocument{%
  \def\setminus{\mathbin{\backslash}}\let\smallsetminus\setminus
  \def\vdots{\mathord{\vbox{\baselineskip3.2pt\lineskiplimit0pt\kern4pt\hbox{.}\hbox{.}\hbox{.}}}}%
  \def\ddots{\mathinner{\mkern1mu\raise6.5pt\hbox{.}\mkern2mu\raise3.6pt\hbox{.}\mkern2mu\raise0.7pt\hbox{.}\mkern1mu}}%
  \def\triangle{\mathord{\scalebox{0.9}{$\symup{\Delta}$}}}%
  \def\nabla{\mathord{\raisebox{\depth}{\scalebox{1}[-1]{$\symup{\Delta}$}}}}%
  \def\checkmark{\popcheck{popdark}}%
  \def\star{\mathbin{\ast}}\def\bigstar{\mathord{\ast}}%
  \def\diamond{\mathbin{\scalebox{0.6}{$\Diamond$}}}%
  \def\bigcup{\mathop{\mathchoice{\raisebox{-0.3em}{\scalebox{1.7}{$\cup$}}}{\scalebox{1.25}{$\cup$}}{\cup}{\cup}}\displaylimits}%
  \def\bigcap{\mathop{\mathchoice{\raisebox{-0.3em}{\scalebox{1.7}{$\cap$}}}{\scalebox{1.25}{$\cap$}}{\cap}{\cap}}\displaylimits}%
  \def\lesseqqgtr{\mathrel{\lessgtr}}\def\bigoplus{\mathop{\oplus}\displaylimits}%
}
\providecommand{\ssi}{si et seulement si} % abréviation fréquente
\AtBeginDocument{\providecommand{\wideparen}{\overparen}} % arc de cercle

%% FIGFIX-BEGIN v3 — correctifs globaux des figures (docs/onbuch-generation/tools/figfix)
\usepackage{adjustbox}\usepackage{etoolbox}
% toute figure TikZ/pgfplots plus large que la ligne est réduite à la largeur disponible
\BeforeBeginEnvironment{tikzpicture}{\begin{adjustbox}{max width=\linewidth}}
\AfterEndEnvironment{tikzpicture}{\end{adjustbox}}
% légendes pgfplots lisibles par-dessus les courbes
\pgfplotsset{every axis/.append style={legend style={fill=white,fill opacity=0.92,text opacity=1,draw=black!15,inner sep=3pt}}}
% étiquettes d'arêtes (syntaxe edge["..."]) avec fond blanc
\tikzset{every edge quotes/.append style={fill=white,inner sep=1.2pt}}
% bulles de cartes mentales : texte à la ligne dans la bulle, bulle un peu plus grande
\tikzset{every concept/.append style={text width=2.0cm,align=center,minimum size=2.3cm,inner sep=2pt}}
%% FIGFIX-END
\begin{document}

\pagegarde{La diversité physique du Cameroun}{Géographie}{Tle A}{Module 1 : Le Cameroun}{N01}{2 h + TP 1 (2 h)}{Cette leçon étudie les grandes composantes physiques du Cameroun (relief, climat, végétation, hydrographie, milieu biogéographique) et leurs interactions. À partir de cartes, photographies et données chiffrées, elle montre comment cette diversité physique structure l'espace national et constitue un enjeu majeur pour l'intégration nationale.
\vspace{4pt}
\begin{multicols}{2}\small
\begin{itemize}
\item À la fin de cette leçon, je sais localiser et décrire les grandes unités de relief du Cameroun sur une carte.
\item À la fin de cette leçon, je sais identifier les domaines climatiques du Cameroun et interpréter des données climatiques (températures, précipitations).
\item À la fin de cette leçon, je sais décrire les grands types de végétation et les relier aux climats.
\item À la fin de cette leçon, je sais présenter le réseau hydrographique camerounais et ses principaux bassins versants.
\item À la fin de cette leçon, je sais caractériser les milieux biogéographiques du Cameroun (faune, sols, écosystèmes).
\item À la fin de cette leçon, je sais légender correctement une carte et construire un croquis géographique.
\end{itemize}\end{multicols}}

\begin{sommaire}
\begin{multicols}{2}
\begin{enumerate}
\item Le relief du Cameroun
\item Les climats du Cameroun
\item La végétation du Cameroun
\item L'hydrographie du Cameroun
\item Le milieu biogéographique
\item Diversité physique et intégration nationale (synthèse)
\item Activité d'intégration
\item Méthodes et exercices résolus
\item Exercices
\item Corrigés des exercices
\item Fiche bilan
\end{enumerate}
\end{multicols}
\end{sommaire}

\begin{objectifs}
\begin{itemize}
\item À la fin de cette leçon, je sais localiser et décrire les grandes unités de relief du Cameroun sur une carte.
\item À la fin de cette leçon, je sais identifier les domaines climatiques du Cameroun et interpréter des données climatiques (températures, précipitations).
\item À la fin de cette leçon, je sais décrire les grands types de végétation et les relier aux climats.
\item À la fin de cette leçon, je sais présenter le réseau hydrographique camerounais et ses principaux bassins versants.
\item À la fin de cette leçon, je sais caractériser les milieux biogéographiques du Cameroun (faune, sols, écosystèmes).
\item À la fin de cette leçon, je sais légender correctement une carte et construire un croquis géographique.
\item À la fin de cette leçon, je sais construire un diagramme climatique à partir d'un tableau statistique.
\item À la fin de cette leçon, je sais construire une carte physique du Cameroun sur un trait donné (relief, climat, hydrographie ou végétation).
\item À la fin de cette leçon, je sais expliquer en quoi la diversité physique est un enjeu pour l'intégration nationale.
\end{itemize}
\end{objectifs}

\begin{prerequis}
\begin{itemize}
\item Notions de base sur les coordonnées géographiques (latitude, longitude) et la localisation du Cameroun en Afrique centrale.
\item Lecture simple d'une carte : échelle, légende, orientation (vues en classe de Première).
\item Notions générales sur les climats tropicaux (saisons sèche et pluvieuse).
\item Connaissance des grandes régions administratives du Cameroun et de leurs chefs-lieux.
\item Notion de bassin versant et de cours d'eau (affluent, confluent, embouchure).
\end{itemize}
\end{prerequis}

\newpage

\coursec{Le relief du Cameroun}

\textbf{Situation d'accroche.} Un commerçant de Douala charge son camion de marchandises et prend la route de Maroua. Premier jour : il roule sous la canopée de la forêt dense du Sud, sur un plateau ondulé où la pluie ne semble jamais s'arrêter. Deuxième jour : il gravit le plateau de l'Adamaoua, les arbres se raréfient, la savane herbeuse s'étend à perte de vue, l'air devient plus frais. Troisième jour : il dévale vers les plaines sèches du Grand Nord, où les buttes rocheuses isolées ponctuent un horizon plat, et où il lui faut parfois traverser le Logone en pirogue. À son retour, il affirme avoir « visité plusieurs pays ». Comment un seul territoire national peut-il offrir des paysages aussi contrastés ? Cette diversité physique, qui vaut au Cameroun le surnom d'\cle{« Afrique en miniature »}, est-elle un atout ou un obstacle à l'\cle{intégration nationale} ? Telle est la problématique qui guidera toute cette leçon, et dont le relief constitue le premier élément de réponse.

\textbf{Annonce du plan de la section.} Après avoir situé le Cameroun comme un pays-charnière, nous décrirons les grandes unités de relief, du Sud au Nord : les plaines et bas plateaux du Sud, les hautes terres de l'Ouest, le plateau de l'Adamaoua, les plaines du Nord et les reliefs résiduels du Grand Nord. Nous terminerons par les conséquences humaines de cette organisation du relief.

\courssub{Un pays-charnière entre deux Africaines}

Le Cameroun occupe une position géographique exceptionnelle : il s'étire sur plus de \SI{1200}{km} du sud au nord, entre le \cle{golfe de Guinée} (océan Atlantique) au sud et la \cle{cuvette du lac Tchad} au nord, entre environ $2^\circ$ et $13^\circ$ de latitude nord. Il est à la fois un pays d'\cle{Afrique de l'Ouest} (il partage une longue frontière avec le Nigeria) et d'\cle{Afrique centrale} (il est voisin du Tchad, de la Centrafrique, du Congo, du Gabon et de la Guinée équatoriale). Cette position de charnière explique que le pays concentre, sur environ \SI{475000}{\kilo\meter\squared}, presque toutes les formes de relief, tous les climats et tous les paysages du continent : d'où le surnom d'« Afrique en miniature ».

\begin{definition}[Unité de relief]
Une \cle{unité de relief} est un ensemble de formes de terrain présentant des caractères communs : l'\cle{altitude}, la \cle{structure} (nature et disposition des roches) et le \cle{modelé} (formes sculptées par l'érosion ou construites par le volcanisme). On distingue classiquement les plaines (faible altitude, faibles pentes), les plateaux (surfaces élevées et tabulaires) et les montagnes (fortes altitudes, forts dénivelés).
\end{definition}

Le relief camerounais s'organise en \textbf{cinq grandes unités} disposées en bandes sensiblement sud-nord, avec une charnière centrale : le plateau de l'Adamaoua. L'altitude moyenne du pays est voisine de \SI{650}{m}, mais les contrastes sont extrêmes : de 0 m sur le littoral à \SI{4095}{m} au sommet du mont Cameroun.

\begin{methode}[Décrire un relief]
Pour décrire le relief d'un espace, procède toujours \textbf{du général au particulier} :
\begin{enumerate}
\item \textbf{Identifier les grandes unités} de relief et les localiser (points cardinaux, villes repères) ;
\item \textbf{Décrire les formes remarquables} de chaque unité (sommets, plaines, vallées, buttes) ;
\item \textbf{Citer les altitudes repères} (point culminant, altitudes moyennes) ;
\item \textbf{S'appuyer systématiquement sur la carte} : chaque affirmation doit pouvoir être vérifiée par la légende et les courbes de niveau ou les teintes hypsométriques.
\end{enumerate}
Au Baccalauréat, une description de relief sans altitude chiffrée ni localisation précise est une copie à moitié réussie.
\end{methode}

\courssub{Les plaines et bas plateaux du Sud (0 à 650 m)}

Le Sud du Cameroun est un domaine de \cle{reliefs bas}, adossé à l'océan Atlantique. On y distingue deux ensembles.

\textbf{La plaine côtière.} Étroite bande de 0 à \SI{200}{m} d'altitude, elle borde le golfe de Guinée entre Rio del Rey (frontière nigériane) et Campo. C'est une plaine de \cle{sédimentation récente}, basse, marécageuse, bordée de mangroves et d'estuaires (Wouri, Sanaga, Ntem). Elle accueille les deux grandes villes portuaires du pays : \cle{Douala}, capitale économique, et \cle{Kribi}, où un port en eau profonde est entré en service en 2014 (phase opérationnelle complète en 2018). Le relief plat y a facilité l'urbanisation et l'industrialisation, mais pose des problèmes d'inondation et d'avancée de l'érosion marine.

\textbf{Le plateau sud-camerounais.} Au-delà de la plaine, le terrain s'élève doucement vers un vaste \cle{plateau} compris entre 300 et \SI{800}{m} d'altitude (on retient la fourchette indicative 0 à \SI{650}{m} pour l'ensemble du domaine sud). Il s'étend autour de \cle{Yaoundé} (environ \SI{750}{m}), \cle{Ebolowa}, \cle{Sangmélima} et \cle{Bertoua}. C'est un vieux socle cristallin (roches anciennes du Précambrien) profondément altéré sous climat équatorial humide : le modelé y est fait de \cle{demi-oranges} (collines convexes) séparées par des vallées marécageuses. Ce plateau porte la grande forêt équatoriale et constitue le grenier vivrier du pays (cacao, café, cultures vivrières).

\courssub{Les hautes terres de l'Ouest : le toit du pays}

À l'ouest, le relief se dresse brutalement : c'est le domaine des \cle{hautes terres}, lié à la \cle{ligne volcanique du Cameroun}.

\textbf{Les massifs montagneux de l'Ouest.} Dans les régions de l'Ouest et du Nord-Ouest, d'anciens massifs volcaniques et cristallins culminent entre \num{2000} et \SI{3000}{m} : les \cle{monts Bamboutos} (environ \SI{2740}{m}) et le \cle{mont Oku} (\SI{3011}{m}), deuxième sommet du pays. Ces montagnes, aux pentes raides et aux sols volcaniques fertiles, sont intensément cultivées (maraîchage, pommes de terre, café arabica) et densément peuplées.

\textbf{La ligne volcanique du Cameroun et le mont Cameroun.} La \cle{ligne volcanique du Cameroun} est un alignement de volcans et de massifs d'orientation sud-ouest/nord-est, qui court du golfe de Guinée (île de Bioko) jusqu'à l'Adamaoua, sur plus de \SI{1000}{km}. Son élément le plus spectaculaire est le \cle{mont Cameroun}, volcan actif qui domine la plaine côtière au-dessus de Buéa et Limbé.

\begin{exemplebox}[Le mont Cameroun, point culminant]
Le \cle{mont Cameroun} culmine à \SI{4095}{m} : c'est le \textbf{point le plus élevé d'Afrique de l'Ouest et d'Afrique centrale}. C'est aussi un \textbf{volcan actif} : sa dernière éruption majeure date de \textbf{1999--2000} (des coulées de lave ont alors menacé les environs de Buéa). Ses pentes, arrosées par des pluies parmi les plus abondantes d'Afrique (plus de \SI{10000}{mm} par an à Debundscha, au pied du volcan), portent des plantations de bananes, de thé et de palmiers à huile.
\end{exemplebox}

\begin{attention}[Erreur fréquente au Bac]
Ne confonds pas le \cle{mont Oku} (\SI{3011}{m}, deuxième sommet du Cameroun) avec le \cle{point culminant} du pays, qui est le \cle{mont Cameroun} (\SI{4095}{m}). De même, évite d'écrire que le mont Cameroun est « le plus haut sommet d'Afrique » : ce titre revient au Kilimandjaro (Tanzanie, environ \SI{5895}{m}). Le mont Cameroun est le plus haut sommet d'Afrique \emph{de l'Ouest et centrale}.
\end{attention}

\begin{savaistu}[La ligne volcanique et la tragédie du lac Nyos]
La ligne volcanique du Cameroun s'étend sur plus de \SI{1000}{km}, de l'île de \cle{Bioko} (Guinée équatoriale) dans le golfe de Guinée jusqu'au plateau de l'Adamaoua. Elle comprend des volcans actifs (mont Cameroun), des massifs éteints (Bamboutos, Oku, Manengouba) et de nombreux \cle{lacs de cratère}. L'un d'eux, le \cle{lac Nyos} (région du Nord-Ouest), a provoqué le 21 août \textbf{1986} une catastrophe unique au monde : un dégazage brutal de dioxyde de carbone accumulé au fond du lac a asphyxié environ \num{1700} personnes et des milliers de têtes de bétail. Un système de dégazage artificiel y a depuis été installé.
\end{savaistu}

\courssub{Le plateau de l'Adamaoua : la charnière du pays}

Au centre du pays s'étend le \cle{plateau de l'Adamaoua}, vaste ensemble tabulaire compris entre \num{1000} et \SI{1500}{m} d'altitude, autour de \cle{Ngaoundéré}. Il est bordé de falaises abruptes (les « falaises de l'Adamaoua ») qui marquent nettement la transition avec les plaines du Nord.

Ce plateau joue un double rôle fondamental :
\begin{itemize}
\item C'est le \cle{« château d'eau » du Cameroun} : grâce à son altitude et à ses pluies abondantes, il donne naissance à de nombreux cours d'eau qui rayonnent vers le sud (Sanaga et ses affluents : Lom, Djérem, Mbam) et vers le nord (Bénoué, Faro, Mayo Kébi). Le fleuve \cle{Sanaga}, le plus long fleuve entièrement camerounais (environ \SI{900}{km}), y prend sa source ; ses chutes et rapides (Edéa, Song Loulou) fournissent l'essentiel de l'électricité du pays.
\item C'est une \cle{zone de transition} entre le Sud forestier et humide et le Nord soudanien et sec : transition climatique, végétale (savane herbeuse), mais aussi humaine (zone de contacts entre les peuples du Sud et du Nord, terre d'élevage bovin par excellence).
\end{itemize}

\courssub{Les plaines du Nord et les reliefs résiduels du Grand Nord}

\textbf{Les plaines du Nord.} Au-delà des falaises de l'Adamaoua, le terrain s'abaisse progressivement vers deux grandes plaines d'érosion et de sédimentation, situées à moins de \SI{500}{m}, puis moins de \SI{300}{m} d'altitude :
\begin{itemize}
\item la \cle{plaine de la Bénoué}, au centre-nord (autour de Garoua, environ \SI{250}{m}), drainée par le fleuve Bénoué, affluent du Niger ; c'est une plaine de savane arborée, parsemée d'\cle{inselbergs} ;
\item la \cle{plaine du Logone} et la \cle{cuvette du lac Tchad}, à l'extrême nord (autour de Kousséri et Waza), à moins de \SI{300}{m} d'altitude : une cuvette presque parfaitement plate, inondée en saison des pluies par les débordements du Logone et du Chari, ce qui explique la fertilité des sols (cultures de décrue, pêche, élevage) et l'usage de la pirogue comme moyen de transport saisonnier.
\end{itemize}

\textbf{Les reliefs résiduels du Grand Nord.} Dans ces plaines monotones se dressent des reliefs isolés, témoins d'une érosion millénaire :
\begin{itemize}
\item les \cle{massifs des Mandara}, à la frontière nigériane (région de Mora, Mokolo), chaîne de montagnes granitiques culminant vers \SI{1500}{m} (mont Oupay, environ \SI{1494}{m}), terrassées et intensément cultivées par les montagnards (Mafa, Mofu, Kapsiki) ;
\item les \cle{inselbergs} (« îles montagneuses ») et \cle{buttes témoins}, reliefs rocheux isolés qui émergent des plaines (par exemple les pics de Rhumsiki, aux formes spectaculaires, ou les inselbergs de la plaine de la Bénoué).
\end{itemize}

\begin{popfigure}
\begin{center}
\includegraphics[width=\linewidth,height=9.5cm,keepaspectratio]{N01/cmr_physique.png}
\end{center}
\legende{Les cinq grandes unités de relief se lisent du sud au nord : plateau du Sud, hautes terres de l'Ouest, plateau de l'Adamaoua, plaines du Nord, reliefs du Grand Nord. Carte physique du Cameroun (altitudes en dégradé de vert à brun, légende en haut à gauche ; fleuves en bleu ; limites d'État en gris). On y lit : les hautes terres de l'Ouest (teintes beige-brun, plus de 1 000 m) et le mont Cameroun (tache brune isolée près de la côte) ; le vaste plateau central de l'Adamaoua et du Sud (750 à 1 000 m, teintes jaune pâle) ; les plaines basses du Nord et du bassin du lac Tchad (vert, moins de 400 m) ; la plaine littorale verte face au golfe de Guinée. \\ Source : Wikimedia Commons, Urutseg, CC0.}
\end{popfigure}

\begin{popfigure}
\begin{center}
\begin{tikzpicture}
\begin{axis}[
width=0.95\linewidth, height=6.2cm,
xlabel={Distance depuis Douala (km, ordre de grandeur)},
ylabel={Altitude (m)},
xmin=0, xmax=1200, ymin=0, ymax=4500,
xtick={0,200,400,600,800,1000,1200},
ytick={0,500,1000,1500,2000,2500,3000,3500,4000},
grid=both, grid style={popline!40},
axis line style={popdark},
]
\addplot[popcurve, smooth, fill=popblueL, fill opacity=0.5] coordinates {
(0,10) (60,50) (120,300) (200,700) (260,750) (320,900) (380,1100) (460,1150) (540,1200) (620,1100) (700,800) (780,400) (860,250) (940,300) (1020,350) (1100,400) (1200,350)
} \closedcycle;
\node[font=\scriptsize, popdark] at (axis cs:200,950) {Plateau du Sud (Yaoundé)};
\node[font=\scriptsize, popdark] at (axis cs:520,1500) {Adamaoua (1000--1500 m)};
\node[font=\scriptsize, popdark] at (axis cs:880,700) {Plaines du Nord ($<$ 300 m)};
\node[font=\scriptsize, popdark] at (axis cs:60,300) {Douala};
\node[font=\scriptsize, popdark] at (axis cs:1130,600) {Maroua};
\draw[dashed, poporange, thick] (axis cs:700,1100) -- (axis cs:700,0) node[below, font=\scriptsize, poporange]{falaises};
\end{axis}
\end{tikzpicture}
\end{center}
\legende{Profil topographique schématique sud-nord Douala $\rightarrow$ Maroua (environ \SI{1200}{km}). On lit trois grandes marches : le plateau du Sud ($\approx$ 650--750 m), le plateau de l'Adamaoua (\num{1000}--\SI{1500}{m}) franchi par un abrupt (falaises), puis les plaines du Nord ($<$ 300 m). Le mont Cameroun (\SI{4095}{m}), situé à l'ouest de ce profil, n'y apparaît pas : attention à ne pas le placer sur cet axe.}
\end{popfigure}

\courssub{Les conséquences humaines du relief}

Le relief n'est pas un décor neutre : il structure la vie des Camerounais.

\textbf{Des contrastes d'accessibilité.} Les plaines et plateaux du Sud ont facilité la construction des routes et du chemin de fer (\cle{Transcamerounais} Douala--Yaoundé--Ngaoundéré). En revanche, les falaises de l'Adamaoua, les pentes des hautes terres de l'Ouest et les massifs des Mandara ont longtemps isolé certaines régions ; les pistes des monts Mandara restent difficiles en saison des pluies. Le relief explique aussi pourquoi le rail s'arrête à Ngaoundéré : au-delà, la desserte du Grand Nord repose sur la route.

\textbf{Une répartition différenciée des activités.} Chaque unité de relief porte des activités spécifiques : port, industries et plantations sur la plaine côtière ; cultures vivrières et rentières (cacao, café) sur le plateau du Sud ; maraîchage et café arabica sur les sols volcaniques fertiles des hautes terres ; élevage bovin sur le plateau de l'Adamaoua ; coton, cultures de décrue et pêche dans les plaines du Nord ; cultures en terrasses dans les Mandara. Le relief commande ainsi la \cle{complémentarité des régions}, base de l'intégration nationale.

\textbf{Des risques liés au relief.} Le volcanisme expose les populations aux \cle{éruptions} (mont Cameroun, 1999--2000) et aux dégazages de lacs de cratère (Nyos, 1986 ; Monoun, 1984). Les fortes pentes des hautes terres et des Mandara, combinées aux pluies intenses, provoquent des \cle{glissements de terrain} (par exemple dans la région de l'Ouest, quartier de Gouache à Bafoussam en 2019) et une érosion spectaculaire des sols.

\begin{exoresolu}[Décrire le relief d'une région : l'Adamaoua]
\textbf{Énoncé.} À partir de la carte des unités de relief, décris en cinq lignes le relief de la région de l'Adamaoua et indique une conséquence humaine de ce relief.
\tcblower
\textbf{Solution détaillée.} L'Adamaoua est un \textbf{plateau} (forme tabulaire) situé au centre du Cameroun, autour de Ngaoundéré. Son altitude est élevée, comprise entre \num{1000} et \SI{1500}{m}. Il est bordé au nord et au sud par des \textbf{falaises abruptes} qui le séparent des plaines de la Bénoué et du plateau du Sud. Son modelé est doux, fait de vastes surfaces planes couvertes de savane herbeuse. Conséquence humaine : l'altitude et les pluies en font le « château d'eau » du pays (sources de la Sanaga et de la Bénoué) et une grande zone d'\textbf{élevage bovin}, tandis que les falaises ont longtemps freiné les communications entre le Nord et le Sud.
\end{exoresolu}

\begin{exercice}[À toi de jouer]
Sur un croquis vierge du Cameroun, place et légende : les cinq unités de relief, le mont Cameroun, le mont Oku, les massifs des Mandara, et les villes de Douala, Yaoundé, Ngaoundéré, Garoua et Maroua. Rédige ensuite un titre et une légende organisée (numéros ou flèches).
\end{exercice}

\begin{corrige}[Exercice : croquis des unités de relief]
Le croquis doit respecter la disposition sud-nord : (1) plateau du Sud au centre-sud, (2) hautes terres de l'Ouest en bordure occidentale avec le mont Cameroun au sud-ouest et le mont Oku au nord-ouest, (3) plateau de l'Adamaoua au centre, (4) plaines de la Bénoué puis du Logone au nord, (5) massifs des Mandara à l'extrême nord-ouest, à la frontière du Nigeria. Les villes se placent : Douala sur la côte, Yaoundé sur le plateau du Sud, Ngaoundéré sur l'Adamaoua, Garoua dans la plaine de la Bénoué, Maroua dans la plaine du Logone. La légende doit être \textbf{organisée} (teintes pour les unités, symboles pour les sommets et les villes) et accompagnée d'un titre et d'une flèche du nord.
\end{corrige}

\begin{aretenir}[L'essentiel de la section]
\begin{itemize}
\item Le Cameroun, pays-charnière entre Afrique de l'Ouest et Afrique centrale, présente \textbf{cinq grandes unités de relief} disposées en bandes sud-nord.
\item \textbf{Sud} : plaine côtière (Douala, Kribi) et plateau sud-camerounais (Yaoundé, 0 à 650 m).
\item \textbf{Ouest} : hautes terres volcaniques (Bamboutos, Oku à \SI{3011}{m}) et \textbf{mont Cameroun} (\SI{4095}{m}, point culminant, volcan actif).
\item \textbf{Centre} : plateau de l'Adamaoua (\num{1000}--\SI{1500}{m}), « château d'eau » et zone de transition.
\item \textbf{Nord} : plaines de la Bénoué et du Logone ($<$ 300 m), cuvette du lac Tchad ; reliefs résiduels : massifs des Mandara, inselbergs, buttes témoins.
\item Le relief explique les contrastes d'accessibilité, la répartition des activités et certains risques (éruptions, glissements de terrain).
\end{itemize}
\end{aretenir}

\begin{savaistu}[Complément Bac : relief et tectonique]
Pour valoriser ta copie, tu peux relier brièvement le relief à la \cle{tectonique} : la plaine de la Bénoué occupe une \textbf{fosse d'effondrement} (fossé de la Bénoué), héritée de l'ouverture de l'Atlantique Sud au Mésozoïque, et la ligne volcanique du Cameroun correspond à une \textbf{ligne de fracture} profonde de la croûte terrestre, réactivée depuis des millions d'années. Reste bref : au Baccalauréat, une phrase bien placée suffit à montrer que tu comprends l'origine des formes ; l'essentiel de l'exigence porte sur la localisation et la description.
\end{savaistu}

\coursec{Les climats du Cameroun}

Dans la section précédente, nous avons parcouru le relief du Cameroun, des plaines côtières aux hauts plateaux de l'Adamaoua et aux massifs volcaniques de l'Ouest. Ce relief n'est pas un décor inerte : il modifie les températures, capte les nuages et organise la répartition des pluies. C'est précisément ce que cette section va montrer. Imagine un voyage en bus de Limbé à Maroua en janvier : tu quittes une ville humide et verte sous une chaleur moite, tu traverses des savanes fraîches sur les hauteurs, et tu arrives dans une plaine poussiéreuse où l'air est sec et brûlant. En une journée, tu as traversé plusieurs mondes climatiques. Comprendre cette diversité, c'est comprendre pourquoi le Cameroun est surnommé « l'Afrique en miniature ».

\begin{definition}[Le climat]
Le \cle{climat} est l'état moyen de l'atmosphère — \cle{températures}, \cle{précipitations}, \cle{vents}, humidité — observé en un lieu donné sur une longue période (au moins 30 ans selon l'Organisation météorologique mondiale). Il ne faut pas le confondre avec le \emph{temps qu'il fait}, qui désigne l'état de l'atmosphère à un instant précis.
\end{definition}

\courssub{Les facteurs du climat camerounais}

Quatre grands facteurs expliquent la diversité climatique du pays.

\textbf{1. La latitude.} Le Cameroun s'étend d'environ 2° N (embouchure du Ntem, Campo) à 13° N (rive du lac Tchad). Plus on s'éloigne de l'équateur, plus la durée de la saison sèche s'allonge et plus les pluies diminuent : c'est le facteur qui organise le grand zonage sud-nord.

\textbf{2. L'altitude.} La température diminue en moyenne d'environ 0,6 °C tous les 100 mètres d'élévation. C'est pourquoi Ngaoundéré (environ 1 100 m), Dschang (environ 1 400 m) ou Bamenda (environ 1 500 m) connaissent des températures modérées malgré leur position tropicale, et pourquoi les sommets comme le mont Oku (3 011 m) ou le mont Cameroun (4 095 m) ont un climat de montagne, frais toute l'année.

\textbf{3. La proximité de l'océan.} Le golfe de Guinée fournit une masse d'air chaude et humide. Le littoral connaît une humidité permanente et des pluies abondantes ; l'influence océanique s'atténue vers l'intérieur des terres.

\textbf{4. Les vents saisonniers : la mousson et l'harmattan.} Deux flux d'air opposés rythment l'année camerounaise :
\begin{itemize}
\item la \cle{mousson}, vent humide du sud-ouest venu de l'océan Atlantique, apporte les pluies ;
\item l'\cle{harmattan}, vent sec et chaud du nord-est venu du Sahara, apporte la sécheresse et la poussière.
\end{itemize}

\begin{propriete}[Le rôle du front intertropical (FIT)]
La limite de contact entre la mousson humide et l'harmattan sec s'appelle le \cle{front intertropical} (FIT). Il migre vers le nord au printemps (jusqu'à la latitude du lac Tchad vers août) puis redescend vers le sud en automne. Là où le FIT passe et séjourne, il pleut ; là où il s'absente longtemps, la saison sèche est longue. C'est ce va-et-vient qui explique l'alternance des saisons et le raccourcissement de la saison des pluies du sud vers le nord. \emph{Complément Bac : savoir citer le FIT dans une dissertation est un plus très apprécié du correcteur.}
\end{propriete}

\begin{savaistu}[L'harmattan, un vent qui paralyse]
En janvier-février, l'harmattan charge l'air du Grand Nord de poussières sahéliennes. La visibilité peut tomber sous quelques centaines de mètres, ce qui perturbe parfois le trafic aérien à Maroua et à Garoua, irrite les voies respiratoires et dessèche les cultures. Les habitants parlent d'un « ciel blanc » caractéristique.
\end{savaistu}

\courssub{Le climat équatorial (Sud)}

Le climat \cle{équatorial} occupe le sud du pays, grossièrement au sud d'une ligne Yaoundé–Bertoua : régions du Sud, du Centre (sauf sa frange nord), du Littoral, du Sud-Ouest et de l'Est.

\textbf{Ses caractéristiques.}
\begin{itemize}
\item \textbf{Pluies abondantes} : de 1 500 à 1 700 mm à Yaoundé jusqu'à environ 10 000 mm par an à Debundscha, au pied du mont Cameroun — l'un des records mondiaux de pluviosité.
\item \textbf{Quatre saisons} : une grande saison des pluies (septembre-novembre), une grande saison sèche (décembre-février), une petite saison des pluies (mars-juin) et une petite saison sèche (juillet-août). Il pleut donc pratiquement toute l'année.
\item \textbf{Températures élevées et régulières} : moyenne annuelle autour de 25 °C, avec une amplitude thermique annuelle très faible (2 à 3 °C seulement).
\end{itemize}

\begin{attention}[Erreur fréquente]
Ne dis jamais que le climat équatorial n'a « pas de saisons ». Il en possède \textbf{quatre} : petite et grande saison sèche, petite et grande saison des pluies. Dire le contraire est une erreur classique sanctionnée au Baccalauréat. Ce qui est vrai, c'est que les \emph{températures} varient peu — mais pas les pluies.
\end{attention}

\courssub{Le climat tropical humide de transition}

Entre le Sud équatorial et le Nord sec s'intercale un climat de transition, dit \cle{tropical humide} (ou « soudanien humide »), qui couvre le plateau de l'Adamaoua et les hautes terres de l'Ouest (Grassfields : régions de l'Ouest et du Nord-Ouest).

\textbf{Ses caractéristiques.}
\begin{itemize}
\item \textbf{Deux saisons nettes} : une saison des pluies (avril à octobre) et une saison sèche (novembre à mars), car le FIT ne fait plus que passer.
\item \textbf{Pluies encore abondantes} : 1 500 à 2 000 mm par an (environ 1 600 mm à Ngaoundéré).
\item \textbf{Températures modérées par l'altitude} : moyennes de 20 à 22 °C à Ngaoundéré (1 100 m) ; nuits franches et fraîches en saison sèche, surtout sur les hautes terres de l'Ouest où le mercure peut descendre sous 15 °C.
\end{itemize}

Ce climat, ni trop chaud ni trop sec, est très favorable à l'élevage (Adamaoua) et aux cultures vivrières et maraîchères (Ouest).

\courssub{Le climat tropical sec (Nord)}

Au nord de l'Adamaoua, dans les régions du Nord et de l'Extrême-Nord (plaine du Diamaré, bassin de la Bénoué), règne le climat \cle{tropical sec}.

\textbf{Ses caractéristiques.}
\begin{itemize}
\item \textbf{Une longue saison sèche} de 6 à 8 mois (octobre-novembre à avril-mai), pendant laquelle l'harmattan domine.
\item \textbf{Pluies faibles et irrégulières} : 700 à 1 200 mm par an, concentrées sur quelques mois d'été (juin à septembre), souvent sous forme d'averses violentes.
\item \textbf{Fortes amplitudes thermiques} : journées très chaudes (35 à 40 °C en mars-avril) et nuits relativement fraîches en saison sèche ; l'amplitude annuelle atteint 8 à 10 °C, bien plus qu'au Sud.
\end{itemize}

\courssub{Le climat soudano-sahélien de l'Extrême-Nord}

À l'approche du lac Tchad (au nord de Kousséri et de Waza), le climat devient \cle{soudano-sahélien}, presque aride :
\begin{itemize}
\item pluies \textbf{inférieures à 800 mm}, parfois moins de 600 mm autour du lac, et très aléatoires d'une année à l'autre ;
\item saison sèche de 8 à 9 mois, avec une sécheresse marquée ;
\item chaleur extrême : les maxima dépassent régulièrement 40 °C en avril-mai, et l'évaporation est intense.
\end{itemize}
C'est la zone la plus fragile du pays face aux aléas climatiques : les populations y pratiquent l'élevage, la culture de mil et de sorgho, et les cultures de décrue sur les berges du Logone.

\begin{exemplebox}[De Debundscha à Maroua : le contraste camerounais]
\cle{Debundscha}, village du Sud-Ouest au pied du versant occidental du mont Cameroun, reçoit environ \textbf{10 000 mm} de pluie par an — l'un des records mondiaux. \cle{Maroua}, capitale de l'Extrême-Nord, n'en reçoit qu'environ \textbf{800 mm}. Rapport : plus de 12 fois plus de pluie à Debundscha qu'à Maroua, pour deux localités d'un même pays ! C'est l'illustration chiffrée de la diversité climatique camerounaise.
\end{exemplebox}

\courssub{Les nuances locales : montagnes et versants}

Le zonage général est modifié localement par le relief :
\begin{itemize}
\item \textbf{Climat de montagne} : au-dessus de 1 800-2 000 m (monts Oku, Bamboutos, Manengouba, sommet du mont Cameroun), les températures sont fraîches toute l'année (moyennes de 10 à 15 °C), le brouillard est fréquent et le gel peut apparaître sur les crêtes les plus hautes.
\item \textbf{Hyper-pluviosité du versant occidental du mont Cameroun} : la montagne barre la route à la mousson océanique ; l'air humide est forcé de s'élever le long du versant, se refroidit, et la vapeur d'eau se condense en pluies abondantes : c'est l'\cle{effet de barrière} (ou effet orographique). Le versant exposé au vent (« au vent ») est très arrosé, tandis que le versant opposé (« sous le vent »), où l'air redescend asséché, est nettement plus sec. C'est ce mécanisme qui explique les records de Debundscha.
\end{itemize}

\courssub{Lire et interpréter un diagramme climatique}

\begin{methode}[Construire et lire un diagramme climatique]
\begin{enumerate}
\item Porter les \textbf{mois} en abscisse (de janvier à décembre).
\item Porter les \textbf{précipitations} en barres, avec l'échelle à gauche (en mm).
\item Porter les \textbf{températures} en courbe, avec l'échelle à droite (en °C).
\item Légender le graphique (titre, axes, unités).
\item Pour l'interprétation : repérer les mois secs (barres basses) et les mois pluvieux, calculer le total annuel des pluies, repérer les températures maximales et minimales et calculer l'amplitude thermique annuelle.
\end{enumerate}
\end{methode}

\begin{popfigure}
\begin{center}
\begin{tikzpicture}
\begin{axis}[
  width=0.92\linewidth, height=6.2cm,
  ybar, bar width=5pt,
  ymin=0, ymax=350,
  ylabel={Précipitations (mm)},
  symbolic x coords={J,F,M,A,Ma,Jn,Jl,Ao,S,O,N,D},
  xtick=data,
  ylabel style={popblue},
  axis y line*=left, axis x line*=bottom,
  legend style={at={(0.5,1.02)},anchor=south,legend columns=2,draw=none,font=\small}
]
\addplot+[ybar,fill=popblue,draw=popblue] coordinates {(J,20)(F,45)(M,130)(A,180)(Ma,220)(Jn,160)(Jl,70)(Ao,100)(S,250)(O,290)(N,190)(D,30)};
\addlegendentry{Pluies Yaoundé (mm)}
\addplot+[ybar,fill=poporange,draw=poporange] coordinates {(J,0)(F,0)(M,2)(A,30)(Ma,90)(Jn,140)(Jl,190)(Ao,200)(S,120)(O,25)(N,0)(D,0)};
\addlegendentry{Pluies Maroua (mm)}
\end{axis}
\begin{axis}[
  width=0.92\linewidth, height=6.2cm,
  ymin=15, ymax=40,
  ylabel={Températures (°C)},
  symbolic x coords={J,F,M,A,Ma,Jn,Jl,Ao,S,O,N,D},
  xtick=data,
  ylabel style={popdark},
  axis y line*=right, axis x line=none,
  legend style={at={(0.5,-0.18)},anchor=north,legend columns=2,draw=none,font=\small}
]
\addplot[popcurve,mark=*,mark size=1.2pt] coordinates {(J,24.5)(F,25.5)(M,25.5)(A,25.3)(Ma,24.5)(Jn,23.5)(Jl,23)(Ao,23)(S,23.5)(O,24)(N,24.3)(D,24.3)};
\addlegendentry{Températures Yaoundé (°C)}
\addplot[popcurve,mark=square*,mark size=1.2pt,poporange] coordinates {(J,23)(F,26)(M,30)(A,33)(Ma,33)(Jn,30)(Jl,27)(Ao,26)(S,27)(O,29)(N,27)(D,24)};
\addlegendentry{Températures Maroua (°C)}
\end{axis}
\end{tikzpicture}
\end{center}
\legende{Diagrammes climatiques comparés de Yaoundé (climat équatorial) et Maroua (climat tropical sec). Barres : précipitations mensuelles (échelle de gauche) ; courbes : températures moyennes mensuelles (échelle de droite). Valeurs approchées, d'après les normales des stations météorologiques camerounaises.}
\end{popfigure}

\textbf{Lecture du document.} À Yaoundé, les barres sont hautes presque toute l'année, avec deux maxima (mai et octobre) et deux minima (janvier, juillet-août) : c'est la signature des \textbf{quatre saisons} du climat équatorial. Le total annuel des barres avoisine 1 700 mm. La courbe de température est presque plate (23 à 25,5 °C) : faible amplitude. À Maroua, les pluies se concentrent entre mai et octobre avec un seul maximum (août), pour un total annuel d'environ 800 mm, et il ne pleut quasiment pas de novembre à mars : \textbf{deux saisons} seulement. La courbe de température est en cloche, avec un maximum de 33 °C en avril-mai : forte amplitude. Ngaoundéré, station intermédiaire, présenterait un profil à deux saisons mais des températures plus douces (effet de l'altitude) et des pluies plus abondantes qu'à Maroua.

\begin{popfigure}
\begin{center}
\includegraphics[width=\linewidth,height=9cm,keepaspectratio]{N01/cmr_koppen.png}
\end{center}
\legende{Carte des climats du Cameroun d'après la classification de Köppen-Geiger (légende en anglais à droite). On lit une succession du sud au nord : climats tropicaux humides (Af et Am, bleus) au Sud et sur le littoral, climat tropical à saison sèche de type savane (Aw, bleu clair) sur tout le centre et le nord du pays, puis climats arides (BSh orange, BWh rouge) à l'Extrême-Nord vers le lac Tchad. Les petites taches jaunes (Csb) signalent les hautes terres de l'Ouest, plus fraîches. Les isohyètes ne sont pas figurées ; on retrouve les bandes climatiques à l'aide des seuils de pluie vus dans la leçon. \\ Source : Wikimedia Commons, Beck et al., CC BY 4.0.}
\end{popfigure}

\begin{exoresolu}[Comparer deux stations climatiques]
\textbf{Énoncé.} À partir du diagramme comparatif ci-dessus : 1) Calcule l'amplitude thermique annuelle à Yaoundé et à Maroua. 2) Identifie le domaine climatique de chaque station en justifiant par deux indices.
\tcblower
\textbf{Solution détaillée.}
1) À Yaoundé : $T_{max} \approx 25{,}5~°C$ (février-mars), $T_{min} \approx 23~°C$ (juillet-août), donc amplitude $\approx 25{,}5 - 23 = 2{,}5~°C$. À Maroua : $T_{max} \approx 33~°C$ (avril-mai), $T_{min} \approx 23~°C$ (janvier), donc amplitude $\approx 33 - 23 = 10~°C$.
2) Yaoundé appartient au domaine \textbf{équatorial} : pluies réparties sur presque toute l'année avec deux maxima (quatre saisons) et amplitude thermique très faible (2,5 °C). Maroua appartient au domaine \textbf{tropical sec} : pluies concentrées sur 5 mois (mai-octobre) avec une longue saison sèche de 7 mois, et forte amplitude thermique (10 °C).
\end{exoresolu}

\begin{aretenir}[L'essentiel de la section]
\begin{itemize}
\item Le climat du Cameroun résulte de la latitude, de l'altitude, de l'océan et du jeu mousson/harmattan autour du front intertropical.
\item Quatre domaines du sud au nord : \textbf{équatorial} (4 saisons, 1 500 à environ 10 000 mm), \textbf{tropical humide} (2 saisons, 1 500-2 000 mm), \textbf{tropical sec} (saison sèche de 6-8 mois, 700-1 200 mm), \textbf{soudano-sahélien} (moins de 800 mm, chaleur extrême).
\item Le relief crée des nuances locales : climat de montagne et hyper-pluviosité du versant occidental du mont Cameroun par effet de barrière (Debundscha, environ 10 000 mm).
\item Un diagramme climatique se lit en croisant barres de pluies et courbe de températures : nombre de saisons, total pluviométrique, amplitude thermique.
\end{itemize}
\end{aretenir}

\begin{exercice}[Pour t'entraîner]
Voici les moyennes annuelles de Ngaoundéré : environ 1 600 mm de pluies, saison des pluies d'avril à octobre, température moyenne de 22 °C. 1) À quel domaine climatique appartient cette station ? Justifie. 2) Pourquoi ses températures sont-elles plus basses que celles de Garoua, située plus au nord mais à environ 250 m d'altitude seulement ? 3) Construis, sur papier millimétré, le diagramme climatique de Ngaoundéré en suivant la méthode du cours (données mensuelles fournies par ton professeur).
\end{exercice}

\coursec{La végétation du Cameroun}

Dans la section précédente, nous avons vu que le climat du Cameroun s'organise en bandes successives : du climat équatorial très humide du Sud au climat sahélien sec de l'Extrême-Nord, en passant par les climats tropicaux de transition. Or, en géographie physique, il existe une loi simple et puissante : \emph{là où il pleut beaucoup et longtemps, la forêt s'installe ; là où la saison sèche s'allonge, les arbres s'espacent et les herbes dominent}. La végétation est donc le \cle{miroir du climat}. C'est exactement ce que l'on observe en traversant le Cameroun de Kribi à Kousséri : la forêt dense cède la place à la savane arborée, puis à la savane herbeuse, puis aux steppes à épineux. Cette section décrit ces grands types de végétation, explique leur répartition et analyse leurs dynamiques.

\courssub{Le principe général : la zonation latitudinale de la végétation}

\begin{definition}[Zonation latitudinale]
La \cle{zonation latitudinale} est la disposition des paysages végétaux en \textbf{bandes successives sensiblement parallèles à l'équateur}, liée à la diminution progressive des précipitations et à l'allongement de la saison sèche quand on s'éloigne de l'équateur vers le nord. Au Cameroun, on passe ainsi, du Sud vers le Nord, de la \textbf{forêt dense humide} à la \textbf{savane}, puis à la \textbf{steppe}.
\end{definition}

L'intuition est facile à saisir : chaque plante a besoin d'une quantité minimale d'eau. Les grands arbres de la forêt équatoriale exigent des pluies abondantes (plus de \SI{1500}{mm} par an) et réparties sur l'année ; les graminées de la savane supportent une longue saison sèche ; seuls les épineux et les plantes très résistantes survivent autour du lac Tchad, où il pleut moins de \SI{600}{mm} par an. La végétation traduit donc, comme sur une carte, les données climatiques étudiées précédemment.

\begin{experience}[Observer la zonation : démarche d'investigation]
\textbf{Document :} une série de quatre photographies prises le long de l'axe Yaoundé--Ngaoundéré--Maroua (route nationale N1), complétée par les données pluviométriques des villes traversées.
\begin{enumerate}
\item \textbf{Observer :} à Yaoundé, végétation dense, arbres hauts et serrés ; à Ngaoundéré, arbres isolés dans une mer d'herbes hautes ; à Garoua, herbes plus basses, arbres rares ; à Maroua, buissons épineux et sols nus.
\item \textbf{Comparer :} les pluies diminuent d'environ \SI{1600}{mm} à Yaoundé à moins de \SI{800}{mm} à Maroua, et la saison sèche passe de 3 à plus de 7 mois.
\item \textbf{Interpréter :} la végétation change parce que l'eau disponible diminue ; la zonation observée sur le terrain confirme la zonation climatique.
\item \textbf{Conclure :} la végétation est un \emph{indicateur} du climat ; on peut cartographier l'une à partir de l'autre.
\end{enumerate}
\end{experience}

\begin{popfigure}
\begin{center}
\includegraphics[width=\linewidth,height=9.5cm,keepaspectratio]{N01/afr_veg.jpg}
\end{center}
\legende{Carte de la végétation naturelle de l'Afrique (légende en anglais), sur laquelle le Cameroun est situé au fond du golfe de Guinée. On y lit la zonation en bandes : forêt dense humide (vert foncé) au sud du pays, savanes arborées puis herbeuses (verts plus clairs) au centre, steppe et semi-désert (rose) au nord vers le Tchad, et le désert plus au nord encore. Les taches violettes indiquent la végétation de montagne. À l'échelle du Cameroun, la végétation se lit comme une transition sud-nord qui suit la baisse des pluies. \\ Source : Wikimedia Commons, NCpedia (carte de la CIA), CC0.}
\end{popfigure}

\courssub{La forêt dense humide du Sud}

Le Sud du Cameroun (régions du Sud, du Centre, de l'Est, du Littoral et une partie du Sud-Ouest) est couvert par la \cle{forêt dense humide}, qui appartient au grand ensemble forestier du bassin du Congo, deuxième massif forestier du monde après l'Amazonie.

\textbf{La forêt équatoriale sempervirente.} Le mot \emph{sempervirente} signifie « toujours verte » : grâce aux pluies réparties sur l'année, les arbres ne perdent jamais tous leurs feuilles en même temps. La forêt est structurée en \textbf{étages} : des arbres géants de 40 à \SI{50}{m} (les émergents) dominent une canopée continue vers 25--\SI{30}{m}, sous laquelle vivent des arbustes et des plantes d'ombre. C'est une forêt d'une richesse botanique exceptionnelle : on y compte des centaines d'essences, dont des bois précieux très recherchés — \textbf{ébène}, \textbf{acajou}, \textbf{iroko}, \textbf{sapelli} — et des bois légers destinés au contreplaqué, comme l'\textbf{ayous} (obèche), longtemps première essence exportée du Cameroun.

\textbf{La forêt littorale et la mangrove.} Sur la côte, la forêt prend des formes particulières : la forêt littorale, et surtout la \cle{mangrove} des estuaires (Wouri, Rio del Rey), forêt amphibie dont les arbres (palétuviers) s'enracinent dans la vase des zones de marée.

\begin{exemplebox}[Une forêt encore vaste mais qui recule]
La forêt couvre encore près de \textbf{40 \%} du territoire camerounais (environ 18 à 20 millions d'hectares selon les estimations — ordre de grandeur). Mais le pays perd plusieurs \textbf{centaines de milliers d'hectares} de forêt par décennie, sous l'effet combiné de l'agriculture itinérante, de l'exploitation forestière et de l'urbanisation. La forêt camerounaise est donc à la fois un trésor national et un espace vulnérable.
\end{exemplebox}

\courssub{Les savanes arborées et herbeuses de la zone intermédiaire}

Entre la forêt du Sud et les steppes du Nord s'étend le domaine de la \cle{savane}, paysage d'herbes hautes parsemé d'arbres plus ou moins nombreux.

\textbf{La savane guinéenne.} Sur le plateau de l'Adamaoua et dans une partie de la région de l'Ouest, la savane guinéenne est une \textbf{savane arborée} : graminées hautes (parfois plus de \SI{1,5}{m}) et arbres dispersés (karités, nérés). Les pluies y restent relativement abondantes (environ \SI{1500}{mm} à Ngaoundéré), mais la saison sèche de 4 à 5 mois empêche la forêt de se reconstituer, d'autant que les \textbf{feux de brousse} annuels détruisent les jeunes arbres.

\textbf{Les prairies d'altitude des hautes terres.} Sur les hautes terres de l'Ouest et les monts Bamboutos, l'altitude (plus de \SI{1500}{m}) crée un climat frais où la végétation devient une \textbf{prairie de montagne} : herbes courtes, fougères arborescentes dans les vallons. Ces prairies, très convoitées, sont aujourd'hui largement transformées en champs et en pâturages.

\courssub{Les savanes sèches et les steppes du Nord}

\textbf{La savane soudanienne.} Dans les régions du Nord et de l'Extrême-Nord (plaine de la Bénoué), la savane devient \textbf{herbeuse} : les arbres se raréfient (acacias, baobabs, rôniers), les herbes jaunissent et sèchent pendant les 6 à 7 mois de saison sèche. Le baobab, arbre emblématique, stocke l'eau dans son tronc massif : une adaptation remarquable à la sécheresse.

\textbf{Les steppes sahéliennes.} Autour du lac Tchad (bassin de la Logone et du Chari, plaine des Masa), la végétation se réduit à des \textbf{buissons épineux} (acacias épineux) et à des herbes éphémères qui ne verdissent que pendant les quelques semaines de pluie. C'est le domaine sahélien, fragile, menacé par l'avancée de la désertification.

\courssub{Les formations particulières}

Certaines formations végétales échappent à la zonation latitudinale car elles dépendent de conditions locales (eau, altitude) :

\begin{itemize}
\item les \textbf{mangroves} de l'estuaire du Wouri et du Rio del Rey, forêts de palétuviers ancrées dans la vase, entre mer et terre ;
\item les \textbf{prairies de montagne} des hautes terres de l'Ouest et du mont Cameroun, liées à l'altitude ;
\item les \textbf{galeries-forêts}, rubans de forêt dense qui bordent les cours d'eau jusque dans la savane et même la steppe (par exemple le long de la Logone), grâce à l'eau disponible en permanence dans le sol.
\end{itemize}

\begin{popfigure}
\begin{center}
\begin{tikzpicture}[scale=1.0]
% Axe horizontal : distance sud->nord
\draw[-{Stealth},thick,popdark] (0,0) -- (12.4,0) node[below left,font=\scriptsize]{Golfe de Guinée \quad $\longrightarrow$ \quad Lac Tchad};
\draw[-{Stealth},thick,popdark] (0,0) -- (0,4.6) node[above,font=\scriptsize]{Pluies annuelles};
% Courbe des pluies (décroissante vers le nord)
\draw[very thick,popblue] plot[smooth] coordinates {(0.3,4.2) (2.5,4.0) (5,3.2) (7.5,2.4) (10,1.6) (12,1.1)};
\node[font=\scriptsize,popblue] at (3.2,4.45) {environ \SI{3000}{mm}};
\node[font=\scriptsize,popblue] at (11.2,1.7) {moins de \SI{600}{mm}};
% Bandes de végétation en bas
\fill[popgreen!55] (0.3,0.05) rectangle (3.6,0.75);
\fill[popgreen!30] (3.6,0.05) rectangle (6.4,0.75);
\fill[popgold!45] (6.4,0.05) rectangle (9.4,0.75);
\fill[poporange!40] (9.4,0.05) rectangle (12,0.75);
\node[font=\scriptsize] at (1.95,0.4) {Forêt dense};
\node[font=\scriptsize] at (5,0.4) {Savane arborée};
\node[font=\scriptsize] at (7.9,0.4) {Savane herbeuse};
\node[font=\scriptsize] at (10.7,0.4) {Steppe épineuse};
% Saison sèche (barres du haut)
\node[font=\scriptsize,popdark] at (1.95,3.0) {saison sèche : 2--3 mois};
\node[font=\scriptsize,popdark] at (5,2.55) {4--5 mois};
\node[font=\scriptsize,popdark] at (7.9,2.0) {6--7 mois};
\node[font=\scriptsize,popdark] at (10.7,1.35) {8--9 mois};
% Flèches de correspondance
\draw[-{Stealth},popmuted,dashed] (1.95,2.85) -- (1.95,0.9);
\draw[-{Stealth},popmuted,dashed] (5,2.4) -- (5,0.9);
\draw[-{Stealth},popmuted,dashed] (7.9,1.85) -- (7.9,0.9);
\draw[-{Stealth},popmuted,dashed] (10.7,1.2) -- (10.7,0.9);
\end{tikzpicture}
\end{center}
\legende{Schéma en coupe sud--nord : correspondance entre la diminution des pluies (courbe bleue), l'allongement de la saison sèche et la succession des types de végétation, du golfe de Guinée au lac Tchad. Ordres de grandeur indicatifs.}
\end{popfigure}

\begin{methode}[Analyser une photographie de paysage végétal]
\begin{enumerate}
\item \textbf{Décrire} d'abord la végétation dominante : hauteur, densité, présence d'arbres ou d'herbes, couleurs (verdoyant ou sec).
\item \textbf{Relier} cette végétation au climat : abondance des pluies, longueur de la saison sèche, altitude éventuelle.
\item \textbf{Identifier} les activités humaines visibles : cultures, élevage, exploitation forestière, habitations.
\item \textbf{Conclure} en nommant le type de paysage (forêt dense, savane arborée, steppe...) et en le localisant dans une région du Cameroun.
\end{enumerate}
\end{methode}

\begin{exoresolu}[Analyse d'une photographie de paysage]
\textbf{Énoncé.} Une photographie prise en janvier près de Ngaoundéré montre une étendue d'herbes hautes jaunies, quelques arbres isolés au feuillage clairsemé, un troupeau de zébus conduit par un berger, et une colonne de fumée à l'horizon. Décris ce paysage, identifie le type de végétation et explique-le.
\tcblower
\textbf{Solution détaillée.}
\begin{enumerate}
\item \emph{Description :} la végétation est dominée par des herbes hautes et sèches (couleur jaune en saison sèche, janvier) ; les arbres sont rares et dispersés ; la présence de fumée signale un feu de brousse.
\item \emph{Relation au climat :} il s'agit d'un climat tropical de transition (environ \SI{1500}{mm} de pluies, mais 4 à 5 mois de saison sèche) : l'eau suffit aux herbes et à quelques arbres, mais pas à une forêt dense.
\item \emph{Activités humaines :} l'élevage de zébus (l'Adamaoua est la première région d'élevage du Cameroun) et les feux de brousse allumés par les éleveurs pour renouveler l'herbe ; ces feux empêchent la régénération des arbres.
\item \emph{Conclusion :} c'est une \textbf{savane arborée guinéenne}, paysage typique du plateau de l'Adamaoua, façonné autant par le climat que par l'élevage et le feu.
\end{enumerate}
\end{exoresolu}

\courssub{Dynamiques, menaces et enjeux}

La végétation camerounaise n'est pas figée : elle recule ou se dégrade sous l'effet de plusieurs pressions.

\begin{itemize}
\item La \textbf{déforestation} : l'agriculture itinérante sur brûlis, les plantations industrielles (hévéa, palmier à huile, bananier) et l'exploitation forestière font reculer la forêt du Sud, surtout autour des axes routiers.
\item Les \textbf{feux de brousse} : répétés chaque année dans la savane, ils appauvrissent les sols et bloquent la régénération des arbres.
\item L'\textbf{extension agricole et pastorale} : dans l'Ouest très peuplé, les prairies de montagne et les galeries-forêts sont grignotées par les champs.
\item La \textbf{désertification} dans l'Extrême-Nord : la combinaison de la sécheresse, de la coupe de bois de chauffe et du surpâturage fait progresser les steppes et les sols nus vers le sud, autour du lac Tchad dont la superficie a fortement diminué depuis les années 1960.
\end{itemize}

Face à ces menaces, l'État camerounais mène une politique de \textbf{conservation} et de \textbf{reboisement} : création de parcs nationaux et de réserves de faune, aménagement forestier durable (concessions avec plans d'aménagement), campagnes de reboisement (notamment la « ceinture verte » dans le Nord pour fixer les sols et freiner la désertification).

\begin{savaistu}[La mangrove, sentinelle du littoral]
La mangrove camerounaise (estuaires du Wouri et du Rio del Rey) \textbf{protège les côtes de l'érosion} marine grâce à ses racines échasses qui fixent la vase, et elle sert de \textbf{nurserie} : de nombreux poissons et crevettes y passent leur jeunesse avant de gagner la mer. Pourtant, elle est fortement menacée par l'urbanisation de Douala, la coupe pour le bois de fumaison du poisson et la pollution portuaire. Détruire la mangrove, c'est fragiliser à la fois la côte et la pêche.
\end{savaistu}

\begin{attention}[La savane n'est pas un paysage « pauvre »]
Erreur fréquente : assimiler toute la savane à un espace stérile ou abandonné. C'est faux ! La savane de l'Adamaoua est une \textbf{grande zone d'élevage} (des millions de têtes de bétail) et d'\textbf{agriculture} (maïs, igname, coton plus au nord). La savane soudanienne est le grenier à coton du Cameroun. Un paysage de savane est un espace \emph{exploité et productif}, simplement adapté à une saison sèche longue.
\end{attention}

\begin{aretenir}[L'essentiel de la section]
\begin{itemize}
\item La végétation du Cameroun suit les climats en \textbf{zonation latitudinale} : forêt dense humide au Sud, savane arborée puis herbeuse au Centre et au Nord, steppes sahéliennes à épineux à l'Extrême-Nord.
\item La forêt équatoriale sempervirente, riche en essences précieuses (ébène, acajou, ayous), couvre encore près de 40 \% du pays.
\item Des formations particulières (mangroves, prairies d'altitude, galeries-forêts) dépendent de conditions locales d'eau et d'altitude.
\item Déforestation, feux de brousse et désertification menacent ces paysages ; la conservation (parcs) et le reboisement (ceinture verte) tentent d'y répondre.
\end{itemize}
\end{aretenir}

\begin{exemplebox}[Complément Bac : un parc national par grande zone]
Pour enrichir tes exemples à l'examen, retiens un parc national par grande zone de végétation : le parc de \textbf{Waza} (Extrême-Nord, savane soudanienne et steppe, célèbre pour ses éléphants et ses lions), le parc de la \textbf{Bénoué} (Nord, savane arborée, réserve de biosphère de l'UNESCO) et le parc de \textbf{Korup} (Sud-Ouest, l'une des plus anciennes forêts tropicales d'Afrique, d'une richesse botanique exceptionnelle). Citer ces exemples précis valorise toujours une copie.
\end{exemplebox}

\begin{exercice}[Pour t'entraîner]
Sur un croquis vierge du Cameroun, place et légende les quatre grandes zones de végétation, deux villes par zone, et indique par des flèches les deux menaces principales qui pèsent sur la forêt du Sud et sur les steppes de l'Extrême-Nord. Rédige ensuite un paragraphe de dix lignes expliquant pourquoi la végétation du Cameroun est un bon indicateur de ses climats.
\end{exercice}

\coursec{L'hydrographie du Cameroun}

Au terme de la section précédente, tu as vu que la végétation du Cameroun se distribue en grandes bandes nord-sud : forêt dense au Sud, savanes au Centre et au Nord, steppes à l'Extrême-Nord. Or cette zonation végétale ne s'explique pas seulement par le climat : elle épouse aussi la présence ou l'absence de l'eau courante. Là où les rivières sont permanentes et nombreuses, la forêt prospère ; là où elles tarissent plusieurs mois par an, la savane s'impose. Nous allons donc maintenant étudier l'\cle{hydrographie}, c'est-à-dire l'ensemble des eaux de surface d'un territoire : fleuves, rivières, lacs et retenues artificielles. Tu découvriras que le Cameroun est, à juste titre, surnommé le \cle{« château d'eau » de l'Afrique centrale}.

\courssub{Un réseau dense et contrasté}

Observe d'abord une situation concrète. En juillet, un élève de Foumban (Ouest) traverse en quelques kilomètres le Noun, puis plusieurs ruisseaux tumultueux ; la même semaine, un élève de Kousséri (Extrême-Nord) constate que le Logone n'est qu'un mince ruban d'eau entre des bancs de sable, et que nombre de petits cours d'eau des environs sont complètement à sec. Même pays, deux réalités hydrologiques opposées.

\begin{definition}[Bassin versant, ligne de partage des eaux et régime]
Un \cle{bassin versant} (ou bassin hydrographique) est l'ensemble du territoire drainé par un cours d'eau et ses affluents : toute goutte de pluie tombée dans ce territoire finit, en principe, par rejoindre ce cours d'eau. Une \cle{ligne de partage des eaux} est la limite (souvent une crête ou un plateau surélevé) qui sépare deux bassins versants voisins. Un \cle{cours d'eau intermittent} ne coule qu'une partie de l'année (saison des pluies) ; un cours d'eau \cle{permanent} coule toute l'année. Le \cle{régime} d'un cours d'eau décrit l'évolution de son débit au cours de l'année.
\end{definition}

Le réseau hydrographique camerounais présente deux grands contrastes, directement hérités du climat étudié précédemment :

\begin{itemize}
\item \textbf{Au Sud} (régions du Sud, du Centre, du Littoral, de l'Ouest, du Nord-Ouest et de l'Est), les précipitations abondantes (1 500 à plus de 4 000 mm par an) et régulières alimentent un réseau \cle{dense} de cours d'eau permanents, au débit élevé toute l'année. Les rivières y sont nombreuses, encaissées dans la forêt, et souvent interrompues par des rapides et des chutes lorsqu'elles franchissent les reliefs (escarpement atlantique, seuils volcaniques).
\item \textbf{Au Nord} (Adamaoua septentrionale, Nord, Extrême-Nord), les pluies faibles (moins de 1 000 mm, parfois moins de 400 mm autour du lac Tchad) et concentrées sur quelques mois (mai-septembre) n'alimentent qu'un réseau \cle{lâche}, fait de cours d'eau intermittents (les \emph{mayo}, en langue peule) qui débordent violemment en saison des pluies puis tarissent en saison sèche. Seuls la Bénoué, le Logone et la Chari, alimentés par des régions plus arrosées (Adamaoua, massif de l'Adamaoua, Nord-Cameroun, RCA/Tchad), restent permanents.
\end{itemize}

\begin{attention}[Ne pas confondre densité du réseau et longueur des fleuves]
Une erreur fréquente consiste à écrire que « le Nord n'a pas de grands fleuves ». C'est faux : la Bénoué et le Logone sont de longs cours d'eau (plus de 1 000 km chacun, bien que leur cours camerounais soit plus court). Ce qui manque au Nord, c'est la \emph{densité} du réseau (peu d'affluents) et la \emph{permanence} des écoulements (régime intermittent pour les mayo). Au brouillon, distingue toujours les deux idées : nombre de cours d'eau par km\textsuperscript{2} (densité) et régime (permanent ou intermittent).
\end{attention}

\courssub{Le plateau de l'Adamaoua, « château d'eau » du Cameroun}

Pourquoi tant de fleuves naissent-ils au cœur du pays ? L'explication tient au relief étudié dans la première section : le plateau de l'Adamaoua, vaste surface tabulaire culminant entre 1 000 et 1 500 m d'altitude, reçoit des pluies relativement abondantes (environ 1 500 mm par an) et constitue une immense \cle{ligne de partage des eaux}. De ses flancs, les cours d'eau s'écoulent dans trois directions opposées, alimentant trois grands bassins versants :

\begin{itemize}
\item vers le \textbf{sud et l'ouest}, la Sanaga et ses affluents (Mbam, Noun) descendent vers l'océan Atlantique (golfe de Guinée) ;
\item vers le \textbf{nord}, la Bénoué et le Mayo Kébi s'échappent vers le Nigeria ; la Bénoué rejoint le Niger, le Mayo Kébi relie (en crue) le bassin de la Bénoué à celui du Logone/Chari ;
\item vers le \textbf{sud-est}, les eaux rejoignent, par la Lom et la Kadeï, le bassin du Congo (fleuve Sangha).
\end{itemize}

C'est ce rôle de « distributeur » des eaux dans trois directions qui vaut à l'Adamaoua son surnom de \cle{château d'eau} : comme le château d'eau d'une ville alimente tous les quartiers, le plateau alimente tous les bassins du pays.

\begin{popfigure}
\begin{center}
\begin{tikzpicture}[scale=1.0]
% Plateau de l'Adamaoua
\fill[poppurpleL,draw=poppurple,thick] (0,0) ellipse (2.2 and 1.1);
\node[font=\small\bfseries,text=poppurple] at (0,0.2) {Plateau de l'Adamaoua};
\node[font=\footnotesize,text=poppurple] at (0,-0.35) {(1 000 -- 1 500 m)};
% Flèche Atlantique (sud-ouest)
\draw[->,line width=3pt,popblue] (-1.6,-0.7) .. controls (-2.5,-1.5) .. (-4.2,-2.4);
\node[font=\small,text=popblue,align=center] at (-4.6,-1.7) {Bassin atlantique\\ \footnotesize Sanaga, Wouri, Ntem, Nyong};
% Flèche Tchad (nord)
\draw[->,line width=3pt,poporange] (0.3,1.0) -- (0.9,3.1);
\node[font=\small,text=poporange,align=center] at (2.6,2.9) {Bassin du Tchad\\ \footnotesize Bénoué (Niger), Logone, Chari, Mayo Kébi};
% Flèche Congo (sud-est)
\draw[->,line width=3pt,popgreen] (1.7,-0.6) .. controls (2.8,-1.5) .. (4.3,-2.3);
\node[font=\small,text=popgreen,align=center] at (4.7,-1.5) {Bassin du Congo\\ \footnotesize Sangha, Dja, Ngoko, Lom, Kadeï};
% Pluie
\foreach \x in {-1.2,-0.4,0.4,1.2} {\draw[popblue,->,thick] (\x,2.1) -- (\x,1.35);}
\node[font=\footnotesize,text=popblue] at (0,2.35) {Pluies ($\approx$ 1 500 mm/an)};
\end{tikzpicture}
\legende{Schéma du partage des eaux sur le plateau de l'Adamaoua : les pluies tombées sur le plateau alimentent trois bassins versants (Atlantique, Congo, Tchad) dont les eaux s'écoulent dans trois directions opposées.}
\end{center}
\end{popfigure}

\courssub{Les trois grands bassins versants}

\textbf{Le bassin atlantique (ou golfe de Guinée).} C'est le plus vaste (environ 60 \% du territoire) et le plus riche en eau du pays. Il rassemble les cours d'eau qui se jettent directement dans le golfe de Guinée :

\begin{itemize}
\item la \cle{Sanaga}, plus long fleuve \emph{entièrement} camerounais (environ 920 km), naît sur l'Adamaoua (affluents Lom, Djérem), reçoit le Mbam (rive gauche) et le Noun (rive droite), puis traverse la forêt avant de se jeter dans l'Atlantique au sud d'Édéa ; ses chutes et rapides (Édéa, Song Loulou, Nachtigal) en font la colonne vertébrale énergétique du pays ;
\item le \cle{Wouri}, formé par la confluence du Nkam et du Makombé (près de Yabassi), traverse Douala et son estuaire accueille le port autonome de Douala, premier port d'Afrique centrale ;
\item le \cle{Nyong}, fleuve de forêt au régime très régulier (pluviométrique), qui traverse Akonolinga et Mbalmayo avant de se jeter dans l'océan au sud de Kribi ;
\item le \cle{Ntem}, au sud, qui marque une partie de la frontière avec le Gabon et la Guinée équatoriale ;
\item la \cle{Dibamba}, fleuve côtier qui se jette dans l'estuaire du Wouri ;
\item le \cle{Mungo}, au pied du mont Cameroun, qui marque la frontière avec le Nigeria sur une petite portion.
\end{itemize}

\textbf{Le bassin du Congo.} À l'Est et au Sud-Est (environ 15 \% du territoire), la \cle{Sangha} et ses affluents (la \cle{Ngoko}, qui forme une partie de la frontière avec le Congo-Brazzaville, le \cle{Dja}, célèbre pour sa réserve de faune classée au patrimoine mondial de l'UNESCO, la \cle{Kadeï} et la \cle{Lom}) appartiennent au bassin du fleuve Congo. Leurs eaux, après un long détour vers l'est et le nord, rejoignent l'Atlantique... par l'embouchure du Congo, à des milliers de kilomètres (RDC).

\textbf{Le bassin du Tchad (versant septentrional).} Au Nord (environ 25 \% du territoire), deux sous-ensembles distincts :
\begin{itemize}
\item la \cle{Bénoué}, qui prend sa source sur l'Adamaoua, traverse Garoua, coule vers l'ouest et rejoint le \emph{Niger} au Nigeria (à Lokoja) — c'est donc un sous-bassin du Niger. Son affluent le \cle{Mayo Kébi}, en période de crue exceptionnelle, peut déverser une partie de ses eaux vers le Logone (liaison inter-bassins) ;
\item le couple \cle{Logone}-\cle{Chari}. Le Logone (formé par la Mbéré et la Vina sur l'Adamaoua) longe la frontière tchadienne ; la Chari vient de RCA. Ils se rejoignent à N'Djaména (Tchad) et se jettent dans le lac Tchad.
\end{itemize}

\begin{attention}[La Bénoué ne se jette pas dans le lac Tchad]
Erreur très fréquente au Baccalauréat : la Bénoué n'alimente \emph{pas} le lac Tchad. Elle coule vers l'ouest et rejoint le fleuve Niger au Nigeria. Seuls le Logone et la Chari (qui se rejoignent à N'Djaména) alimentent directement le lac Tchad. Retiens : « Bénoué = Niger ; Logone + Chari = Tchad ». Le terme « bassin du Tchad » dans les manuels camerounais désigne souvent le \emph{versant septentrional} du château d'eau adamaouan (exutoire vers la dépression tchadienne), mais rigoureusement, la Bénoué appartient au bassin du Niger.
\end{attention}

\begin{popfigure}
\begin{center}
\includegraphics[width=\linewidth,height=9.5cm,keepaspectratio]{N01/cmr_map.jpg}
\end{center}
\legende{Carte générale du Cameroun (relief ombré, légende en haut à gauche) montrant les fleuves en bleu : la Bénoué (Benue) qui traverse le Nord et rejoint le Nigeria vers le Niger ; le Logone et le Chari qui rejoignent le lac Tchad au nord ; la Sanaga qui descend vers le golfe de Guinée ; la Dja et la Ngoko au sud-est qui gagnent le bassin du Congo. Les retenues de Lagdo et de Mbakaou sont aussi visibles. Les trois versants (Atlantique, Congo, Tchad) se repèrent à la direction d'écoulement de ces cours d'eau. \\ Source : Wikimedia Commons, CIA, domaine public.}
\end{popfigure}

\courssub{Les lacs du Cameroun}

Le Cameroun compte trois grandes familles de lacs :

\begin{itemize}
\item le \cle{lac Tchad}, à l'Extrême-Nord, vaste lac peu profond (profondeur moyenne 1,5 à 4 m) partagé avec le Tchad, le Niger et le Nigeria, dont le Cameroun n'occupe qu'une petite partie (archipel des îles et rives sud) ;
\item les \cle{lacs de cratère} de la dorsale camerounaise (région de l'Ouest et du Nord-Ouest) : lacs Nyos, Monoun, Oku, Barombi Mbo, Awing, nés du volcanisme et souvent très profonds (Nyos : > 200 m) ; certains sont méromictiques (eaux profondes non mélangées, riches en gaz dissous) ;
\item les \cle{lacs de barrage} (retenues artificielles) créées par les barrages hydroélectriques ou de régulation : Mbakaou (sur le Noun, affluent Sanaga), Bamendjing (sur le Mbam, affluent Sanaga), Lagdo (sur la Bénoué), Lom Pangar (sur la Lom, affluent Sanaga, mis en service 2021), Memve'ele (sur le Ntem, en construction/achèvement).
\end{itemize}

\begin{savaistu}[La régression spectaculaire du lac Tchad et le risque limnique]
Le lac Tchad est passé d'environ 25 000 km\textsuperscript{2} dans les années 1960 (période « Grand Tchad ») à moins de 2 500 km\textsuperscript{2} certaines années sèches (période « Petit Tchad », ordres de grandeur variables selon les saisons et les sources). Les causes sont à la fois naturelles (sécheresses répétées du Sahel depuis les années 1970, forte évaporation) et humaines (prélèvements croissants pour l'irrigation sur le Logone et la Chari, barrage de Maga sur le Logone). Cette régression menace la pêche, l'agriculture de décrue et l'élevage de millions de personnes, et constitue un enjeu de coopération entre les quatre pays riverains, réunis dans la Commission du bassin du lac Tchad (CBLT), qui étudie le projet de transfert d'eau du bassin du Congo (Oubangui). Autre particularité : le lac Nyos a connu en août 1986 une catastrophe limnique (dégazage brutal de dioxyde de carbone CO\textsubscript{2} piégé au fond) qui a causé la mort d'environ 1 700 personnes et 3 500 têtes de bétail ; des tuyaux de dégazage ont depuis été installés pour prévenir une récidive.
\end{savaistu}

\courssub{Les enjeux de l'hydrographie camerounaise}

L'eau courante est une ressource stratégique à plusieurs titres :

\begin{itemize}
\item \textbf{Énergie hydroélectrique.} Les chutes de la Sanaga (escarpement atlantique) alimentent les centrales d'\cle{Édéa} (1950s, 260 MW), de \cle{Song Loulou} (1980s, 388 MW) et de \cle{Nachtigal} (2023-2024, 420 MW), qui fournissent l'essentiel de l'électricité du Réseau Interconnecté Sud (RIS : Douala, Yaoundé, Ouest, Centre, Sud, Est), tandis que le barrage de \cle{Lagdo} (sur la Bénoué, 72 MW) alimente le Réseau Isolé Nord (RIN). Le potentiel hydroélectrique camerounais est considéré comme le deuxième d'Afrique subsaharienne (derrière la République démocratique du Congo), estimé à plus de 20 000 MW, encore largement inexploité : de nouveaux barrages (Lom Pangar - réservoir de régulation 6 milliards m\textsuperscript{3}, Memve'ele, Bini à Warak) sont en construction ou en service pour réguler le débit et produire.
\item \textbf{Navigation.} Limitée par les rapides (seuils), elle reste possible sur le bas Wouri (accès des navires de mer au port de Douala, chenal entretenu par dragage) et, saisonnièrement (juillet-novembre), sur la Bénoué entre Garoua et le Nigeria (port de Garoua) en période de hautes eaux.
\item \textbf{Pêche.} Le lac Tchad, le Logone, la Sanaga, le Mbam, le Noun et les lacs de barrage (Mbakaou, Bamendjing, Lagdo) soutiennent une pêche artisanale et semi-industrielle vitale pour les populations riveraines (poissons : capitaine, tilapia, carpe, silure).
\item \textbf{Irrigation.} Dans le Nord sec, les eaux du Logone et de la Bénoué irriguent les périmètres de la \cle{SEMRY} (Société d'Exploitation et de Modernisation de la Riziculture de Yagoua, riziculture autour de Yagoua, Maga, Kaélé) et les plaines de la Maga (barrage de Maga sur le Logone).
\item \textbf{Approvisionnement en eau potable.} Les captages sur la Sanaga (usine de traitement d'Ekombitié, Nkombou) alimentent Yaoundé ; le Wouri et la Dibamba alimentent Douala ; la Bénoué alimente Garoua et Maroua (via le barrage de Maga et le canal de dérivation).
\item \textbf{Risques naturels.} Les crues du Logone et de la Chari provoquent régulièrement des inondations dévastatrices dans la plaine de l'Extrême-Nord (Kaélé, Yagoua, Kousséri, Maroua), détruisant champs, habitations et infrastructures : la ressource peut devenir une menace (risque inondation, risque érosion des berges).
\end{itemize}

\begin{exemplebox}[La Sanaga, fleuve nourricier du Sud-Cameroun]
Avec ses quelque 920 km entièrement sur le territoire national, la Sanaga draine la majeure partie du Centre, de l'Ouest et de l'Adamaoua (bassin \textasciitilde 133 000 km\textsuperscript{2}). Ses rapides au franchissement de l'escarpement atlantique (Édéa, Song Loulou, Nachtigal) ont permis la construction de barrages qui fournissent plus de 90 \% de l'électricité consommée sur le Réseau Interconnecté Sud (RIS), alimentant l'usine d'aluminium ALUCAM (Édéa), les cimenteries, les villes de Douala et Yaoundé. En amont, les retenues de Mbakaou (sur le Noun, 1970s) et de Bamendjing (sur le Mbam, 1970s), complétées par Lom Pangar (sur la Lom, 2021, 6 milliards m\textsuperscript{3}), régulent son débit à l'étiage, garantissant la production électrique annuelle. Un seul fleuve assure ainsi l'essentiel de la production électrique et de l'eau potable de la moitié sud du pays.
\end{exemplebox}

\begin{propriete}[Complément Bac : un potentiel hydroélectrique de premier plan africain]
Le potentiel hydroélectrique du Cameroun est estimé comme le deuxième d'Afrique subsaharienne, derrière celui de la République démocratique du Congo (fleuve Congo). Cet atout, lié à la conjonction d'un relief accidenté (chutes et rapides sur l'escarpement atlantique et la dorsale) et de précipitations abondantes au Sud, constitue un argument majeur pour le développement national : électrification rurale, industrialisation (aluminium ALUCAM, cimenteries CIMENCAM/DANGOTE, agro-industrie) et intégration des régions dans l'économie nationale (interconnexion RIS-RIN prévue via Lom Pangar/Nachtigal). La démonstration détaillée n'est pas exigible, mais tu dois savoir mobiliser cet argument dans une synthèse sur l'intégration nationale ou les potentialités économiques.
\end{propriete}

\begin{methode}[Légender une carte hydrographique (préparation au TP 1)]
Pour construire et légender une carte de l'hydrographie du Cameroun, procède par étapes :
\begin{enumerate}
\item Trace le contour du pays et repère l'orientation (flèche du nord) et une échelle indicative (ex. 1 cm pour 200 km).
\item Reporte les fleuves principaux (Sanaga, Bénoué, Logone, Wouri, Ntem, Nyong, Sangha) par des traits bleus \textbf{épais}, et les affluents principaux (Mbam, Noun, Lom, Djérem, Dja, Ngoko, Kadeï, Chari, Vina, Mbéré) par des traits bleus \textbf{fins} : la hiérarchie doit être visible.
\item Figure les lacs naturels (Tchad, Nyos, Barombi Mbo, Oku, Monoun) par des surfaces bleues claires et les lacs de barrage (Mbakaou, Bamendjing, Lagdo, Lom Pangar, Memve'ele) par des surfaces bleues hachurées ou un symbole distinct (carré noir sur le trait du fleuve).
\item Colorie légèrement chaque bassin versant d'une teinte différente (bleu atlantique, vert Congo, orange Tchad/Nord) et nomme-les en capitales (BASIN ATLANTIQUE, BASIN DU CONGO, BASSIN DU TCHAD / VERSANT NORD).
\item Rédige la légende \textbf{sous} la carte, en reprenant chaque symbole dans l'ordre : 1. Fleuves principaux, 2. Affluents, 3. Lacs naturels, 4. Lacs de barrage / Barrages, 5. Limites des bassins versants, 6. Villes repères, 7. Flèche Nord, 8. Échelle. Donne un titre précis : « Carte hydrographique du Cameroun : réseau, bassins versants et aménagements ».
\end{enumerate}
\end{methode}

\begin{exoresolu}[Localiser et classer les cours d'eau camerounais]
Énoncé. Voici une liste de cours d'eau : Sanaga, Bénoué, Sangha, Logone, Nyong, Dja, Ntem, Chari. 1) Classe-les selon leur bassin versant (exutoire final). 2) Indique lequel est le plus long fleuve entièrement camerounais. 3) Explique pourquoi la Bénoué est rattachée au « bassin du Tchad » (versant nord) dans la classification camerounaise alors qu'elle ne se jette pas dans le lac Tchad.
\tcblower
Solution. 1) \textbf{Bassin atlantique (golfe de Guinée)} : Sanaga, Nyong, Ntem, Wouri (ils se jettent directement dans le golfe de Guinée). \textbf{Bassin du Congo} : Sangha et son affluent le Dja (rejoignent le fleuve Congo puis l'Atlantique en RDC). \textbf{Bassin du Tchad / Versant septentrional} : Bénoué, Logone, Chari (s'écoulent vers le nord, vers la dépression tchadienne). 2) Le plus long fleuve \emph{entièrement} camerounais est la \textbf{Sanaga}, avec environ 920 km (le Logone et la Bénoué sont plus longs mais traversent plusieurs pays). 3) La Bénoué est rattachée au versant septentrional car elle draine le Nord du pays (Adamaoua, Nord) et s'écoule vers la dépression tchadienne ; mais rigoureusement, elle rejoint le fleuve Niger au Nigeria (à Lokoja) : c'est un sous-bassin du Niger. Seuls le Logone et la Chari alimentent directement le lac Tchad. La classification « bassin du Tchad » désigne donc ici le \emph{versant nord du château d'eau adamaouan} (exutoire vers le nord), et non le seul exutoire lacustre. C'est une classification régionale (pays) et non strictement hydrographique mondiale.
\end{exoresolu}

\begin{exercice}[Pour t'entraîner - TP 1 type]
Sur un croquis muet du Cameroun (fourni ou tracé à main levée) : 1) Place et nomme la Sanaga, la Bénoué, le Logone, la Sangha, le Nyong, le Wouri et le lac Tchad ; 2) Hachure les trois bassins versants de couleurs différentes (bleu, vert, orange) ; 3) Marque par un symbole distinct (carré noir, triangle) les barrages d'Édéa, Song Loulou, Nachtigal, Lagdo et Lom Pangar ; 4) Rédige la légende complète et donne un titre à ton croquis. Tu pourras comparer ta production avec les critères de la méthode ci-dessus.
\end{exercice}

En résumé, l'hydrographie camerounaise se caractérise par un réseau dense et permanent au Sud, lâche et intermittent au Nord, organisé autour du château d'eau de l'Adamaoua qui partage les eaux entre trois bassins (Atlantique, Congo, Tchad/Niger). Cette richesse hydraulique, source d'énergie (2\textsuperscript{e} potentiel africain), de poissons, d'eau d'irrigation et de navigation, constitue un facteur essentiel de l'intégration nationale (interconnexion énergétique Nord-Sud, désenclavement) — mais aussi, comme le montrent les inondations du Logone et la régression du lac Tchad, un défi majeur de gestion durable et de coopération régionale.

\coursec{Le milieu biogéographique}

Après avoir analysé l'hydrographie, qui structure le territoire par ses réseaux fluviaux et lacustres, nous devons maintenant embrasser le territoire dans sa globalité vivante. Les fleuves que nous venons d'étudier ne sont pas de simples lignes bleues sur une carte : ils irriguent des sols, abreuvent une végétation, soutiennent une faune et conditionnent des activités humaines. C'est précisément l'objet du milieu biogéographique : comprendre comment le relief, le climat, la végétation, les sols et la faune interagissent pour former des ensembles cohérents, dynamiques et fragiles que sont les \cle{écosystèmes} camerounais.

\courssub{Définition et approche systémique du milieu biogéographique}

\begin{definition}[Milieu biogéographique et écosystème]
Le \textbf{milieu biogéographique} est l'ensemble des interactions réciproques entre les composantes physiques (relief, climat, hydrographie, sols) et biologiques (végétation, faune, micro-organismes) d'un territoire donné, sous l'influence croissante de l'\textbf{anthroposystème} (les sociétés humaines et leurs activités). 

À l'échelle locale, l'unité fonctionnelle de base est l'\textbf{écosystème} : un ensemble formé par un \textbf{biotope} (milieu physique : sol, climat, eau, relief) et une \textbf{biocénose} (ensemble des êtres vivants : flore, faune, micro-organismes), unis par des flux d'énergie et de matière (cycles biogéochimiques, chaînes alimentaires). L'écosystème est un système \emph{ouvert}, \emph{dynamique} et \emph{autorégulé} (homéostasie) tant que les perturbations ne dépassent pas les seuils de résilience.
\end{definition}

Pour saisir la complexité du Cameroun, souvent qualifié d'\cle{« Afrique en miniature »} en raison de la coexistence de presque tous les grands types de milieux africains, il faut raisonner en \textbf{système}. Aucun élément n'agit isolément : le climat commande l'altération des roches et la formation des sols ; le sol filtre l'eau et stocke les nutriments pour la végétation ; la végétation protège le sol de l'érosion, recycle la matière organique et offre l'habitat (niche écologique) à la faune ; la faune disperse les graines, pollinise, régule les populations et enrichit le sol ; l'homme, enfin, prélève, transforme, fragilise ou protège cet ensemble.

\begin{popfigure}
\begin{tikzpicture}[scale=0.9, every node/.style={font=\small}]
% Central node
\node[draw=popdark, thick, circle, minimum size=2.2cm, fill=popcream, align=center] (center) {MILIEU\\BIOGÉOGRAPHIQUE\\CAMEROUNAIS};

% Satellite nodes
\node[draw=popblue, thick, ellipse, minimum width=2.2cm, minimum height=1.2cm, fill=popblueL, align=center] (climat) at (110:3.8) {CLIMAT\\(Températures, Pluies,\\Saisonnalité)};
\node[draw=poporange, thick, ellipse, minimum width=2.2cm, minimum height=1.2cm, fill=poporangeL, align=center] (sol) at (10:3.8) {SOL\\ (Type, Fertilité,\\Profondeur, Drainage)};
\node[draw=popgreen, thick, ellipse, minimum width=2.2cm, minimum height=1.2cm, fill=popgreenL, align=center] (veg) at (-70:3.8) {VÉGÉTATION\\(Forêt, Savane,\\Steppe, Mangrove)};
\node[draw=poppurple, thick, ellipse, minimum width=2.2cm, minimum height=1.2cm, fill=poppurpleL, align=center] (faune) at (-130:3.8) {FAUNE\\(Mammifères, Oiseaux,\\Poissons, Insectes)};
\node[draw=poppink, thick, ellipse, minimum width=2.2cm, minimum height=1.2cm, fill=poppinkL, align=center] (homme) at (180:3.8) {HOMME\\(Activités, Aménagements,\\Prélèvements, Protection)};

% Bidirectional arrows between center and satellites
\foreach \sat/\col in {climat/popblue, sol/poporange, veg/popgreen, faune/poppurple, homme/poppink} {
    \draw[-Stealth, thick, \col] (center) -- (\sat);
    \draw[-Stealth, thick, \col, dashed] (\sat) -- (center);
}

% Interactions between satellites (simplified)
\draw[popmuted, very thin, dashed, ->] (climat) to[bend left=15] (sol);
\draw[popmuted, very thin, dashed, ->] (sol) to[bend left=15] (veg);
\draw[popmuted, very thin, dashed, ->] (veg) to[bend left=15] (faune);
\draw[popmuted, very thin, dashed, ->] (faune) to[bend left=15] (sol);
\draw[poppink, thick, ->] (homme) to[bend right=10] (climat);
\draw[poppink, thick, ->] (homme) to[bend left=10] (veg);
\draw[poppink, thick, ->] (homme) to[bend right=10] (faune);
\draw[poppink, thick, ->] (homme) to[bend left=10] (sol);

% Legend
\node[anchor=north west, align=left, font=\footnotesize, text width=5cm] at (-5.5,-4.2) {
\textbf{Légende des interactions :}\\
\color{popblue} Climat $\leftrightarrow$ Milieu \quad
\color{poporange} Sol $\leftrightarrow$ Milieu \quad
\color{popgreen} Végétation $\leftrightarrow$ Milieu\\
\color{poppurple} Faune $\leftrightarrow$ Milieu \quad
\color{poppink} Homme $\rightarrow$ Composantes (pression forte)\\
\color{popmuted} Interactions naturelles indirectes (flèches pointillées)
};
\end{tikzpicture}
\legende{Schéma systémique des interactions au sein du milieu biogéographique camerounais. Le climat est le facteur moteur principal ; le sol est le support et la mémoire ; la végétation est le producteur et le structurant ; la faune est le consommateur et le régulateur ; l'homme est le gestionnaire (ou le perturbateur) majeur.}
\end{popfigure}

\courssub{Les sols : support édaphique et mémoire du milieu}

Les sols du Cameroun sont le fruit de l'altération des roches mères sous l'action du climat (température, eau) et des organismes vivants. Ils se distribuent en grands domaines latitudinaux et altitudinaux, calqués sur le climat et la végétation.

\begin{propriete}[Grands types de sols du Cameroun]
\begin{itemize}
    \item \textbf{Sols ferrallitiques (Sud, ~60\% du territoire) :} Issus d'altération intense sous climat équatorial (chaleur + humidité permanente). Lessivés en bases (Ca, Mg, K), riches en oxydes de fer et d'aluminium (couleur \textbf{rouge} à jaune), profonds (souvent > 10 m), acides (pH 4--5), pauvres en matière organique en surface (minéralisation rapide). \emph{Fertilité chimique faible, fertilité physique bonne (structure grumeleuse).} Adaptés aux cultures pérennes (cacao, café, hévéa, palmier à huile) sous couvert forestier.
    \item \textbf{Sols ferrugineux tropicaux (Centre-Nord, savanes) :} Sous climat tropical soudano-sahélien (saison sèche marquée). Moins lessivés, horizon d'accumulation de fer (cuirasse ferrugineuse ou \emph{latérite} souvent affleurante). Couleur \textbf{rouge-brun} à \textbf{jaune}. Moins profonds, sensibles à l'érosion hydrique (ruissellement violent) et éolienne. \emph{Fertilité moyenne à faible.} Cultures vivrières (mil, sorgho, maïs, coton) et pâturages.
    \item \textbf{Sols hydromorphes (plaines d'inondation, bas-fonds, deltas) :} Saturés en eau temporairement ou permanemment (Logone, Chari, Wouri, Sanaga, Ndong, bas-fonds des plateaux). Gleyïsés (couleurs gris-bleu, taches de rouille), riches en matière organique, potentiellement très fertiles après drainage/aménagement. Riziculture, cultures de contre-saison, pêche.
    \item \textbf{Sols volcaniques (Andosols) (Ouest, Mont Cameroun, Manengouba, Bamboutos) :} Jeunes, issus de cendres et laves récentes. \textbf{Noirs}, très profonds, structure grumeleuse exceptionnelle, très riches en nutriments (bases, phosphore), forte capacité de rétention d'eau. \emph{Extrêmement fertiles.} Café arabica, thé, légumes, pomme de terre, cultures maraîchères de haute altitude. Fragiles face à l'érosion sur versants pentus.
    \item \textbf{Sols halomorphes (côte, bas Wouri, bas Sanaga, Logone) :} Salés/sodiques (influence marine ou évaporation). Difficiles pour l'agriculture sans aménagements spécifiques (lavage, drainage). Mangroves, palétuviers.
\end{itemize}
\end{propriete}

\begin{exemplebox}[Le couple sol-végétation sur le volcan du Cameroun]
Sur les flancs du Mont Cameroun (4 095 m), on observe une zonation altitudinale parfaite : 
\begin{enumerate}
    \item \textbf{0--800 m :} Andosols jeunes sur laves récentes $\rightarrow$ Forêt dense de plaine (exploitée) $\rightarrow$ Plantations industrielles (CDC : bananiers, palmiers, hévéas).
    \item \textbf{800--1 800 m :} Andosols profonds, humifères, frais $\rightarrow$ Forêt de montagne (espèces endémiques) $\rightarrow$ Caféiers arabica, thé (Tole Tea Estate), cultures vivrières intenses (Buea, Muyuka).
    \item \textbf{1 800--2 500 m :} Sols plus acides, tourbeux $\rightarrow$ Forêt subalpine (bambous, \emph{Podocarpus}) $\rightarrow$ Pâturages extensifs (élevage bovin).
    \item \textbf{> 2 500 m :} Sols squelettiques, gelifs $\rightarrow$ Landes à \emph{Erica} / \emph{Philippia} $\rightarrow$ Aucune agriculture.
\end{enumerate}
Cet exemple illustre la \textbf{zonalité} (climat $\rightarrow$ sol $\rightarrow$ végétation) et l'\textbf{anthropisation} (remplacement de la forêt naturelle par des agrosystèmes).
\end{exemplebox}

\courssub{La faune : diversité, répartition et valeur patrimoniale}

La faune camerounaise est l'une des plus riches d'Afrique (près de \textbf{300 espèces de mammifères}, plus de \textbf{900 espèces d'oiseaux}, faune ichtyologique diverse). Sa répartition suit strictement les grands biomes (forêt dense, savane, montagne, zones humides).

\begin{exemplebox}[Grande faune forestière vs grande faune de savane]
\begin{center}
\begin{tabularx}{\linewidth}{|l|X|X|}\hline
\rowcolor{popblueL}
\textbf{Milieu} & \textbf{Espèces emblématiques} & \textbf{Aires de refuge principales} \\ \hline
\textbf{Forêt dense humide} (Sud, Est) & Gorille de l'Ouest (\emph{Gorilla gorilla}), Chimpanzé (\emph{Pan troglodytes}), Éléphant de forêt (\emph{Loxodonta cyclotis} -- plus petit, défenses droites), Bongo (\emph{Tragelaphus eurycerus}), Céphalophes, Mandrill (\emph{Mandrillus sphinx}), Drill, Pangolin géant, Calao, Perroquet gris. & Parc national de \textbf{Korup} (SW), Réserve de faune du \textbf{Dja} (Est/Sud), Parc de \textbf{Lobéké} (Est), Parc de \textbf{Boumba Bek} / \textbf{Nki} (Est), Sanctuaire de \textbf{Mengamé} (Sud). \\ \hline
\textbf{Savane boisée / herbeuse} (Nord, Adamaoua) & Lion (\emph{Panthera leo}), Éléphant de savane (\emph{Loxodonta africana} -- plus grand, défenses courbées), Girafe (\emph{Giraffa camelopardalis}), Buffle, Hippotrague, Cob de Buffon, Damalisque, Phacochère, Autruche, Marabout, Vautours. & Parc national de \textbf{Waza} (Extrême-Nord), Parc national de \textbf{Bénoué} (Nord), Parc national de \textbf{Bouba Ndjida} (Nord), Réserve de faune du \textbf{Bénoué} (zones cynégétiques). \\ \hline
\textbf{Montagne / Forêt subalpine} (Ouest, SW) & Drill (Monts Mandara), Chimpanzé (Monts Bamenda, Kupe), Espèces endémiques (crapauds, caméléons, oiseaux : \emph{Tauraco bannermani}, \emph{Platysteira laticincta}). & Mont Cameroun, Monts Bamboutos, Mont Manengouba, Monts Bakossi, Mont Kupe, Monts Mandara. \\ \hline
\textbf{Zones humides / Littoral} & Lamantin (\emph{Trichechus senegalensis}), Crocodile du Nil, Hippopotame, Tortues marines (plages Sud), Avifaune migratrice (Paléarctique), Poissons (faune du Lac Tchad, fleuves). & Réserve de \textbf{Kalamaloué} (Logone), Lac Tchad (partie camerounaise), Embouchures Wouri/Sanaga/Ntem, Plages de \textbf{Campo} / \textbf{Ebodjé}. \\ \hline
\end{tabularx}
\end{center}
\end{exemplebox}

\begin{exemplebox}[Le Parc national de Waza : un sanctuaire sahélien menacé]
Créé en 1934 (réserve) puis parc en 1968, le \textbf{Parc national de Waza} (région de l'Extrême-Nord, département du Logone-et-Chari) couvre \textbf{environ 170 000 ha} (1 700 km²). C'est l'une des dernières grandes réserves de faune de savane soudano-sahélienne d'Afrique centrale. Il est classé \textbf{Réserve de biosphère UNESCO} (1979) et site \textbf{Ramsar} (zone humide d'importance internationale, 2006) pour sa plaine d'inondation (Yaérés) alimentée par le Logone.
\textbf{Richesse faunique :} Éléphants (migrations saisonnières vers Kalamaloué/Logone), lions (population étudiée), girafes (sous-espèce \emph{G. c. peralta}, menacée), kob, damalisques, autruches, avifaune aquatique exceptionnelle (pélicans, spatules, sternes -- jusqu'à 20 000 oiseaux d'eau).
\textbf{Menaces :} Braconnage (ivoire, viande de brousse), pression pastorale (transhumance), assèchement des yaérés (barrages en amont sur le Logone : Maga, Yagoua), conflits homme-faune (éléphants détruisant les cultures riveraines), insécurité (Boko Haram limitant la surveillance).
\end{exemplebox}

\courssub{Les écosystèmes remarquables et le réseau d'aires protégées}

Le Cameroun a engagé très tôt une politique de conservation. Le réseau d'aires protégées couvre aujourd'hui \textbf{environ 20 \% du territoire national} (près de 95 000 km² sur 475 442 km²), l'un des taux les plus élevés d'Afrique. Il comprend parcs nationaux, réserves de faune, réserves de biosphère, sanctuaires, zones cynégétiques et forêts communautaires.

\begin{savaistu}[La Réserve de faune du Dja : joyau forestier mondial]
Inscrite au \textbf{Patrimoine mondial de l'UNESCO} depuis \textbf{1987} (critères ix et x), la Réserve de faune du Dja (régions de l'Est et du Sud, ~526 000 ha) est l'un des plus vastes et des mieux conservés massifs de forêt dense humide de basse altitude d'Afrique (90\% de forêt primaire). Elle entoure presque entièrement la rivière Dja (affluent du Congo). Elle abrite une biodiversité exceptionnelle : \textbf{> 1 500 espèces de plantes}, \textbf{107 mammifères} (dont gorilles, chimpanzés, éléphants de forêt, bongos, mandrills, pangolins), \textbf{320 oiseaux}, et des populations humaines (Baka « pygmées » et Bantous) dont les droits d'usage traditionnels sont reconnus. C'est un site pilote pour la gestion participative et la recherche écologique (station de la \emph{Station de recherche du Dja} gérée par l'ECOFAC / UE).
\end{savaistu}

\begin{popfigure}
\begin{center}
\includegraphics[width=\linewidth,height=9.5cm,keepaspectratio]{N01/cmr_parcs.jpg}
\end{center}
\legende{Carte des parcs nationaux du Cameroun sur fond de relief (légende d'altitudes en haut à gauche ; noms en espagnol pour le titre, noms des parcs en français). On repère les parcs sahéliens et soudaniens du Nord (Waza, Kalamaloué, Bénoué, Faro, Bouba Ndjida), les parcs de forêt dense du Sud et de l'Est (Dja, Campo-Ma'an, Boumba Bek, Nki, Lobéké) et les parcs de montagne ou littoraux de l'Ouest (Mont Cameroun, Korup, Takamanda). Les aires protégées couvrent les grands biomes du pays. \\ Source : Wikimedia Commons, Teogomez, CC0.}
\end{popfigure}

\courssub{Les déséquilibres : une dynamique sous forte pression anthropique}

Le milieu biogéographique camerounais n'est pas figé ; il est le siège de dynamiques rapides, largement anthropiques, qui menacent sa résilience.

\begin{attention}[Erreur fréquente : présenter le milieu comme statique]
\emph{Ne jamais écrire : « La forêt est là, la savane est là, les animaux y vivent. »} 
Le milieu est \textbf{dynamique} : la limite forêt-savane migre (avancée de la savane par les feux), les sols s'érodent ou s'enrichissent, les populations animales fluctuent. L'action humaine (déforestation, braconnage, feux, agriculture, urbanisation, barrages, changement climatique) est devenue le \textbf{facteur dominant} de l'évolution actuelle. Au Bac, une analyse qui ignore l'anthropisation et la dynamique temporelle est incomplète.
\end{attention}

\begin{propriete}[Principaux déséquilibres observés]
\begin{itemize}
    \item \textbf{Déforestation et dégradation forestière :} Perte nette de couverture forestière estimée à \textbf{~0,5 à 1 \% par an} (données Global Forest Watch / OSFAC). Causes : agriculture itinérante sur brûlis (petits exploitants), agro-industrie (palmier à huile, hévéa -- ex. Herakles Farms / SG Sustainable Oils dans le Sud-Ouest, controversé), exploitation forestière illégale ou non durable, coupe de bois de feu/charbon (villes : Yaoundé, Douala, Bafoussam), infrastructures (routes, barrages : Lom Pangar, Nachtigal, Memve'ele).
    \item \textbf{Braconnage et trafic d'espèces :} Criminel organisé (réseaux transfrontaliers). Cible : ivoire (éléphants de forêt et de savane), écailles de pangolin (espèces \emph{Manis gigantea}, \emph{Phataginus tricuspis}, \emph{Phataginus tetradactyla}), viande de brousse (approvisionnement villes), plumes/peaux, perroquets gris (\emph{Psittacus erithacus} -- CITES Annexe I). Effet « forêts vides » (\emph{empty forest syndrome}) : forêts intactes en structure mais défaunées.
    \item \textbf{Avancée de la désertification / dégradation des terres (Nord) :} Poussée vers le sud de l'isohyète 600 mm. Surpâturage, feux de brousse précoces, déboisement pour cultures de rente (coton -- SODECOTON) et vivrières, surexploitation nappes phréatiques. Conséquences : érosion éolienne (vents d'harmattan), ensablement cours d'eau (Logone, Chari), perte de productivité, insécurité alimentaire, migrations.
    \item \textbf{Conflits homme-faune (CHF) :} En augmentation aux interfaces aires protégées / zones riveraines. Éléphants (déprédation cultures : maïs, manioc, bananiers, cacao), hippopotames (riziculture bas-fonds), primates (champs), carnivores (prédation bétail). Réponses : clôtures électriques (testées Waza, Bénoué), tranchées, répulsifs (piment, ruches), indemnisation (peu effective), corridors écologiques.
    \item \textbf{Pollution et espèces invasives :} Pollution industrielle (Douala-Bonabéri : métaux lourds, hydrocarbures dans Wouri), pollution agricole (nitrates, pesticides dans nappes et fleuves), espèces invasives (\emph{Chromolaena odorata} -- « herbe à la fièvre », \emph{Mimosa pigra}, \emph{Eichhornia crassipes} -- jacinthe d'eau sur Wouri, Sanaga, Logone, Lac Tchad).
    \item \textbf{Changement climatique :} Augmentation T° (+0,5 à 1°C depuis 1950), pluviométrie plus erratique (saisons des pluies plus courtes/intenses au Sud, plus incertaines au Nord), fréquence accrue d'événements extrêmes (inondations Logone 2012, 2022 ; sécheresses Nord). Impact : modification aires de distribution espèces, phénologie, productivité agricole, risque incendies.
\end{itemize}
\end{propriete}

\courssub{La protection et la valorisation : enjeux de développement durable}

Face à ces menaces, le Cameroun a construit un arsenal juridique et institutionnel, soutenu par des partenaires internationaux (UE, AFD, KfW, GEF, WWF, WCS, FTNS -- \emph{Fondation pour l'Environnement et le Développement au Cameroun}).

\begin{methode}[Caractériser un milieu biogéographique (démarche type Bac)]
Pour décrire et expliquer un milieu biogéographique (ex. : « Le massif du Mont Cameroun », « La plaine du Logone », « La réserve du Dja »), suis ces \textbf{4 étapes} :
\begin{enumerate}
    \item \textbf{Localiser et situer :} Coordonnées, région, département, communes, altitude, contexte régional (carte mentale).
    \item \textbf{Analyser les facteurs physiques (Biotope) :} Climat (type, températures, pluviométrie, saisonnalité) $\rightarrow$ Relief (géologie, morphologie) $\rightarrow$ Sols (type, propriétés, aptitudes) $\rightarrow$ Hydrographie (réseau, régime, nappes).
    \item \textbf{Décrire la biocénose :} Végétation (type, structure, espèces dominantes/endémiques, dynamique) ; Faune (groupes clés, espèces emblématiques/menacées, migrations, chaînes alimentaires).
    \item \textbf{Mettre en évidence les INTERACTIONS et l'ANTHROPISATION :} Cycles (eau, carbone, nutriments) ; Relations proies-prédateurs, dispersion, pollinisation ; \textbf{Activités humaines} (agriculture, élevage, chasse, cueillette, tourisme, exploitation forestière, minière, urbaine) ; \textbf{Gestion / Protection} (statut juridique, plan de gestion, efficacité, conflits, valorisation économique -- écotourisme, PES, produits forestiers non ligneux).
\end{enumerate}
\emph{Astuce :} Utilise un \textbf{tableau croisé} (Climat $\times$ Sol $\times$ Végétation $\times$ Faune $\times$ Homme) pour structurer ta démonstration écrite ou ton croquis.
\end{methode}

\begin{aretenir}[Bilan : Diversité physique et intégration nationale]
La diversité biogéographique du Cameroun (forêts, savanes, montagnes, zones humides, volcans) est un \textbf{capital naturel majeur} :
\begin{itemize}
    \item \textbf{Réservoir de biodiversité :} > 8 000 espèces végétales (dont ~156 endémiques), ~300 mammifères, ~900 oiseaux, forte endémicité montagnarde (Monts Cameroun, Bakossi, Bamboutos, Manengouba, Kupe, Mandara).
    \item \textbf{Fournisseur de services écosystémiques :} Régulation climatique (puits de carbone forêt du bassin du Congo), régulation hydrique (château d'eau : Sanaga, Bénoué, Logone, Nyong, Ntem), fertilité sols (volcans, alluvions), ressources génétiques (pharmacopée, amélioration variétale).
    \item \textbf{Levier économique :} Bois d'œuvre (premier produit d'exportation hors pétrole), agro-industrie (cacao, café, coton, banane, palme, hévéa), pêche, \textbf{écotourisme} (potentiel sous-exploité : Waza, Korup, Dja, Mont Cameroun, Lobéké), services écosystémiques payants (REDD+, PES).
    \item \textbf{Facteur d'unité / de fragmentation :} Ressources partagées (bassins versants transfrontaliers : Lac Tchad, Fleuve Congo, Sanaga) $\rightarrow$ coopération (CBLT, CICOS, CEMAC). Mais aussi inégalités spatiales (Nord pauvre en ressources / Sud riche), conflits d'usage (agriculteurs/éleveurs, locaux/industriels, conservation/développement).
\end{itemize}
L'intégration nationale passe par une \textbf{gestion concertée, équitable et durable} de ce patrimoine commun.
\end{aretenir}

\begin{savaistu}[Biodiversité et écotourisme : un argument chiffré]
Le Cameroun abrite \textbf{plus de 8 000 espèces de plantes vasculaires} (dont ~156 endémiques strictes) et \textbf{près de 300 espèces de mammifères} (dont 27 primates -- 3 grands singes : gorille, chimpanzé, mandrill/drill). Cette « Afrique en miniature » concentre 92\% des écosystèmes africains. 
L'\textbf{écotourisme} (observation gorilles/chimpanzés à Campo Ma'an, Lobéké, Dja ; safaris photo Waza/Bénoué/Bouba Ndjida ; trekking Mont Cameroun, Monts Mandara ; birdwatching -- 900+ espèces) est un levier de \textbf{valorisation non extractive} de la biodiversité. Il génère des revenus directs (droits d'entrée, guides, hébergement, artisanat) et incite à la conservation (partage des bénéfices avec communautés riveraines -- ex. \emph{Community Based Natural Resource Management} autour de Waza, Bénoué, Dja). Au Bac, citer \textbf{le couple biodiversité/écotourisme} comme argument de développement durable et d'intégration (création d'emplois locaux, fierté nationale, ouverture internationale) est très valorisé.
\end{savaistu}

\courssub{Exercice résolu : Analyse d'un document sur la dynamique forêt-savane}

\begin{exoresolu}[Étude de la dynamique de l'écotone forêt-savane dans l'Adamaoua]
\textbf{Document :} Extrait d'une étude télédétection (Landsat 1986--2020) sur la transition forêt-savane dans le département de la Mbéré (Région de l'Adamaoua).
\begin{itemize}
    \item 1986 : Forêt galerie + îlots forestiers = 38\% ; Savane boisée/arborée = 52\% ; Cultures/Jachères = 10\%.
    \item 2020 : Forêt galerie + îlots = 22\% ; Savane = 45\% ; Cultures/Jachères = 30\% ; Surfaces nues/érodées = 3\%.
    \item Témoignage paysan (Ngoundal, 2021) : « Avant, la forêt arrivait au bord du champ. Maintenant, c'est la savane qui brûle chaque année. Le sol durcit, le maïs rend moins. Les éléphants ne passent plus. »
\end{itemize}

\tcblower
\textbf{Solution détaillée :}

\textbf{1. Identification du milieu biogéographique :} Zone de contact (écotone) entre le domaine forestier guinéo-congolais (Sud) et le domaine soudanien (Nord), sur le plateau de l'Adamaoua (alt. 1 000--1 500 m). Climat tropical de transition (Aw) : 1 500--1 800 mm/an, saison sèche 5--6 mois. Sols : ferrugineux lessivés sur socle cristallin, cuirasses ferrugineuses affleurantes. Végétation potentielle : mosaïque forêt-savane maintenue par le feu.

\textbf{2. Analyse des changements (1986--2020) :}
\begin{itemize}
    \item \textbf{Régression forestière drastique :} -16 points (perte de 42\% de la surface forestière initiale). Causes : extension des cultures (maïs, igname, arachide, coton), feux de brousse annuels (préparation champs, chasse, pâturage) qui empêchent la régénération forestière et favorisent la savane pyrogène, exploitation bois d'œuvre/feu.
    \item \textbf{Extension des surfaces cultivées :} Triplement (10\% $\rightarrow$ 30\%). Preuve de la pression démographique et agricole (immigration interne, politiques cotonnières SODECOTON).
    \item \textbf{Dégradation de la savane :} Baisse relative (52\% $\rightarrow$ 45\%) mais apparition de surfaces nues/érodées (3\%) $\rightarrow$ perte de productivité, érosion hydrique (ruissellement sur sols nus), formation de glacis.
    \item \textbf{Disparition de la grande faune :} Témoignage + perte d'habitat/connectivité (fragmentation forêt) + braconnage $\rightarrow$ extinction locale éléphants, grands ongulés.
\end{itemize}

\textbf{3. Mise en évidence des interactions (Système) :}
\begin{align*}
\text{Climat (saison sèche marquée)} &+ \text{Feux anthropiques récurrents} \rightarrow \text{Sélection espèces pyrophiles (savane)} \\
\text{Déforestation} &\rightarrow \text{Perte couverture végétale} \rightarrow \text{Battance sols ferrugineux} \rightarrow \text{Ruissellement/Érosion} \rightarrow \text{Appauvrissement sols} \rightarrow \text{Baisse rendements} \rightarrow \text{Nouveaux défrichements} \quad \text{(Boucle de rétroaction positive = dégradation)}
\end{align*}

\textbf{4. Réponses / Gestion :} Classement forêts résiduelles (Forêts classées de la Mbéré), promotion agroforesterie (parcs à \emph{Faidherbia albida}, \emph{Vitellaria paradoxa} -- karité), brûlis dirigés précoces, corridors écologiques vers parc de Mbam-Djerem, intensification agricole durable (SCV -- Semis Direct sous Couvert Végétal).

\textbf{Conclusion :} Ce milieu biogéographique de transition est \textbf{hautement dynamique et fragile}. L'anthropisation rapide (34 ans) a basculé le système d'une mosaïque forestière-savane relativement stable vers un agrosystème savanicole dégradé. La résilience est compromise sans inversion des pratiques (gestion intégrée feu/sol/eau/biodiversité).
\end{exoresolu}

\begin{exercice}[Application : Le massif forestier du Sud-Ouest]
À partir de vos connaissances et de la carte des aires protégées (figure ci-dessus), réalisez une fiche de synthèse du \textbf{massif forestier du Sud-Ouest} (Parcs de Korup, Takamanda, Bakossi, Mont Cameroun, Baie de Ndian, Sanctuaire de Banyang-Mbo).
Structurez votre réponse selon la \textbf{Méthode} (4 étapes) :
\begin{enumerate}
    \item Localisation et contexte (régions, départements, relief, climat).
    \item Biotope (géologie, sols, hydrographie).
    \item Biocénose (végétation : types, endémisme ; faune : primates, oiseaux, herpétofaune).
    \item Interactions et enjeux de conservation (pressions : agro-industrie CDC/SGSOC, palmiers, exploitation forestière, braconnage, populations locales ; réponses : statuts protégés, corridors, écotourisme, REDD+, gestion communautaire).
\end{enumerate}
\textit{Indication :} Ce massif est un \textbf{hotspot de biodiversité mondial} (Forêts de la Guinée du Golfe). Citez des espèces endémiques (ex. : \emph{Prunus africana}, \emph{Cephalophus rufilatus}, \emph{Tauraco bannermani}, \emph{Xenopus longipes} -- lac Bermin).
\end{exercice}

\begin{corrige}[Exercice n°1 : Le massif forestier du Sud-Ouest]
\textbf{1. Localisation :} Régions du Sud-Ouest (départements : Ndian, Kupe-Manengouba, Meme, Lebialem, Fako) et du Littoral (dépt. Moungo -- partie Nord). Frontière nigériane (Cross River). Massif volcanique (ligne du Cameroun) : Mont Cameroun (4 095 m), Monts Bakossi, Mont Kupe, Monts Rumpi, Monts Bamboutos (extrémité). Climat équatorial de montagne / transition : pluviométrie exceptionnelle (2 500--10 000 mm/an au pied du Mont Cameroun, Debundscha = l'un des lieux les plus arrosés du monde), températures douces en altitude.
\textbf{2. Biotope :} Roches volcaniques (basaltes, andésites, trachytes) récentes à anciennes $\rightarrow$ \textbf{Andosols} (fertiles, profonds) sur versants ; sols ferrallitiques lessivés en piedmont/plaines ; sols hydromorphes (marais, mangroves baie de Ndian). Hydrographie dense : radiaire sur volcans (Mungo, Wouri, Dibamba, Ndian, Meme, Cross River), côtière. Lacs de cratère (Barombi Mbo, Bermin, Dissoni -- endémismes poissons).
\textbf{3. Biocénose :} \textbf{Végétation :} Forêt dense de plaine (espèces : \emph{Aucoumea klaineana} -- okoumé, \emph{Terminalia superba} -- fraké, \emph{Entandrophragma} -- acajous), forêt de montagne (épiphytes, mousses, \emph{Podocarpus}, \emph{Prunus africana} -- pygeum, médicament), forêt subalpine (bambous \emph{Arundinaria}, \emph{Erica}), landes d'altitude. \textbf{Endémisme végétal très fort} (Monts Bakossi/Kupe : ~2 400 spp, 10\% endémiques). \textbf{Faune :} Primates : Gorille de Cross River (\emph{Gorilla gorilla diehli} -- CR, < 300 ind.), Chimpanzé Nigeria-Cameroun (\emph{Pan troglodytes ellioti} -- EN), Drill (\emph{Mandrillus leucophaeus} -- EN), Preuss's guenon (\emph{Cercopithecus preussi}). Oiseaux : \emph{Tauraco bannermani} (Bannerman's turaco -- EN, monts Bamenda/Bakossi), \emph{Platysteira laticincta} (Banded wattle-eye). Herpétofaune : \emph{Xenopus longipes} (Lac Bermin -- CR), \emph{Wolterstorffina} spp. (crapauds endémiques). Éléphants de forêt, buffles nains, céphalophes.
\textbf{4. Interactions / Enjeux :} \textbf{Pressions :} Agro-industrie historique (CDC -- palmiers, bananes, hévéas, thé) + nouvelle concession SGSOC (Herakles Farms -- 73 000 ha, contentieux foncier/environnemental) ; Exploitation forestière (UFA -- Unités Forestières d'Aménagement) ; Braconnage (viande de brousse villes : Kumba, Limbe, Douala) ; Prélèvements \emph{Prunus africana} (écorce -- surexploitation passée, maintenant règlementée CITES) ; Croissance démographique / urbanisation (Limbe, Kumba, Buea). \textbf{Conservation :} Réseau parcs (Korup 1986 -- 126 000 ha, 1\textsuperscript{er} parc forêt pluviale Afrique ; Takamanda 2008 -- 67 000 ha, transfrontalier Cross River NP Nigeria ; Bakossi 2007 -- 29 000 ha ; Mont Cameroun 2009 -- 58 000 ha ; Banyang-Mbo 1996 -- 66 000 ha). Projets : PSMNR-SWR (Programme de Gestion Durable des Ressources Naturelles - Sud-Ouest, GIZ/KfW), REDD+ (pilote), écotourisme (Korup, Mont Cameroun, lacs de cratère), gestion communautaire (forêts communautaires, comités de gestion). \textbf{Enjeu majeur :} Connectivité écologique (corridors entre parcs) vs fragmentation par plantations/routes. Intégration populations locales (Bakossi, Banyang, Mbo, Bakweri) dans la gouvernance.
\end{corrige}

\coursec{Diversité physique et intégration nationale (synthèse)}

Dans les sections précédentes, nous avons vu que le milieu biogéographique camerounais résulte de la combinaison du relief, du climat, de la végétation et des sols : chaque grande zone naturelle abrite des écosystèmes spécifiques, de la forêt dense du Sud aux savanes sèches de l'Extrême-Nord. Il est temps maintenant de prendre du recul et de rassembler toutes les pièces du puzzle. Cette synthèse répond à une question essentielle, qui est au cœur de la famille de situations « l'intégration nationale » : \emph{la diversité physique du Cameroun est-elle un atout ou une contrainte pour la construction de la nation ?} Tu verras qu'elle est les deux à la fois, et que toute l'habileté de l'aménageur consiste à transformer les contraintes en atouts.

\courssub{Bilan : un pays, plusieurs mondes physiques}

Avant de raisonner, faisons l'inventaire de ce que nous avons appris dans les cinq sections précédentes. Le Cameroun se présente comme une mosaïque remarquablement organisée.

\begin{exemplebox}[Un territoire aux dimensions contrastées]
Le Cameroun s'étend sur environ \num{475000} km² et s'allonge sur près de \num{1200} km du sud (la pointe de Campo, vers 2°20' N) au nord (le lac Tchad, vers 13° N). Cet allongement méridien, perpendiculaire aux zones climatiques, explique à lui seul l'ampleur des contrastes : pendant que le Sud baigne dans l'humidité équatoriale, l'Extrême-Nord connaît une saison sèche de sept à huit mois. Ajoute à cela un relief qui culmine à \num{4095} m (mont Cameroun) et plonge sous le niveau de la mer sur le littoral, et tu obtiens un condensé de géographie physique.
\end{exemplebox}

\begin{aretenir}[Les quatre chiffres-clés de la diversité physique]
\begin{itemize}
\item \cle{Cinq grandes unités de relief} : les plaines et plateaux littoraux ; les plateaux du Sud (650--900 m) ; les hautes terres de l'Ouest (ligne volcanique du Cameroun) ; les hauts plateaux de l'Adamaoua (1000--1200 m) ; les plaines et monts du Nord (Bénoué, Mandara, Logone).
\item \cle{Quatre domaines climatiques} : équatorial (sud du 4°--5° N, pluies abondantes toute l'année) ; soudanien humide (Adamaoua, saison sèche courte) ; soudanien sec (Nord, saison sèche longue) ; tropical de montagne (hautes terres de l'Ouest, fraîcheur et pluies orographiques).
\item \cle{Trois grands ensembles végétaux} : la forêt dense (Sud) ; la savane arborée et herbeuse (Adamaoua et Ouest) ; la steppe et la savane sèche (Extrême-Nord).
\item \cle{Trois bassins hydrographiques} : le bassin atlantique (Sanaga, Wouri, Nyong, Ntem) ; le bassin du Tchad (Logone, Chari, El Beid) ; le bassin du Congo (Dja, Boumba, Ngoko, Sangha).
\end{itemize}
\end{aretenir}

\begin{popfigure}
\begin{center}
\includegraphics[width=\linewidth,height=9.5cm,keepaspectratio]{N01/cmr_physique.png}
\end{center}
\legende{Fond de carte : carte physique du Cameroun (altitudes, fleuves, limites d'État). Le croquis de synthèse attendu superpose sur ce fond : 1. la forêt dense (climat équatorial) au Sud ; 2. les savanes (climat soudanien humide, plateau de l'Adamaoua) au centre ; 3. la steppe et la savane sèche (climat soudanien sec) au Nord ; 4. la ligne volcanique de l'Ouest (climat de montagne) ; 5. les limites approximatives des domaines climatiques ; 6. les trois directions d'écoulement des eaux (Atlantique, Tchad, Congo). Échelle : 1 cm $\approx$ 200 km. \emph{Exercice : reproduis ce croquis et organise la légende en trois rubriques (relief, climat, végétation).} \\ Source : Wikimedia Commons, Urutseg, CC0.}
\end{popfigure}

\begin{attention}[Erreur fréquente au Bac : la carte « mille-feuille »]
Beaucoup de candidats, voulant montrer qu'ils « savent tout », empilent sur une même carte le relief, les climats, la végétation, l'hydrographie, les villes, les routes... Résultat : un croquis illisible, sanctionné par le correcteur. \textbf{Règle d'or : une bonne carte physique ne traite qu'un trait à la fois.} Le croquis de synthèse ci-dessus est un exercice de superposition à visée pédagogique ; dans une copie, tu ne représenteras que le thème demandé par le sujet, avec 3 à 5 rubriques de légende maximum.
\end{attention}

\courssub{La diversité physique : une richesse pour la nation}

Commençons par l'intuition. Imagine un marché où chaque région du Cameroun apporterait ce qu'elle produit le mieux : le Sud déposerait son cacao et son bois, l'Ouest son café et ses légumes, l'Adamaoua sa viande, le Nord son coton et ses oignons, le littoral son poisson et ses bananes. Ce marché imaginaire, c'est exactement ce que permet la diversité physique : des \cle{complémentarités} entre régions.

\begin{definition}[Complémentarité régionale]
On appelle \cle{complémentarité régionale} la situation dans laquelle des régions aux milieux physiques différents produisent des biens différents et échangent entre elles, chacune spécialisée dans ce que son milieu lui permet de produire le mieux. La diversité physique crée ainsi une interdépendance économique qui renforce la solidarité nationale.
\end{definition}

Concrètement, cette richesse se décline en quatre registres :

\begin{itemize}
\item \textbf{Les complémentarités agricoles.} Le cacao et l'hévéa prospèrent dans la forêt humide du Sud (Centre, Sud, Est) ; le café arabica et les cultures vivrières de montagne (pommes de terre, légumes) caractérisent les hautes terres de l'Ouest ; l'Adamaoua, vaste savane herbeuse épargnée par la mouche tsé-tsé, est le premier pôle d'\cle{élevage} bovin du pays (plusieurs millions de têtes, ordre de grandeur : 6 à 8 millions) ; le Nord fournit le \cle{coton} (filière organisée par la SODECOTON), le mil et le niébé. Aucune région ne peut se passer des autres.
\item \textbf{Le potentiel touristique.} Plages de Kribi et de Limbe, parcs nationaux (Waza dans l'Extrême-Nord, Bénoué, Campo-Ma'an, Lobéké), montagnes de l'Ouest et leurs chefferies, chutes d'Ekom-Nkam... La variété des paysages est un argument touristique unique en Afrique centrale.
\item \textbf{Les ressources hydrauliques.} Le régime régulier des fleuves du Sud (Sanaga notamment) offre un potentiel hydroélectrique considérable, estimé à environ \num{20000} MW (ordre de grandeur, potentiel théorique), le deuxième d'Afrique subsaharienne après la RDC.
\item \textbf{La diversité biologique.} Forêts du bassin du Congo et savanes soudaniennes font du Cameroun l'un des pays les plus riches d'Afrique en biodiversité, avec des emblèmes comme les gorilles des forêts du Sud-Est et les éléphants des savanes du Nord.
\end{itemize}

\begin{savaistu}[Pourquoi dit-on « l'Afrique en miniature » ?]
Le surnom du Cameroun, \emph{« Afrique en miniature »}, vient précisément de cette concentration exceptionnelle : sur un seul territoire de taille moyenne se succèdent, du golfe de Guinée au lac Tchad, presque tous les paysages du continent — forêt équatoriale, savanes humides et sèches, steppes sahéliennes, hautes montagnes volcaniques, littoral mangrove. Peu de pays africains offrent un tel échantillon. C'est une fierté nationale... mais aussi, nous allons le voir, un défi.
\end{savaistu}

\courssub{La diversité physique : une contrainte à maîtriser}

Renversons maintenant la perspective. La même diversité qui enrichit le pays complique aussi son organisation. Le raisonnement est simple : plus les milieux sont contrastés, plus il est coûteux de les relier, de les peupler harmonieusement et de les protéger.

\begin{itemize}
\item \textbf{Des contrastes d'accessibilité.} La forêt dense du Sud-Est, les reliefs accidentés de l'Ouest et les plaines inondables du Logone rendent la construction et l'entretien des routes difficiles. En saison des pluies, certaines pistes de l'Est ou de l'Extrême-Nord deviennent impraticables : c'est le problème de l'\cle{enclavement}.
\item \textbf{Des déséquilibres de peuplement.} Les milieux les plus accueillants (littoral, hautes terres de l'Ouest, plateau du Centre) concentrent l'essentiel de la population : l'Ouest atteint des densités rurales parmi les plus élevées d'Afrique (plus de 200 hab/km² dans certains départements, ordre de grandeur), tandis que l'Est et l'Adamaoua restent faiblement peuplés (souvent moins de 10 hab/km²). La diversité physique produit donc des « Cameroun » démographiques très inégaux.
\item \textbf{Des coûts d'aménagement élevés.} Franchir les collines de l'Ouest, stabiliser les sols latéritiques détrempés du Sud, protéger les berges du Logone : chaque kilomètre d'infrastructure coûte plus cher dans un milieu difficile.
\item \textbf{Des risques naturels différenciés.} Chaque zone a ses aléas : éruptions et séismes sur la ligne volcanique (mont Cameroun, éruptions récentes en 1999 et 2000), inondations dans le bassin du Logone (fortes crues récurrentes, notamment en 2012, 2020 et 2022), sécheresses et avancée de la désertification dans l'Extrême-Nord, glissements de terrain dans les montagnes de l'Ouest (par exemple dans la région de Bafoussam, en 2019).
\end{itemize}

\begin{attention}[Piège de rédaction]
Ne présente jamais la diversité physique comme une « fatalité » ou une « malédiction ». Au Bac, le correcteur attend un raisonnement nuancé : une contrainte n'est pas une impossibilité, c'est un défi qui appelle une réponse d'aménagement. Évite aussi le déterminisme géographique (« le climat explique la pauvreté ») : le milieu influence les activités humaines, il ne les commande pas.
\end{attention}

\courssub{Les réponses de l'État : aménager pour unir}

Face à ces contraintes, l'État camerounais a engagé, depuis l'indépendance, une politique volontariste d'aménagement du territoire. Trois grandes familles d'actions peuvent être retenues.

\textbf{1. Le désenclavement par les infrastructures.} Le réseau routier relie progressivement les capitales régionales (axe Douala--Yaoundé--Bafoussam, route Yaoundé--Ngaoundéré, bitumage progressif de l'axe Ngaoundéré--Garoua--Maroua). Le chemin de fer \cle{Transcamerounais}, construit en deux étapes (Douala--Yaoundé en 1927 sous tutelle allemande puis française, prolongé jusqu'à Ngaoundéré en 1974), franchit sur environ \num{886} km presque toutes les grandes unités physiques du pays : c'est littéralement une ligne d'intégration nationale, qui transporte vers le port de Douala le coton du Nord, le bois du Sud et les produits de l'Adamaoua.

\textbf{2. La maîtrise de l'eau et de l'énergie.} Les grands barrages hydroélectriques exploitent le potentiel des fleuves : Song-Loulou (384 MW) et Edéa (environ 264 MW) sur la Sanaga, Lagdo (72 MW) sur la Bénoué — qui régularise aussi les crues et soutient l'irrigation —, et le barrage de Memve'ele (environ 211 MW) sur le Ntem, entré en service récemment. Ces ouvrages alimentent les villes et les industries et réduisent la dépendance énergétique.

\textbf{3. La conservation du patrimoine naturel.} L'État a créé un réseau d'aires protégées (une vingtaine de parcs nationaux, dont Waza, Bénoué, Korup, Campo-Ma'an, Lobéké) couvrant une part importante du territoire, ainsi que des réserves forestières et des zones de chasse. La loi de 1994 sur les forêts et la faune encadre l'exploitation. Objectif : faire de la diversité biologique un héritage durable plutôt qu'une ressource pillée.

\begin{exoresolu}[Exercice type Bac : commentaire de tableau]
\textbf{Énoncé.} Voici un tableau simplifié de correspondance entre milieux physiques et activités humaines au Cameroun.

\begin{center}
\begin{tabularx}{\linewidth}{|l|X|X|X|}
\hline
\rowcolor{popblueL} \textbf{Grande zone} & \textbf{Relief -- climat} & \textbf{Végétation} & \textbf{Activités dominantes} \\
\hline
Sud forestier & Plateaux 650--900 m ; équatorial humide & Forêt dense & Cacao, bois, vivrier \\
\hline
Ouest & Hautes terres volcaniques ; montagne & Savane arborée, galeries & Café, légumes, forte densité \\
\hline
Adamaoua & Hauts plateaux 1000--1200 m ; soudanien humide & Savane herbeuse & Élevage bovin, maïs \\
\hline
Nord & Plaines et inselbergs ; soudanien sec & Savane sèche, steppe & Coton, mil, pêche (Logone) \\
\hline
\end{tabularx}
\end{center}

À l'aide du tableau et de tes connaissances : 1) montre que le milieu physique oriente les activités humaines ; 2) explique en quoi cette spécialisation renforce l'intégration nationale.
\tcblower
\textbf{Solution détaillée.} 1) On observe une correspondance nette entre milieu et production : l'humidité équatoriale du Sud permet les cultures pérennes de forêt (cacao) ; la fraîcheur des hautes terres de l'Ouest convient au café arabica et aux cultures maraîchères ; la savane herbeuse de l'Adamaoua, sans tsé-tsé, se prête à l'élevage extensif ; la saison sèche longue du Nord favorise le coton et les céréales soudaniennes. Le milieu n'impose pas, mais il \emph{oriente} : chaque région exploite ses avantages physiques. 2) Cette spécialisation crée des besoins réciproques : le Nord a besoin du cacao et du bois du Sud, le Sud consomme la viande de l'Adamaoua et le coton transformé du Nord. Les échanges interrégionaux, acheminés par la route et le Transcamerounais vers Douala, tissent une solidarité économique qui est un fondement concret de l'unité nationale. La diversité physique devient ainsi, grâce aux échanges, un ciment de la nation.
\end{exoresolu}

\courssub{TP 1 : construire une carte physique du Cameroun}

Le TP 1 te demande de construire une carte physique du Cameroun sur \emph{un seul} trait au choix : les unités de relief, les domaines climatiques, l'hydrographie ou la végétation. Voici la démarche complète.

\begin{methode}[Réussir un croquis géographique (règles générales)]
\begin{enumerate}
\item \textbf{Tracer le fond de carte} : contour simplifié du Cameroun, respectant l'allongement nord-sud (environ 1,5 fois plus long que large).
\item \textbf{Donner un titre} précis, en haut : il annonce le thème et l'espace (ex. : « Les domaines climatiques du Cameroun »).
\item \textbf{Organiser la légende} en rubriques hiérarchisées : d'abord les \cle{aplats} (couleurs des zones), puis les \cle{symboles} (points, surfaces), enfin les \cle{figurés} linéaires (flèches, limites).
\item \textbf{Orienter} la carte (flèche du nord) et placer quelques \textbf{noms repères} : Douala, Yaoundé, Ngaoundéré, Garoua, Maroua, océan Atlantique, lac Tchad.
\item \textbf{Vérifier la cohérence} : chaque élément figuré sur la carte doit apparaître dans la légende, et réciproquement.
\end{enumerate}
\end{methode}

\begin{methode}[Démarche spécifique du TP 1]
\begin{enumerate}
\item \textbf{Choisir un seul trait} (relief, climat, hydrographie ou végétation) et s'y tenir : toute information étrangère au thème est une faute de méthode.
\item \textbf{Hiérarchiser l'information} : 3 à 5 rubriques de légende maximum (ex. pour la végétation : forêt dense / savane arborée / savane sèche et steppe / mangrove littorale).
\item \textbf{Soigner la lisibilité avant la couleur} : un croquis au crayon bien construit, avec des limites nettes et des noms bien placés, vaut mieux qu'une carte bariolée et confuse. La couleur vient en dernier, avec des teintes logiques (vert pour la forêt, jaune-orangé pour les zones sèches, bleu pour l'eau).
\item \textbf{Comparer avec un atlas} puis auto-évaluer : titre ? légende complète ? nord ? noms repères ?
\end{enumerate}
\end{methode}

\begin{experience}[Protocole du TP 1 (séance de 2 h)]
\textbf{Matériel} : fond de carte vierge du Cameroun (ou papier calque sur atlas), crayon, règle, crayons de couleur, atlas géographique. \textbf{Protocole} : (1) 10 min — choix du trait et repérage des documents de l'atlas ; (2) 20 min — tracé du contour et des limites internes au crayon ; (3) 30 min — report des informations (aplats, symboles, figurés) ; (4) 20 min — rédaction du titre et construction de la légende ; (5) 20 min — placement des noms repères, orientation, relecture croisée avec un camarade ; (6) 20 min — mise en couleur et présentation orale de 2 minutes devant la classe. \textbf{Observations attendues} : chaque élève doit pouvoir justifier chaque limite tracée (ex. : « la limite forêt/savane suit approximativement le 4°--5° N et le rebord de l'Adamaoua »).
\end{experience}

\courssub{Conclusion : la diversité, condition et enjeu de l'intégration nationale}

Revenons à notre question de départ. La diversité physique du Cameroun est un \cle{atout} par les complémentarités agricoles, le potentiel touristique, hydraulique et biologique qu'elle offre ; elle est une \cle{contrainte} par les contrastes d'accessibilité, les déséquilibres de peuplement, les coûts d'aménagement et les risques naturels différenciés qu'elle impose. Mais la leçon essentielle est ailleurs : c'est précisément parce que les milieux sont différents que les régions ont besoin les unes des autres. La route, le rail, le barrage et l'aire protégée sont les instruments qui transforment une juxtaposition de milieux en un territoire national solidaire. \textbf{La maîtrise de la diversité physique est donc une condition de l'unité et de l'intégration nationale} : un pays « Afrique en miniature » ne peut être fort que s'il est un pays relié, aménagé et partagé.

\begin{propriete}[Complément Bac : construire un plan dialectique]
En dissertation de géographie, l'opposition « atout / contrainte » de la diversité physique permet de construire un \cle{plan dialectique} très efficace en trois parties : I. La diversité physique, une richesse potentielle (complémentarités, ressources) ; II. ...mais aussi une contrainte pour l'unité (enclavement, déséquilibres, risques) ; III. Les réponses de l'État : aménager la diversité pour construire la nation (infrastructures, barrages, conservation). Ce schéma est transposable à de nombreux sujets du module « Le Cameroun ».
\end{propriete}

\begin{exercice}[Pour t'entraîner]
1) Construis, sur fond de carte vierge, le croquis « Les trois bassins hydrographiques du Cameroun » (TP 1, trait : hydrographie) avec titre, légende de 3 rubriques, nord et cinq noms repères. 2) Rédige un paragraphe argumenté (15 lignes) répondant à la question : « La diversité physique du Cameroun : handicap ou chance pour l'intégration nationale ? » 3) Calcule : si l'Extrême-Nord reçoit environ 400 mm de pluie par an et Debundscha (pied du mont Cameroun) plus de \num{10000} mm, quel est le rapport entre ces deux pluviométries ? Que traduit-il ?
\end{exercice}

\newpage

\coursec{Activité d'intégration}

\coursec{TP 1 : Construction d'une carte physique du Cameroun}

\courssub{Objectifs de l'activité}

À l'issue de ce travail pratique, tu seras capable de :
\begin{itemize}
\item \cle{inventorier} et \cle{classer} des informations géographiques issues de documents variés (cartes d'atlas, tableaux statistiques, photographies) ;
\item \cle{élaborer une légende hiérarchisée} de 3 à 5 rubriques, conforme aux règles de la cartographie scolaire ;
\item \cle{reporter} des informations sur un fond de carte muet à l'aide d'aplats de couleurs et de symboles adaptés ;
\item \cle{confronter} plusieurs cartes thématiques pour construire un croquis de synthèse ;
\item \cle{argumenter} une réponse à une question de géographie : la diversité physique du Cameroun est-elle un atout ou un handicap pour l'intégration nationale ?
\end{itemize}

\courssub{Situation}

\textbf{Contexte.} Dans le cadre de la préparation de la fête de l'Unité nationale, le club de géographie du lycée de Bafoussam a été sollicité par le proviseur pour réaliser une exposition intitulée \emph{« Un pays, mille paysages : la diversité physique du Cameroun »}. Les visiteurs, élèves et parents, doivent pouvoir comprendre en quelques minutes comment le relief, le climat, la végétation et les eaux structurent l'espace national, de la plaine du Logone aux sommets de la dorsale camerounaise.

La classe de Terminale est divisée en \textbf{quatre groupes}, chacun chargé d'un trait physique : groupe 1 (unités de relief), groupe 2 (domaines climatiques), groupe 3 (hydrographie), groupe 4 (végétation). Chaque groupe reçoit un \textbf{dossier documentaire} composé d'un fond de carte muet du Cameroun, d'une carte de l'atlas scolaire, d'un tableau de données et de deux photographies de paysages.

\textbf{Document 1 — Tableau de données communes (extraits).}

\begin{center}
\begin{tabularx}{\linewidth}{|l|X|}
\hline
\rowcolor{popblueL} \textbf{Donnée} & \textbf{Valeur repère} \\
\hline
Altitude du mont Cameroun (Fako) & environ \SI{4095}{m} (point culminant du pays ; certaines sources récentes donnent environ \SI{4040}{m}) \\
\hline
Altitude des monts Mandara & environ \SI{1494}{m} (mont Oupay, point culminant du massif) \\
\hline
Altitude des hauts plateaux de l'Ouest & \num{1500} à \SI{2500}{m} environ (monts Bamboutos : environ \SI{2740}{m}) \\
\hline
Altitude du plateau de l'Adamaoua & \num{1000} à \SI{1300}{m} environ \\
\hline
Altitude de la plaine du Logone & moins de \SI{300}{m} \\
\hline
Pluies annuelles à Debundscha (côte du Sud-Ouest) & environ \SI{10000}{mm} (ordre de grandeur ; l'une des stations les plus arrosées d'Afrique) \\
\hline
Pluies annuelles à Yaoundé & environ \SI{1600}{mm} \\
\hline
Pluies annuelles à Ngaoundéré & environ \SI{1500}{mm} \\
\hline
Pluies annuelles à Maroua & environ \SI{800}{mm} \\
\hline
Pluies annuelles à Kousséri (Extrême-Nord) & moins de \SI{600}{mm} \\
\hline
Longueur de la Sanaga & environ \SI{918}{km} (plus long fleuve entièrement camerounais) \\
\hline
Longueur du Logone & environ \SI{1000}{km} (avec le Chari, en direction du lac Tchad) \\
\hline
\end{tabularx}
\end{center}

\textbf{Document 2 — Photographies (décrites).} Photographie A : forêt dense et sombre, arbres géants serrés, sous-bois humide, prise près de Kribi (région du Sud). Photographie B : étendue d'herbes jaunies parsemée de karités et de baobabs, sol rougeâtre, prise entre Garoua et Maroua (région du Nord).

\textbf{Problème.} Comment construire, à partir de ce dossier, une carte physique lisible et correctement légendée qui montre la diversité physique du Cameroun ? Et, au-delà, cette diversité renforce-t-elle ou fragilise-t-elle l'unité du pays ?

\courssub{Consignes}

\begin{enumerate}
\item \textbf{Observer et inventorier.} Dans ton groupe, lis attentivement tous les documents. Dresse la liste des informations géographiques utiles à ton trait physique (noms, localisations, valeurs chiffrées) et classe-les en catégories.
\item \textbf{Élaborer la légende.} Construis une légende hiérarchisée de 3 à 5 rubriques : choisis un titre général de légende, ordonne les rubriques (par exemple du plus haut au plus bas, du plus humide au plus sec), et attribue à chacune une couleur ou un symbole cohérent avec les conventions cartographiques.
\item \textbf{Reporter sur le fond de carte.} Reporte les informations par aplats de couleurs (pour les zones) et par symboles (points pour les sommets et villes, traits bleus pour les cours d'eau, flèches pour les flux). Respecte les localisations relatives (le mont Cameroun au sud-ouest, la plaine du Logone à l'extrême nord, etc.).
\item \textbf{Titrer et commenter.} Rédige un titre précis (thème + espace + éventuellement date) et une courte légende commentée de 3 à 5 lignes expliquant l'organisation spatiale mise en évidence.
\item \textbf{Présenter et confronter.} Présente ta carte en 3 minutes. Puis, avec la classe, confronte les quatre cartes : quelles superpositions observes-tu entre relief, climat, végétation et hydrographie ?
\item \textbf{Synthétiser.} Construis collectivement un croquis de synthèse intitulé « Diversité physique et intégration nationale », puis réponds en une quinzaine de lignes à la question : cette diversité est-elle un atout ou un handicap pour l'unité du pays ?
\end{enumerate}

\courssub{Ressources à mobiliser}

\begin{aretenir}[Rappels utiles pour le TP]
\begin{itemize}
\item \textbf{Unités de relief} : plaine côtière et basse vallée de la Sanaga ; dorsale camerounaise (mont Cameroun, environ \SI{4095}{m}, monts Bamboutos, Manengouba) ; hauts plateaux de l'Ouest ; plateau de l'Adamaoua ; plateaux du Sud ; monts Mandara ; plaines du Nord (bas-Chari, Logone, Diamaré).
\item \textbf{Domaines climatiques} : climat équatorial (avec notamment le sous-type guinéen hyperhumide de la côte et le sous-type camerounien du Centre-Sud, à 4 saisons) ; climat tropical humide ou soudanien (Adamaoua) ; climat tropical sec ou sahélien (Nord et Extrême-Nord, 2 saisons).
\item \textbf{Végétation} : forêt dense humide (Sud) ; savane arborée et herbeuse (Adamaoua, Nord) ; steppe et formations sahéliennes (Extrême-Nord) ; mangrove côtière ; prairies d'altitude sur la dorsale.
\item \textbf{Hydrographie} : réseau atlantique (Sanaga, Nyong, Ntem, Wouri, Dibamba, Mungo) ; réseau du Congo (Dja, Sangha, Ngoko) ; réseau du Tchad (Logone, Chari) ; la Bénoué, qui traverse Garoua, est un affluent du Niger au Nigeria ; lacs (Tchad, Ossa, Barombi Mbo, Nyos).
\item \textbf{Règles de cartographie} : une légende se lit du haut vers le bas, du plus remarquable au moins remarquable ; le bleu est réservé à l'eau ; les couleurs chaudes suggèrent l'aridité ou l'altitude, les couleurs froides l'humidité ; tout symbole figurant sur la carte doit figurer dans la légende, et réciproquement.
\end{itemize}
\end{aretenir}

\courssub{Démarche et corrigé commenté}

\begin{exoresolu}[TP 1 — Construire une carte physique du Cameroun]
Énoncé : à partir du dossier documentaire, construire et légender une carte physique du Cameroun sur le trait choisi par ton groupe, puis participer au croquis collectif de synthèse et au débat final.
\tcblower
\textbf{Étape 1 — Inventorier et classer.} Le groupe 2 (climats), par exemple, extrait du tableau les pluies repères : Debundscha $\approx$ \SI{10000}{mm}, Yaoundé $\approx$ \SI{1600}{mm}, Ngaoundéré $\approx$ \SI{1500}{mm}, Maroua $\approx$ \SI{800}{mm}, Kousséri $<$ \SI{600}{mm}. Classées du plus humide au plus sec, ces valeurs dessinent un gradient sud-nord : c'est le critère de la légende. Le groupe 1 classe les altitudes : $> \SI{2000}{m}$ (dorsale camerounaise et sommets des hauts plateaux), \num{1000}--\SI{2000}{m} (hauts plateaux de l'Ouest, Adamaoua), \num{300}--\SI{1000}{m} (plateaux du Sud), $< \SI{300}{m}$ (plaines côtières et plaines du Nord).

\textbf{Étape 2 — Élaborer la légende.} Exemple pour le groupe 2 : titre de légende « Domaines climatiques » ; rubrique 1 : climat équatorial hyperhumide (côte) en vert foncé ; rubrique 2 : climat équatorial camerounien (Centre-Sud-Est) en vert clair ; rubrique 3 : climat tropical humide (Adamaoua) en jaune-vert ; rubrique 4 : climat tropical sec (Nord, Extrême-Nord) en orange. La hiérarchie suit le gradient d'humidité : cohérent et lisible.

\textbf{Étape 3 — Reporter sur le fond de carte.} Les limites sont tracées d'après l'atlas : la limite équatorial/tropical passe au niveau de l'Adamaoua ; la limite tropical humide/tropical sec passe au nord de Garoua. Les aplats sont posés en couleurs légères, les villes repères (Yaoundé, Ngaoundéré, Maroua, Kousséri) sont localisées par des points nommés, et les totaux pluviométriques sont inscrits en petits chiffres. Le nord est indiqué par une flèche, l'échelle par un trait gradué.

\textbf{Étape 4 — Titre et légende commentée.} Titre : « Les domaines climatiques du Cameroun : un gradient d'humidité du sud vers le nord ». Commentaire : « L'humidité décroît régulièrement de la côte (environ \SI{10000}{mm} à Debundscha) vers l'extrême nord (moins de \SI{600}{mm} à Kousséri). Le sud est équatorial, le centre tropical humide, le nord tropical sec : le Cameroun résume à lui seul les climats de l'Afrique intertropicale, d'où son surnom d'Afrique en miniature. »

\textbf{Étape 5 — Confrontation des quatre cartes.} La superposition mentale des quatre productions révèle des correspondances : la forêt dense occupe le domaine équatorial humide ; la savane couvre l'Adamaoua tropical ; la steppe suit le climat sec du nord ; la dorsale camerounaise perturbe ce zonage (prairies d'altitude, pluies orographiques abondantes sur le versant exposé aux vents océaniques, d'où les records de Debundscha). Le réseau hydrographique épouse le relief : la Sanaga draine le centre vers l'Atlantique, la Bénoué draine le nord-ouest vers le Niger, le Logone et le Chari drainent l'extrême nord vers le lac Tchad.

\textbf{Étape 6 — Croquis de synthèse et débat.} Le croquis collectif oppose un « Sud vert et humide » (forêt, climat équatorial, réseau atlantique) à un « Nord sec et ouvert » (savane, steppe, climat tropical, réseau tchadien), séparés par la « charnière adamaouaïenne ». Réponse argumentée attendue : la diversité est un handicap potentiel (contrastes de milieux, éloignement des extrémités, coûts de désenclavement), mais surtout un atout : complémentarités agricoles (cacao et café au Sud, coton et élevage au Nord), variété des ressources, image d'« Afrique en miniature » qui nourrit le sentiment d'unité nationale. L'intégration se construit précisément en reliant ces espaces complémentaires (routes, chemin de fer Transcamerounais, corridors).
\end{exoresolu}

\begin{popfigure}
\begin{center}
\includegraphics[width=\linewidth,height=9.5cm,keepaspectratio]{N01/cmr_physique.png}
\end{center}
\legende{Fond de carte : carte physique du Cameroun (altitudes, fleuves et limites d'État, sans noms de lieux). Le croquis de synthèse attendu situe sur ce fond : 1. la dorsale camerounaise à l'Ouest (teintes brunes) ; 2. le plateau de l'Adamaoua au centre ; 3. les plaines du Nord et le lac Tchad (vert) ; 4. les plateaux du Sud ; 5. les monts Mandara à l'Extrême-Nord ; puis les villes repères (Douala, Yaoundé, Ngaoundéré, Maroua) et la Sanaga. L'élève ajoute lui-même noms et figurés. \\ Source : Wikimedia Commons, Urutseg, CC0.}
\end{popfigure}

\begin{attention}[Erreurs fréquentes à éviter]
\begin{itemize}
\item Oublier le titre, l'orientation ou l'échelle : une carte sans ces éléments est incomplète, même bien coloriée.
\item Mettre dans la légende des éléments absents de la carte (ou l'inverse) : légende et carte doivent se correspondre exactement.
\item Utiliser le bleu pour autre chose que l'eau, ou des couleurs qui contredisent le message (vert pour la steppe, par exemple).
\item Confondre la Bénoué (qui rejoint le Niger au Nigeria) avec un affluent direct du lac Tchad : seuls le Logone et le Chari alimentent le lac Tchad.
\item Placer le mont Cameroun dans l'Extrême-Nord : il se situe au sud-ouest, sur la côte, dans la région du Sud-Ouest.
\item Écrire que la Sanaga est « le plus long fleuve d'Afrique centrale » : elle est seulement le plus long fleuve entièrement camerounais (environ \SI{918}{km}).
\end{itemize}
\end{attention}

\courssub{Bilan de l'activité}

Ce TP t'a permis de passer du statut de lecteur de cartes à celui de \cle{producteur de cartes}. En inventoriant, classant puis reportant les données, tu as vérifié concrètement que les quatre traits physiques ne sont pas indépendants : le relief commande les pluies, le climat commande la végétation, et l'hydrographie épouse la topographie. La confrontation des quatre cartes a fait émerger la grande opposition Nord sec / Sud humide et le rôle de charnière de l'Adamaoua. Enfin, le débat final t'a montré qu'un même fait géographique peut être lu comme un handicap ou comme un atout : tout dépend de l'organisation politique et économique de l'espace. C'est précisément le sens de la famille de situations « intégration nationale » : la diversité physique du Cameroun, loin de diviser, offre les complémentarités sur lesquelles l'unité du pays se construit.

\begin{exercice}[Pour aller plus loin]
Sur un fond de carte muet, construis seul(e) la carte « Les domaines climatiques du Cameroun » en 20 minutes, avec titre, légende hiérarchisée de 4 rubriques, orientation et échelle. Rédige ensuite un commentaire de 10 lignes montrant comment ce zonage climatique explique la répartition de la végétation.
\end{exercice}

\begin{corrige}[Pour aller plus loin]
La carte comporte quatre aplats ordonnés du sud vers le nord : vert foncé (équatorial hyperhumide, côte), vert clair (équatorial camerounien, Centre-Sud-Est), jaune-vert (tropical humide, Adamaoua), orange (tropical sec, Nord et Extrême-Nord). Le commentaire attendu suit trois étapes : (1) décrire le gradient d'humidité décroissant du sud (environ \SI{10000}{mm} à Debundscha) vers le nord (moins de \SI{600}{mm} à Kousséri) ; (2) établir la correspondance entre chaque domaine climatique et sa végétation (forêt dense humide au sud, savane arborée à l'Adamaoua, savane herbeuse puis steppe au nord) ; (3) conclure que la végétation est un indicateur fidèle du climat, la durée de la saison sèche étant le facteur décisif : plus elle est longue, plus la végétation s'éclaircit et se raréfie.
\end{corrige}

\newpage

\coursec{Méthodes et exercices résolus}

\coursec{La diversité physique du Cameroun}

\courssub{Le relief du Cameroun : une architecture en trois étages}

\begin{definition}[Unités morphostructurales]
Le relief camerounais s'organise en trois grandes unités superposées du nord au sud et de l'ouest à l'est, héritées de l'histoire géologique (socle précambrien, volcanisme tertiaire/quaternaire, bassins de sédimentation).
\end{definition}

\begin{exemplebox}[Les trois étages du relief]
\begin{itemize}
\item \textbf{Les plaines et basses terres ($< 300$ m) :} plaine côtière (Douala, Kribi, Tiko), plaine de la Bénoué (Garoua), plaine du Logone et bassin du lac Tchad (Maroua, Kousséri). Ce sont des zones d'accumulation sédimentaire (alluvions, colluvions) propices à l'agriculture et à l'urbanisation.
\item \textbf{Les plateaux ($300$--\num{1000} m) :} plateau sud-camerounais (Yaoundé, Sangmélima, Bertoua, \num{600}--\num{700} m), plateau de l'Adamaoua (Ngaoundéré, \num{1100} m, charnière climatique), plateau du Nord (autour de Garoua). Sur le socle ancien, l'altération a façonné des paysages de \cle{cuestas} (grands escarpments) et de \cle{plaines d'altération}.
\item \textbf{Les hautes terres ($> \num{1000}$ m) :} la \cle{dorsale camerounaise} (alignement volcanique SO-NE : mont Cameroun \num{4095} m, monts Bamboutos \num{2740} m, mont Oku \num{3011} m, mont Manengouba \num{2411} m) et les \cle{monts Mandara} (extrême-nord, granitiques, mont Oupay \num{1494} m). Zones de sources, de forêts de montagne, de cultures de haute altitude (thé, pomme de terre) et de tourisme.
\end{itemize}
\end{exemplebox}

\begin{popfigure}
\begin{center}
\includegraphics[width=\linewidth,height=9.5cm,keepaspectratio]{N01/cmr_physique.png}
\end{center}
\legende{Grandes unités de relief du Cameroun. Carte physique du Cameroun (altitudes en dégradé de vert à brun, légende en haut à gauche ; fleuves en bleu ; limites d'État en gris). On y lit : les hautes terres de l'Ouest (teintes beige-brun, plus de 1 000 m) et le mont Cameroun (tache brune isolée près de la côte) ; le vaste plateau central de l'Adamaoua et du Sud (750 à 1 000 m, teintes jaune pâle) ; les plaines basses du Nord et du bassin du lac Tchad (vert, moins de 400 m) ; la plaine littorale verte face au golfe de Guinée. \\ Source : Wikimedia Commons, Urutseg, CC0.}
\end{popfigure}

\begin{methode}[Décrire une photographie ou un paysage en 5 étapes]
La description est la première étape de tout commentaire de document géographique. Elle doit être \cle{rigoureuse}, \cle{spatialisée} et \cle{neutre} (on décrit avant d'expliquer).
\begin{enumerate}
\item \textbf{Identifier et localiser} : nommer le phénomène (ex. un massif montagneux), le situer (région, département, points cardinaux, repères connus : ville, fleuve, route).
\item \textbf{Donner l'échelle et l'orientation} : préciser s'il s'agit d'un paysage proche ou lointain, et orienter (avant-plan, arrière-plan ; nord, sud...).
\item \textbf{Décrire du général au particulier} : d'abord les grandes unités (plaine, plateau, montagne), puis les détails (formes de versants, sommets, vallées).
\item \textbf{Utiliser le vocabulaire géographique exact} : escarpement, piémont, cuesta, inselberg, cône volcanique, gradin, marigot... et des \cle{données chiffrées} (altitudes, pentes) quand elles sont disponibles.
\item \textbf{Conclure par une première hypothèse explicative} : relier le phénomène observé à une cause (ex. alignement de reliefs $\Rightarrow$ faille ou axe volcanique).
\end{enumerate}
\textbf{Réflexe} : toute affirmation doit être visible sur le document. Jamais d'information « par cœur » non vérifiable dans le document.
\end{methode}

\begin{exoresolu}[Décrire un paysage : les monts Mandara]
\textbf{Énoncé (type Bac).} Photographie : paysage de collines granitiques aux formes arrondies, parsemées de blocs rocheux, séparées par des plaines cultivées (mil, sorgho) ; villages aux cases rondes au pied des reliefs ; végétation clairsemée (quelques arbres, herbes sèches). Localisation : extrême-nord du Cameroun, région de l'Extrême-Nord, à la frontière du Nigeria. Décrivez ce paysage en suivant une démarche rigoureuse (5 à 8 lignes).
\tcblower
\textbf{Solution détaillée.}
\begin{enumerate}
\item \textbf{Localisation} : le paysage se situe dans les monts Mandara, à l'extrême-nord du Cameroun (région de l'Extrême-Nord), le long de la frontière nigériane, à l'ouest de la plaine de Mora.
\item \textbf{Échelle et orientation} : il s'agit d'une vue d'ensemble (moyenne distance) ; les reliefs occupent l'arrière-plan et les plaines cultivées l'avant-plan.
\item \textbf{Description du général au particulier} : l'unité dominante est un relief collinaire (altitudes moyennes de 500 à \num{1000} m, sommets culminant vers \num{1494} m au mont Oupay). Les collines présentent des formes arrondies en boules, caractéristiques des reliefs granitiques, avec de nombreux blocs rocheux (chaos) résultant de l'altération en boules. Entre les collines s'étendent des plaines de piémont mises en culture (mil, sorgho).
\item \textbf{Vocabulaire et éléments humains} : ce sont des \cle{inselbergs} (collines résiduelles) ; l'habitat groupé (cases rondes) occupe le pied des reliefs, signe d'une forte densité humaine et d'une mise en valeur intensive des terres.
\item \textbf{Hypothèse explicative} : les formes arrondies et les chaos de blocs traduisent une \cle{altération chimique} d'un socle granitique ancien sous climat tropical à saisons contrastées, suivie d'un dégagement des parties altérées par l'érosion.
\end{enumerate}
\textbf{Conclusion} : ce paysage des monts Mandara illustre un milieu de socle ancien modelé par l'altération et l'érosion, intensément occupé et cultivé par l'homme.
\end{exoresolu}

\courssub{Les climats du Cameroun : un gradient nord-sud commandé par la latitude et le relief}

\begin{definition}[Domaines climatiques]
Le Cameroun présente une succession zonale du sud au nord, modifiée par l'altitude (Adamaoua, dorsale) : domaine \cle{équatorial} (côtier et de transition), domaine \cle{tropical} (soudanien et sahélien).
\end{definition}

\begin{exemplebox}[Les quatre domaines climatiques]
\begin{center}
\begin{tabularx}{\linewidth}{|l|X|X|X|}
\hline
\rowcolor{popblueL} Domaine & Localisation type & Pluviométrie (mm/an) & Températures / Saisons \\
\hline
Équatorial camerounais (côtier) & Douala, Kribi, Limbé & \num{3000}--\num{4000} (jusqu'à \num{10000} au pied du Mt Cameroun) & $T^\circ$ moy. \num{26}--\num{27}$^{\circ}\mathrm{C}$, amplitude faible. 2 saisons des pluies (mars-juin, sept-nov), 2 saisons sèches courtes. \\
\hline
Équatorial de transition (guinéen) & Yaoundé, Bafia, Bertoua & \num{1500}--\num{2000} & $T^\circ$ moy. \num{24}--\num{25}$^{\circ}\mathrm{C}$. 4 saisons marquées (grande/s petite pluies, grande/petite sèches). \\
\hline
Tropical soudanien & Ngaoundéré, Garoua & \num{1000}--\num{1500} & $T^\circ$ moy. \num{26}--\num{28}$^{\circ}\mathrm{C}$, amplitude nette. 2 saisons : pluies (mai-oct), sèche (nov-avr). \\
\hline
Sahélien (tropical sec) & Maroua, Kousséri & $< \num{800}$ (Maroua \num{800}) & $T^\circ$ moy. \num{28}--\num{30}$^{\circ}\mathrm{C}$, amplitude forte. Saison sèche 7-8 mois (oct-avr), pluies courtes (juil-sept). \\
\hline
\end{tabularx}
\end{center}
\end{exemplebox}

\begin{popfigure}
\begin{center}
\includegraphics[width=\linewidth,height=9cm,keepaspectratio]{N01/cmr_koppen.png}
\end{center}
\legende{Carte des climats du Cameroun d'après la classification de Köppen-Geiger (légende en anglais à droite). On lit une succession du sud au nord : climats tropicaux humides (Af et Am, bleus) au Sud et sur le littoral, climat tropical à saison sèche de type savane (Aw, bleu clair) sur tout le centre et le nord du pays, puis climats arides (BSh orange, BWh rouge) à l'Extrême-Nord vers le lac Tchad. Les petites taches jaunes (Csb) signalent les hautes terres de l'Ouest, plus fraîches. \\ Source : Wikimedia Commons, Beck et al., CC BY 4.0.}
\end{popfigure}

\begin{methode}[Construire et lire un diagramme ombrothermique (Bagnouls-Gaussen)]
Le diagramme ombrothermique représente, mois par mois, les \cle{températures moyennes} (courbe rouge) et les \cle{précipitations} (histogramme bleu).
\begin{enumerate}
\item \textbf{Tracer les axes} : en abscisse, les 12 mois (J à D) ; à gauche, les précipitations $P$ (mm) ; à droite, les températures $T$ ($^{\circ}\mathrm{C}$) avec l'échelle conventionnelle $1^{\circ}\mathrm{C} = 2$ mm (échelle $1:2$).
\item \textbf{Construire l'histogramme des pluies} : une barre par mois, hauteur proportionnelle aux mm.
\item \textbf{Tracer la courbe des températures} : points mensuels reliés.
\item \textbf{Identifier les saisons sèches} : mois où la courbe $P$ passe \emph{sous} la courbe $2T$ (convention $P < 2T$ = saison sèche physiologique).
\item \textbf{Calculer les indicateurs} : $P_{\text{annuel}} = \sum P_i$ ; $T_{\text{moy}} = \frac{\sum T_i}{12}$ ; Amplitude $= T_{\max} - T_{\min}$.
\item \textbf{Conclure sur le type de climat} : nombre/longueur des saisons sèches, régime des pluies (un ou deux maxima).
\end{enumerate}
\end{methode}

\begin{exoresolu}[Exploiter des diagrammes ombrothermiques : Maroua et Douala]
\textbf{Énoncé (type Bac).} Données climatiques simplifiées (valeurs moyennes annuelles).

\begin{center}
\begin{tabularx}{\linewidth}{|l|X|X|X|X|}
\hline
\rowcolor{popblueL} Station & $P_{\text{annuel}}$ (mm) & $T_{\text{moy}}$ ($^{\circ}\mathrm{C}$) & Mois les plus pluvieux & Saison sèche ($P < 2T$) \\
\hline
Maroua & 800 & 28 & août & octobre à avril (7 mois) \\
\hline
Douala & \num{4000} & 26 & juillet--août & décembre à février (3 mois) \\
\hline
\end{tabularx}
\end{center}

1. Calculez la pluviométrie moyenne mensuelle. 2. Comparez les régimes pluviométriques. 3. Identifiez le domaine climatique.
\tcblower
\textbf{Solution détaillée.}
\begin{enumerate}
\item \textbf{Calculs.} $P_m = \frac{P_{\text{annuel}}}{12}$.
\begin{align*}
\text{Maroua : } P_m &= \frac{800}{12} \approx 66,7 \text{ mm/mois} \\
\text{Douala : } P_m &= \frac{4000}{12} \approx 333,3 \text{ mm/mois}
\end{align*}
Douala reçoit en moyenne 5 fois plus de pluie que Maroua ($\frac{4000}{800}=5$).
\item \textbf{Comparaison.} Maroua : régime \cle{tropical soudanien} à un seul maximum (août), saison sèche longue (7 mois), forte amplitude thermique. Douala : régime \cle{équatorial camerounais} (côtier), pluies abondantes quasi permanentes, deux maxima (juin et septembre), saison sèche courte (3 mois) et peu marquée ($P$ reste souvent $> 2T$), amplitude thermique faible.
\item \textbf{Domaines.} Maroua : domaine sahélien (Extrême-Nord). Douala : domaine équatorial camerounais (Littoral).
\end{enumerate}
\textbf{Conclusion} : le gradient climatique Nord-Sud est le facteur premier de la diversité physique camerounaise.
\end{exoresolu}

\begin{popfigure}
\begin{tikzpicture}
\begin{groupplot}[group style={group size=2 by 1, horizontal sep=1.5cm}, width=7cm, height=5cm,
    xmin=0.5, xmax=12.5, ymin=0, ymax=500,
    xtick={1,2,3,4,5,6,7,8,9,10,11,12},
    xticklabels={J,F,M,A,M,J,J,A,S,O,N,D},
    axis y line*=left, ylabel={Précipitations (mm)},
    axis y line*=right, ylabel={Température ($^{\circ}\mathrm{C}$)}, ymin=0, ymax=35,
    legend style={at={(0.5,-0.2)}, anchor=north, legend columns=2}]
% Maroua
\nextgroupplot[title={Maroua (Sahélien)}]
\addplot[ybar, popblue!60, fill=popblue!60] coordinates {(1,0)(2,0)(3,0)(4,5)(5,30)(6,80)(7,180)(8,220)(9,150)(10,40)(11,0)(12,0)};
\addplot[popred, thick, mark=*] coordinates {(1,26)(2,28)(3,31)(4,33)(5,32)(6,30)(7,28)(8,27)(9,27)(10,28)(11,27)(12,26)};
\addlegendentry{$P$ (mm)} \addlegendentry{$T$ ($^{\circ}\mathrm{C}$)}
% Douala
\nextgroupplot[title={Douala (Équatorial côtier)}, ymax=800]
\addplot[ybar, popblue!60, fill=popblue!60] coordinates {(1,30)(2,50)(3,150)(4,250)(5,300)(6,450)(7,650)(8,700)(9,600)(10,350)(11,150)(12,50)};
\addplot[popred, thick, mark=*] coordinates {(1,27)(2,28)(3,28)(4,27)(5,27)(6,26)(7,25)(8,25)(9,25)(10,26)(11,27)(12,27)};
\end{groupplot}
\end{tikzpicture}
\legende{Diagrammes ombrothermiques comparés (Bagnouls-Gaussen, échelle 1:2). Maroua : saison sèche nette 7 mois (barres sous la courbe 2T). Douala : pas de mois sec strict ($P > 2T$ toute l'année), deux maxima pluviométriques.}
\end{popfigure}

\courssub{La végétation du Cameroun : expression du climat et du relief}

\begin{definition}[Domaines phytogéographiques]
La végétation naturelle suit le gradient pluviométrique (climat) et l'altitude (relief). On distingue cinq grands ensembles, souvent dégradés par l'action humaine (agriculture, feu, urbanisation).
\end{definition}

\begin{exemplebox}[Les grands domaines de végétation]
\begin{itemize}
\item \textbf{Forêt dense humide sempervirente :} Sud, Littoral, Est, Sud-Ouest (au sud de la ligne Yaoundé--Bertoua). Essences : okoumé, ayous, sapelli, iroko. Biodiversité maximale. Exploitée par les compagnies forestières (ex. Pallisco, SEFAC).
\item \textbf{Savane arborée / forêt claire guinéenne :} Transition Centre (Yaoundé), Adamaoua (Ngaoundéré). Arbres espacés (daniellia, lophira) + strate herbacée. Zone de cultures pérennes (café, cacao, ananas).
\item \textbf{Savane herbeuse soudanienne :} Nord (Garoua, bassin Bénoué). Herbes hautes (Andropogon, Hyparrhenia), arbres isolés (karité, néré, fromager). Élevage transhumant, cultures cotonnières (SODECOTON).
\item \textbf{Steppe sahélienne (ocre) :} Extrême-Nord (Maroua, Logone). Herbes courtes, épineux (acacias, balanites), sols nus en saison sèche. Cultures de mil, sorgho, niébé ; pâturages.
\item \textbf{Mangrove :} Estuaires du Wouri (Douala), Rio del Rey (Limbé), Ndian. Palétuviers (Rhizophora). Protection côtière, nurseries halieutiques, bois de chauffe.
\end{itemize}
\end{exemplebox}

\begin{popfigure}
\begin{center}
\includegraphics[width=\linewidth,height=9.5cm,keepaspectratio]{N01/afr_veg.jpg}
\end{center}
\legende{Carte de la végétation naturelle de l'Afrique (légende en anglais), sur laquelle le Cameroun est situé au fond du golfe de Guinée. On y lit la zonation en bandes : forêt dense humide au sud du pays, savanes arborées puis herbeuses (verts plus clairs) au centre, steppe et semi-désert (rose) au nord vers le Tchad, et le désert plus au nord encore. Les taches violettes indiquent la végétation de montagne. À l'échelle du Cameroun, la végétation se lit comme une transition sud-nord qui suit la baisse des pluies. \\ Source : Wikimedia Commons, NCpedia (carte de la CIA), CC0.}
\end{popfigure}

\begin{methode}[Construire une légende de carte physique (TP 1)]
Une légende est la clé de lecture de la carte. Pour le TP 1 (carte physique : relief, climat, hydrographie ou végétation) :
\begin{enumerate}
\item \textbf{Analyser le sujet} : identifier le \cle{trait} demandé et les 3 à 5 catégories à représenter (ex. végétation : forêt dense, savane arborée, savane herbeuse, steppe, mangrove).
\item \textbf{Choisir le mode de représentation} : \cle{aplats de couleurs} (dégradé) pour zonages continus ; \cle{ponctuels} pour villes/sommets ; \cle{linéaires} pour fleuves/axes.
\item \textbf{Respecter les conventions} : dégradé clair$\to$foncé pour altitude/humidité croissante ; bleu pour l'eau ; vert pour végétation. Une couleur = une seule signification.
\item \textbf{Ordonner la légende} : ordre logique (altitudes croissantes, nord$\to$sud, humide$\to$sec).
\item \textbf{Compléter la carte} : titre précis (quoi + où), flèche du nord, échelle indicative, noms des repères (villes, sommets, fleuves).
\end{enumerate}
\textbf{Réflexe} : chaque symbole sur la carte figure dans la légende, et réciproquement.
\end{methode}

\begin{exoresolu}[TP guidé : Carte des domaines de végétation du Cameroun]
\textbf{Énoncé.} Construisez la carte « Les grands domaines de végétation du Cameroun » : fond de carte, 4 domaines + mangrove, légende complète.
\tcblower
\textbf{Solution détaillée.}
\begin{enumerate}
\item \textbf{Fond} : contour simplifié (triangle allongé), nord, échelle (1 cm $\approx$ 100 km), villes repères (Douala, Yaoundé, Ngaoundéré, Garoua, Maroua, Bamenda).
\item \textbf{Catégories (ordre sud$\to$nord)} :
\begin{itemize}
\item Forêt dense humide : Sud, Est, Littoral, Sud-Ouest.
\item Savane arborée / forêt claire : Centre (transition), Adamaoua.
\item Savane herbeuse soudanienne : Nord, Bénoué.
\item Steppe sahélienne (ocre) : Extrême-Nord, Logone, pourtour Lac Tchad.
\item Mangrove : estuaires Wouri, Rio del Rey, Ndian.
\end{itemize}
\item \textbf{Report des limites} : frontière forêt/savane $\approx$ parallèle de Yaoundé (transition) ; Adamaoua = savane arborée d'altitude ; Nord = savane herbeuse ; Extrême-Nord = steppe.
\item \textbf{Légende} : titre, items ordonnés, couleurs conventionnelles, symboles villes/sommets.
\item \textbf{Vérification} : Douala en forêt/mangrove ; Ngaoundéré en savane arborée ; Maroua en steppe.
\end{enumerate}
\textbf{Conclusion} : la carte montre un zonage climatique clair, le relief (Adamaoua, Dorsale) créant des enclaves (forêts de montagne, savanes d'altitude).
\end{exoresolu}

\courssub{L'hydrographie du Cameroun : deux bassins versants majeurs}

\begin{definition}[Réseau hydrographique]
Le réseau s'organise autour de deux bassins versants principaux séparés par la ligne de partage des eaux (Adamaoua, dorsale) : le \cle{bassin atlantique} (océan) et le \cle{bassin endoréique du lac Tchad}. La Bénoué relie le Cameroun au bassin du Niger.
\end{definition}

\begin{exemplebox}[Les grands cours d'eau et plans d'eau]
\begin{center}
\begin{tabularx}{\linewidth}{|l|X|X|X|}
\hline
\rowcolor{popblueL} Cours d'eau / Lac & Longueur (km) & Bassin & Caractéristiques / Aménagements \\
\hline
Sanaga & 890 (entièrement camerounais) & Atlantique & 1$^{\text{er}}$ fleuve national. Barrages : Édéa, Song Loulou, Lom Pangar (régulation). Hydroélectricité (90\% production). \\
\hline
Bénoué & \num{1400} ($\sim$700 au Cameroun) & Niger $\to$ Atlantique & Navigable (Garoua) en saison des pluies. Barrage de Lagdo (irrigation, électricité, pêche). Affluents : Mayo Rey, Faro. \\
\hline
Logone & \num{1000} (partagé Tchad) & Lac Tchad (endoréique) & Frontière Tchad. Inondations saisonnières (yaérés) = agriculture de décrue, pâturage. Affluent Chari. \\
\hline
Nyong & 690 & Atlantique & Régime équatorial, navigable par pirogues. \\
\hline
Wouri & 160 & Atlantique & Estuaire = Port autonome de Douala (95\% trafic maritime). \\
\hline
Lacs de barrage & - & - & Lagdo (Bénoué), Lom Pangar (Lom/Sanaga), Mbakaou (Djérem/Sanaga), Bamendjin (Noun/Sanaga). \\
\hline
Lacs de cratère & - & - & Lac Ossa (Sanaga Maritime), Barombi Mbo (Koumba), Nyos, Monoun (risque gaz CO$_2$). \\
\hline
\end{tabularx}
\end{center}
\end{exemplebox}

\begin{popfigure}
\begin{center}
\includegraphics[width=\linewidth,height=9.5cm,keepaspectratio]{N01/cmr_map.jpg}
\end{center}
\legende{Carte générale du Cameroun (relief ombré, légende en haut à gauche) montrant les fleuves en bleu : la Bénoué (Benue) qui traverse le Nord et rejoint le Nigeria vers le Niger ; le Logone et le Chari qui rejoignent le lac Tchad au nord ; la Sanaga qui descend vers le golfe de Guinée ; la Dja et la Ngoko au sud-est qui gagnent le bassin du Congo. Les retenues de Lagdo et de Mbakaou sont aussi visibles. Les trois versants (Atlantique, Congo, Tchad) se repèrent à la direction d'écoulement de ces cours d'eau. \\ Source : Wikimedia Commons, CIA, domaine public.}
\end{popfigure}

\begin{methode}[Exploiter un tableau statistique géographique]
\begin{enumerate}
\item \textbf{Lire le tableau} : titre, unités, source, date ; lignes (catégories), colonnes (variables).
\item \textbf{Classer} : ordre croissant/décroissant ; regrouper en classes.
\item \textbf{Calculer des indicateurs} : part $\% = \frac{\text{valeur}}{\text{total}} \times 100$ ; taux de variation $\% = \frac{V_f - V_i}{V_i} \times 100$.
\item \textbf{Repérer les faits saillants} : max, min, rapports (double, moitié).
\item \textbf{Interpréter géographiquement} : relier chiffres et réalités spatiales (localisation, facteurs).
\end{enumerate}
\textbf{Réflexe} : toujours citer les chiffres exacts du document avec leur unité.
\end{methode}

\begin{exoresolu}[Classer et interpréter : les fleuves du Cameroun]
\textbf{Énoncé (type Bac).} Tableau des longueurs approximatives.

\begin{center}
\begin{tabularx}{\linewidth}{|l|X|X|}
\hline
\rowcolor{popblueL} Cours d'eau & Longueur (km) & Bassin de destination \\
\hline
Sanaga & 890 & Atlantique \\
\hline
Bénoué & \num{1400} (dont $\sim$700 au Cameroun) & Niger $\to$ Golfe de Guinée \\
\hline
Wouri & 160 & Atlantique \\
\hline
Logone & \num{1000} (partagé) & Lac Tchad (endoréique) \\
\hline
Nyong & 690 & Atlantique \\
\hline
\end{tabularx}
\end{center}

1. Classez par longueur décroissante. 2. Part de la Sanaga dans le total des 5 cours d'eau (longueurs totales). 3. Deux directions d'écoulement et signification pour l'intégration.
\tcblower
\textbf{Solution détaillée.}
\begin{enumerate}
\item \textbf{Classement} : Bénoué (\num{1400}) $>$ Logone (\num{1000}) $>$ Sanaga (890) $>$ Nyong (690) $>$ Wouri (160).
\item \textbf{Part de la Sanaga.} Total $= 1400 + 1000 + 890 + 690 + 160 = \num{4140}$ km.
\[ \text{Part} = \frac{890}{4140} \times 100 \approx 21,5\,\% \]
La Sanaga représente $\approx$ 1/5 de la longueur cumulée.
\item \textbf{Interprétation.} Deux directions : (1) \cle{Bassin atlantique} (Sanaga, Nyong, Wouri, Bénoué via Niger) draine $\sim$ 80\% du territoire. (2) \cle{Bassin endoréique du Lac Tchad} (Logone/Chari) à l'Extrême-Nord. La Sanaga, plus long fleuve \emph{entièrement} camerounais, est la colonne vertébrale énergétique (barrages Édéa, Song Loulou, Lom Pangar) reliant le réservoir nord (Lom Pangar) aux centres de consommation sud (Douala, Yaoundé, Édéa). Le Wouri, bien que court, est vital : port de Douala = porte de l'économie nationale.
\end{enumerate}
\textbf{Conclusion} : l'hydrographie structure l'intégration (énergie, transport fluvial, irrigation, pêche) mais pose le défi de la gestion transfrontalière (Logone, Bénoué, Lac Tchad).
\end{exoresolu}

\courssub{Le milieu biogéographique : biodiversité et aires protégées}

\begin{definition}[Milieu biogéographique]
Ensemble formé par l'interaction du relief, du climat, de la végétation (flore) et de la faune. Le Cameroun, par sa position charnière et son gradient altitudinal, abrite une biodiversité exceptionnelle (\cle{hotspot} de la forêt du Congo et des montagnes du Cameroun).
\end{definition}

\begin{exemplebox}[Principaux écosystèmes et aires protégées emblématiques]
\begin{itemize}
\item \textbf{Forêt dense de plaine :} Parc national de \cle{Korup} (Sud-Ouest, forêt refuge pléistocène, primates), Réserve de la \cle{Dja} (Est, Patrimoine UNESCO, gorilles, chimpanzés, éléphants de forêt), Parc de \cle{Campo Ma'an} (Littoral, façade maritime).
\item \textbf{Savanes et forêts claires (Nord/Adamaoua) :} Parc national de \cle{Waza} (Extrême-Nord, zone inondable Logone, lions, éléphants, girafes, avifaune), Parc de \cle{Bénoué} et \cle{Bouba Ndjida} (Nord, grands mammifères savanicoles, cynégétique), Réserve de faune du \cle{Mbam et Djerem} (transition forêt-savane).
\item \textbf{Montagnes (Dorsale/Mandara) :} Forêts de montagne (Oku, Kilum-Ijim) : endémisme fort (oiseaux, amphibiens, plantes). Mont Cameroun : gradient continu mer-sommet unique en Afrique.
\item \textbf{Zones humides :} Lac Tchad (ramsar), Yaérés du Logone, Estuaire du Wouri (mangrove).
\end{itemize}
\end{exemplebox}

\begin{savaistu}[Le Cameroun, « Afrique en miniature » : un atout biodiversité]
Le Cameroun compte \num{9000} espèces de plantes vasculaires (dont \num{156} endémiques), \num{900} espèces d'oiseaux, \num{300} mammifères. Il possède \num{20} aires protégées majeures couvrant $\sim$ \num{10}\% du territoire (objectif Aichi/COP15 : 30\%). La conservation fait face au braconnage (ivoire, écailles de pangolin), à l'agro-industrie (palmier à huile, hévéa) et aux conflits homme-faune. Le corridor écologique Campo-Ma'an $\leftrightarrow$ Dja $\leftrightarrow$ Tri-national de la Sangha (TNS) est vital pour les grands mammifères.
\end{savaistu}

\courssub{Diversité physique et intégration nationale : synthèse}

\begin{propriete}[Complémentarités et contraintes]
La diversité physique offre des \cle{complémentarités} (échanges Nord-Sud : bétail/céréales vs produits forestiers/industriels ; énergie hydroélectrique Nord$\to$Sud ; port unique Douala/Kribi) mais impose des \cle{contraintes} : enclavement du Nord (routes difficiles saisonnières), risques naturels (coulées de lave Mt Cameroun, émanations CO$_2$ lacs Nyos/Monoun, inondations Logone, érosion côtière), fragmentation des espaces par les frontières.
\end{propriete}

\begin{methode}[Commenter un document géographique au Bac (carte, texte, photo, graphique)]
\begin{enumerate}
\item \textbf{Introduction} : nature/source/date/auteur du document ; localisation du sujet ; \cle{problématique} ; annonce du plan (2 ou 3 parties).
\item \textbf{Développement} : chaque partie = 1 idée force + \textbf{preuves documentaires} (chiffres, noms, localisations lus sur le doc) + \textbf{explication} (facteurs physiques/humains).
\item \textbf{Articulation} : décrire (ce que je vois) $\rightarrow$ expliquer (pourquoi c'est ainsi) $\rightarrow$ interpréter (quelles conséquences).
\item \textbf{Conclusion} : réponse synthétique à la problématique + ouverture (enjeu d'intégration, d'aménagement, de développement durable).
\end{enumerate}
\textbf{Réflexe} : zéro point pour la paraphrase pure. La copie doit \cle{dialoguer avec le document}.
\end{methode}

\begin{exoresolu}[Commenter un texte : « Le Cameroun, Afrique en miniature »]
\textbf{Énoncé (type Bac).} Texte : « De la plaine du lac Tchad aux forêts du Congo, des savanes de l'Adamaoua aux plages de Kribi, le Cameroun rassemble sur son territoire presque toutes les natures de l'Afrique : on l'appelle l'Afrique en miniature. » À l'aide de vos connaissances, commentez ce texte en montrant en quoi la diversité physique justifie cette appellation.
\tcblower
\textbf{Solution détaillée (plan rédigé).}

\textbf{Introduction.} Le document est un extrait de vulgarisation géographique qualifiant le Cameroun d'« Afrique en miniature ». Situé à la charnière Afrique centrale/Afrique de l'Ouest, entre golfe de Guinée ($4^{\circ}$ N) et lac Tchad ($13^{\circ}$ N), le Cameroun condense sur \num{475 442} km$^2$ une diversité exceptionnelle. Problématique : comment le relief, le climat, la végétation et l'hydrographie fondent-ils cette image ? Nous montrerons que le Cameroun réplique le grand gradient continental Nord-Sud, et que ses reliefs et réseaux hydrographiques parachèvent cette mosaïque.

\textbf{I. Un gradient climatique et végétal Nord-Sud, réplique accélérée du continent.} En moins de \num{1300} km, on observe la succession zonale africaine : au sud, domaine \cle{équatorial camerounais} (Douala, \num{4000} mm, forêt dense humide sempervirente : okoumé, ayous) ; au centre, domaine \cle{équatorial de transition} (Yaoundé, \num{1600} mm, savane arborée, cultures cacaoyères) ; au centre-nord, domaine \cle{tropical soudanien} (Ngaoundéré, Garoua, \num{1200} mm, savane herbeuse, karité, néré, coton) ; à l'extrême-nord, domaine \cle{sahélien} (Maroua, \num{800} mm, steppe épineuse, mil/sorgho). Ce zonage reproduit la séquence forêt--savane--steppe--désert qui s'étale sur \num{5000} km en Afrique.

\textbf{II. Relief et hydrographie : une mosaïque de formes et d'eaux structurante.} Le relief oppose les plaines périphériques (côtière, Bénoué, Logone, Tchad) aux plateaux centraux (sud-camerounais, Adamaoua) et à la \cle{dorsale camerounaise}, alignement volcanique SO-NE (Mt Cameroun \num{4095} m, plus haut sommet Afrique de l'Ouest ; Mt Oku \num{3011} m) qui barre la communication Est-Ouest. L'hydrographie s'organise en deux bassins versants majeurs : le \cle{bassin atlantique} (Sanaga, Nyong, Wouri, Bénoué via Niger) drainant le gros du territoire, et le \cle{bassin endoréique du Lac Tchad} (Logone/Chari) au Nord. Cette dualité alimente l'intégration : énergie (Sanaga : Édéa, Song Loulou, Lom Pangar), transport (port de Douala sur Wouri, navigation Bénoué), irrigation (Lagdo, yaérés Logone), pêche.

\textbf{Conclusion.} Par ses climats, végétations, reliefs et eaux, le Cameroun concentre les grands biomes africains : l'appellation « Afrique en miniature » est scientifiquement fondée. Cette diversité est un atout majeur (potentiel agricole, touristique, énergétique, minier) mais impose des politiques d'aménagement du territoire (Transcam, autoroute Yaoundé-Douala, port de Kribi, barrage Lom Pangar, corridor écologique) pour transformer l'hétérogénéité en cohésion nationale.

\end{exoresolu}

\courssub{TP 1 : Construire une carte physique du Cameroun (évaluation pratique)}

\begin{methode}[Démarche du TP cartographique (2 h)]
Le TP 1 porte sur un trait au choix : unités de relief, domaines climatiques, hydrographie ou végétation.
\begin{enumerate}
\item \textbf{Préparer le fond} (15 min) : reproduire le contour (polygone simplifié « en aile d'avion » ou « poule »), flèche du nord, échelle indicative (1 cm = 100 km), repères fixes (Douala, Yaoundé, Garoua, Maroua, Bamenda, Ngaoundéré, Bafoussam, Bertoua).
\item \textbf{Choisir et hiérarchiser l'information} (15 min) : 3 à 5 classes max. Ex. Végétation : Forêt dense / Savane arborée / Savane herbeuse / Steppe / Mangrove.
\item \textbf{Reporter les limites} (45 min) : tracer les frontières des zones d'après l'atlas/carte fournie. Respecter les grandes orientations (bandes Est-Ouest pour climat/végétation ; alignement SO-NE pour dorsale).
\item \textbf{Appliquer la légende} (30 min) : couleurs conventionnelles, ordre logique, pas de symbole orphelin. Vérifier correspondance carte/légende.
\item \textbf{Finaliser} (15 min) : titre complet (Thème + Lieu + Date), noms des repères, signature.
\end{enumerate}
\end{methode}

\begin{exoresolu}[TP guidé : Carte des unités de relief (exemple corrigé)]
\textbf{Énoncé.} Construire la carte « Les grandes unités de relief du Cameroun » (3 classes d'altitude + sommets + villes).
\tcblower
\textbf{Solution détaillée (attendus de correction).}
\begin{enumerate}
\item \textbf{Fond} : contour, nord, échelle, 7 villes repères.
\item \textbf{Classes (aplats dégradé)} :
\begin{itemize}
\item Vert clair : Plaines et basses terres $< 300$ m (côte, Bénoué, Logone, Lac Tchad).
\item Jaune ocre : Plateaux $300$--\num{1000} m (Sud-camerounais, Adamaoua).
\item Brun : Hautes terres $> \num{1000}$ m (Dorsale : Mt Cameroun, Bamboutos, Oku, Manengouba ; Mandara).
\end{itemize}
\item \textbf{Symboles ponctuels} : $\blacktriangle$ noir pour sommets (Mt Cameroun \num{4095} m, Mt Oku \num{3011} m, Mt Oupay \num{1494} m) ; $\bullet$ noir pour villes.
\item \textbf{Légende ordonnée} (bas $\to$ haut) : Plaines / Plateaux / Hautes terres / Sommets / Villes.
\item \textbf{Titre} : « Les grandes unités de relief du Cameroun ».
\end{enumerate}
\textbf{Critères de réussite (grille Bac) :} Fond correct (4 pts) + Légende complète/ordonnée (6 pts) + Report précis des limites (6 pts) + Esthétique/Propreté (4 pts) = 20 pts.
\end{exoresolu}

\begin{attention}[Les 5 erreurs qui coûtent des points au Bac]
\begin{enumerate}
\item \textbf{Production cartographique incomplète :} Carte ou croquis sans titre, sans légende, sans flèche du nord, sans échelle. \textbf{Sanction :} note maximale plafonnée à 10/20 voire 0 si absence totale de légende. \textit{Vérification systématique avant de rendre la copie.}
\item \textbf{Paraphrase au lieu d'exploitation :} Recopier les chiffres du tableau sans les classer, calculer (parts, taux) ou interpréter. \textit{Ex. : « La Sanaga fait 890 km » (0 pt) vs « La Sanaga (890 km), soit 21,5\% du total, est le plus long fleuve entièrement camerounais » (2 pts).}
\item \textbf{Confusions spatiales majeures :} Placer le Mt Oku plus haut que le Mt Cameroun ; situer la forêt dense au Nord ou la steppe au Sud ; rattacher le Logone à l'Atlantique (c'est le bassin endoréique Tchad) ; confondre Bénoué (Niger) et Sanaga (Atlantique). \textit{Apprendre les couples : Douala/Forêt-Mangrove, Ngaoundéré/Savane Adamaoua, Maroua/Steppe, Garoua/Savane Bénoué.}
\item \textbf{Erreurs de calcul « bêtes » :} Oublier $\times 100$ pour un pourcentage ; diviser par le mauvais total ; oublier l'unité (mm, $^{\circ}\mathrm{C}$, km, \%). \textit{Poser la formule : $\text{Part} = \frac{\text{Valeur}}{\text{Total}} \times 100$.}
\item \textbf{Dissociation description/explication :} Décrire sans localiser (pas de région, pas de point cardinal) ; expliquer sans s'appuyer sur le document (réciter le cours hors sujet). \textit{Structure obligatoire : Localisation $\to$ Description chiffrée $\to$ Explication causale.}
\end{enumerate}
\end{attention}

\newpage

\coursec{Exercices}

\courssub{Je vérifie mes connaissances}

\begin{exercice}[QCM : associer une ville à son milieu physique]
Pour chaque ville, choisis la bonne association \emph{unité de relief -- climat -- végétation dominante}. Une seule réponse est correcte par ligne.

\begin{center}
\begin{tabularx}{\linewidth}{|l|X|X|X|}
\hline
\rowcolor{popblueL} \textbf{Ville} & \textbf{Proposition A} & \textbf{Proposition B} & \textbf{Proposition C} \\
\hline
Douala & Plaine côtière ; climat équatorial humide ; forêt dense & Plateau de l'Adamaoua ; climat soudanien ; savane arborée & Hautes terres de l'Ouest ; climat de montagne ; prairies \\
\hline
Ngaoundéré & Plaine côtière ; climat équatorial ; mangrove & Plateau de l'Adamaoua ; climat soudanien humide ; savane arborée & Bassin du lac Tchad ; climat sahélien ; steppe \\
\hline
Maroua & Hautes terres de l'Ouest ; climat tropical humide ; forêt & Plateau de l'Adamaoua ; climat soudanien ; savane & Plaine du Diamaré ; climat soudano-sahélien sec ; savane herbeuse et steppe \\
\hline
Bafoussam & Hautes terres de l'Ouest ; climat tropical de montagne ; savane herbeuse d'altitude & Plaine côtière ; climat équatorial ; forêt & Bassin du lac Tchad ; climat sahélien ; steppe épineuse \\
\hline
\end{tabularx}
\end{center}

Justifie brièvement chacun de tes choix en une phrase.
\end{exercice}

\begin{exercice}[Vrai ou faux ? Justifie ta réponse]
Indique si chaque affirmation est vraie ou fausse, puis justifie en une ou deux phrases.
\begin{enumerate}
\item Le mont Cameroun (\SI{4100}{m} environ) est le point culminant de l'Afrique de l'Ouest et de l'Afrique centrale.
\item Le climat du Cameroun est uniformément équatorial sur l'ensemble du territoire.
\item Le plateau de l'Adamaoua est appelé le \og château d'eau \fg{} du Cameroun parce que de nombreux fleuves y prennent leur source.
\item La forêt dense humide couvre la majeure partie des régions de l'Extrême-Nord et du Nord.
\item Le fleuve Sanaga se jette dans l'océan Atlantique.
\end{enumerate}
\end{exercice}

\begin{exercice}[Définitions]
Définis précisément, en deux ou trois phrases chacune, les termes et expressions suivants :
\begin{enumerate}
\item domaine climatique ;
\item bassin versant ;
\item savane arborée ;
\item ligne de partage des eaux ;
\item milieu biogéographique.
\end{enumerate}
\end{exercice}

\begin{exercice}[Calculs climatiques directs]
On donne les données suivantes pour la station de Maroua (Extrême-Nord) : précipitations annuelles totales : \SI{800}{mm} environ ; mois le plus chaud (avril) : moyenne de \SI{33}{\celsius} ; mois le plus \og frais \fg{} (janvier) : moyenne de \SI{25}{\celsius}. Pour Yaoundé : précipitations annuelles : \SI{1600}{mm} environ ; température moyenne annuelle : \SI{23,5}{\celsius} ; amplitude thermique annuelle : \SI{2}{\celsius} environ.
\begin{enumerate}
\item Calcule l'amplitude thermique annuelle de Maroua.
\item Calcule le rapport entre les précipitations de Yaoundé et celles de Maroua. Exprime le résultat en une phrase : \og Yaoundé reçoit environ \ldots{} fois plus de pluie que Maroua \fg.
\item Compare les amplitudes thermiques des deux stations et propose une explication géographique à cet écart (latitude, altitude, continentalité).
\end{enumerate}
\end{exercice}

\courssub{Je m'entraîne}

\begin{exercice}[Construction de diagrammes climatiques]
Voici les moyennes mensuelles de deux stations camerounaises (T = température moyenne en \SI{}{\celsius} ; P = précipitations en \SI{}{mm}).

\begin{center}
\begin{tabularx}{\linewidth}{|l|X|X|X|X|X|X|X|X|X|X|X|X|}
\hline
\rowcolor{popblueL} \textbf{Mois} & J & F & M & A & M & J & J & A & S & O & N & D \\
\hline
T Yaoundé & 24 & 25 & 25 & 25 & 24 & 23 & 22 & 23 & 23 & 23 & 24 & 24 \\
\hline
P Yaoundé & 20 & 50 & 140 & 170 & 200 & 160 & 70 & 90 & 250 & 290 & 130 & 30 \\
\hline
T Maroua & 25 & 28 & 32 & 33 & 31 & 28 & 26 & 26 & 27 & 29 & 27 & 25 \\
\hline
P Maroua & 0 & 0 & 2 & 25 & 70 & 130 & 220 & 230 & 110 & 25 & 2 & 0 \\
\hline
\end{tabularx}
\end{center}

\begin{enumerate}
\item Construis sur ton cahier les deux diagrammes climatiques (températures en courbe, précipitations en barres ; échelle : \SI{1}{\celsius} $=$ \SI{2}{mm}).
\item Calcule pour chaque station le total annuel des précipitations et l'amplitude thermique annuelle.
\item Identifie le domaine climatique de chaque station (équatorial camerounien pour Yaoundé ; soudanien/soudano-sahélien pour Maroua) en justifiant par le régime des pluies (nombre de saisons, durée de la saison sèche).
\item Explique par deux facteurs géographiques (latitude, distance à la mer, altitude, alizés) les écarts observés entre les deux stations.
\end{enumerate}
\end{exercice}

\begin{exercice}[Localisation sur fond de carte muet]
Sur un fond de carte muet du Cameroun (reproduit à l'échelle de ton choix) :
\begin{enumerate}
\item Localise et nomme les cinq grandes unités de relief : la plaine côtière, les plateaux du Sud, le plateau de l'Adamaoua, les hautes terres de l'Ouest, les plaines du Nord (bénouéenne et tchadienne).
\item Place trois sommets repères : le mont Cameroun (\SI{4100}{m}), le mont Oku (\SI{3011}{m}), les monts Mandara.
\item Trace et nomme les trois grands bassins versants : Atlantique (Sanaga, Nyong, Wouri), Congo (Dja, Sangha), Tchad (Bénoué, Logone).
\item Rédige une légende conforme (un symbole = une information, titre de chaque élément de légende), ajoute le titre de la carte et la flèche du nord.
\end{enumerate}
\end{exercice}

\begin{exercice}[Classement et interprétation d'un tableau hydrographique]
On donne les caractéristiques approximatives de quelques cours d'eau camerounais :

\begin{center}
\begin{tabularx}{\linewidth}{|l|X|X|X|}
\hline
\rowcolor{popblueL} \textbf{Cours d'eau} & \textbf{Longueur (km)} & \textbf{Bassin versant} & \textbf{Débouché} \\
\hline
Sanaga & 890 & Atlantique & Océan Atlantique \\
\hline
Bénoué & 1400 (dont environ 700 au Cameroun) & Tchad & Fleuve Niger (Nigeria) \\
\hline
Dja & 720 & Congo & Rivière Sangha \\
\hline
Logone & 1000 (avec le Chari) & Tchad & Lac Tchad \\
\hline
Nyong & 690 & Atlantique & Océan Atlantique \\
\hline
\end{tabularx}
\end{center}

\begin{enumerate}
\item Classe ces cours d'eau selon leur bassin versant.
\item Quel est le cours d'eau le plus long entièrement ou majoritairement camerounais ? Justifie.
\item Explique pourquoi la Bénoué, bien que coulant au Cameroun, n'atteint ni l'océan Atlantique ni le lac Tchad directement.
\item Montre, à partir du tableau, que le Cameroun appartient à trois ensembles hydrographiques majeurs : quel atout cela représente-t-il pour l'irrigation et l'hydroélectricité ?
\end{enumerate}
\end{exercice}

\begin{exercice}[Croquis : le plateau de l'Adamaoua, \og château d'eau \fg{} du Cameroun]
Construis un croquis du réseau hydrographique camerounais centré sur le plateau de l'Adamaoua.
\begin{enumerate}
\item Esquisse le contour simplifié du Cameroun et hachure ou colorie la zone du plateau de l'Adamaoua (altitude moyenne : \SI{1000}{m} environ).
\item Représente par des flèches bleues les cours d'eau qui naissent sur le plateau et rayonnent vers : le Nord (Bénoué, Vina, Mbéré), le Sud (Sanaga et ses affluents : Lom, Djérem), l'Est (Kadeï, Lom, vers le bassin du Congo).
\item Indique par un symbole le barrage hydroélectrique de Lagdo (sur la Bénoué) et celui de Mbakaou (sur le Djérem).
\item Respecte les impératifs cartographiques : titre, légende ordonnée, flèche du nord, échelle indicative.
\item Rédige un court paragraphe (5 à 8 lignes) expliquant pourquoi ce plateau mérite le nom de \og château d'eau \fg.
\end{enumerate}
\end{exercice}

\courssub{Je me prépare au Bac}

\begin{exercice}[Analyse de documents : deux milieux contrastés sous menace]
\textbf{Document 1} (photographie décrite) : forêt dense de Korup (Sud-Ouest). Végétation luxuriante à plusieurs étages, arbres géants à fûts droits dépassant \SI{40}{m}, lianes, sous-bois sombre et humide ; plus de \SI{3000}{mm} de pluie par an ; le parc national de Korup (environ \SI{1260}{km\squared}) abrite une biodiversité exceptionnelle (environ un quart des espèces de primates d'Afrique).

\textbf{Document 2} (photographie décrite) : savane de Waza (Extrême-Nord). Étendue herbeuse jaunie parsemée d'acacias épineux ; moins de \SI{700}{mm} de pluie concentrés sur 4 mois ; le parc national de Waza (environ \SI{1700}{km\squared}) abrite éléphants, girafes, lions et antilopes.

\textbf{Document 3} (texte) : \og Le lac Tchad, qui couvrait environ \SI{25000}{km\squared} dans les années 1960, n'en couvre plus que \SI{1500}{} à \SI{2500}{km\squared} en saison sèche selon les années. La baisse des pluies, les prélèvements pour l'irrigation et la pression démographique expliquent cette régression qui menace pêcheurs, éleveurs et cultivateurs riverains. \fg

\begin{enumerate}
\item Décris méthodiquement chacun des deux milieux photographiés (végétation, climat, faune) et identifie le domaine biogéographique correspondant.
\item Calcule la perte de superficie du lac Tchad entre les années 1960 et aujourd'hui, en valeur absolue et en pourcentage (prends \SI{2000}{km\squared} comme valeur actuelle moyenne).
\item À partir des trois documents, dégage deux menaces naturelles et deux menaces humaines qui pèsent sur les milieux physiques camerounais.
\item Propose trois mesures concrètes de conservation applicables au Cameroun, en précisant pour chacune le milieu concerné.
\end{enumerate}
\end{exercice}

\begin{exercice}[Étude de cas : l'axe Douala--Ngaoundéré, une traversée des milieux camerounais]
Un transporteur emprunte le corridor Douala--Ngaoundéré (environ \SI{900}{km} par la route, environ \SI{620}{km} par le chemin de fer Transcamerounais). Le tableau ci-dessous résume les milieux traversés.

\begin{center}
\begin{tabularx}{\linewidth}{|l|X|X|X|}
\hline
\rowcolor{popblueL} \textbf{Étape} & \textbf{Relief} & \textbf{Climat} & \textbf{Végétation} \\
\hline
Douala & Plaine côtière ($<$ \SI{200}{m}) & Équatorial humide, 2 saisons de pluies & Forêt dense, mangrove \\
\hline
Yaoundé & Plateau du Sud (700--\SI{800}{m}) & Équatorial de transition, 4 saisons & Forêt dégradée, mosaïque forêt-savane \\
\hline
Ngaoundéré & Plateau de l'Adamaoua (\SI{1100}{m}) & Soudanien humide, 1 saison de pluies & Savane arborée \\
\hline
\end{tabularx}
\end{center}

\begin{enumerate}
\item Présente le document (nature, objet, échelle du phénomène étudié).
\item Décris les transformations du paysage observées par le voyageur de Douala à Ngaoundéré (relief, climat, végétation).
\item Explique le rôle de l'altitude et de la latitude dans l'évolution des climats le long de cet axe.
\item Montre comment cet axe illustre la complémentarité des milieux camerounais (produits échangés entre le Sud forestier et le Nord pastoral : bois, cacao, café contre bétail, coton, oignons).
\item Dégage deux contraintes physiques que la diversité des milieux impose à la construction et à l'entretien de ce corridor (relief, pluies, ensablement des voies au Nord).
\end{enumerate}
\end{exercice}

\begin{exercice}[Question de synthèse : dissertation géographique]
\textbf{Sujet :} \og La diversité physique du Cameroun : atout ou obstacle à l'intégration nationale ? \fg

Tu rédigeras un développement structuré et argumenté comportant :
\begin{enumerate}
\item une introduction avec une définition des termes clés (diversité physique, intégration nationale), une présentation du cadre spatial et une problématique clairement formulée ;
\item une première partie montrant que la diversité physique est un \emph{atout} : complémentarités agricoles (cacao et café du Sud, coton et bétail du Nord), variété des paysages touristiques (plages de Kribi, parcs de Waza et de Korup, mont Cameroun), potentialités hydroélectriques (Sanaga, Bénoué), image d'\og Afrique en miniature \fg ;
\item une deuxième partie montrant qu'elle peut être un \emph{obstacle} : difficultés et coûts de désenclavement (routes coupées par les pluies, relief accidenté de l'Ouest), inégalités de peuplement et de mise en valeur entre Sud et Nord, vulnérabilités climatiques (sécheresses au Nord, inondations à Douala) ;
\item une conclusion équilibrée ouvrant sur les politiques d'aménagement qui transforment les contraintes en atouts (grands travaux : routes, barrages, pont sur le Wouri).
\end{enumerate}
Chaque argument devra être illustré par au moins un exemple précis et localisé (région, ville, fleuve ou infrastructure nommés).
\end{exercice}



\newpage

\coursec{Corrigés des exercices}

\coursec{La diversité physique du Cameroun}

\courssub{Je vérifie mes connaissances}

\begin{corrige}[Exercice 1 — QCM : associer une ville à son milieu physique]
\begin{itemize}
\item \textbf{Douala : proposition A.} Douala est située dans la plaine côtière, au bord de l'estuaire du Wouri ; elle connaît un climat équatorial de type camerounais (deux saisons de pluies, deux saisons sèches) avec des précipitations annuelles supérieures à \SI{3000}{mm} sur le littoral proche (Debundscha) et une végétation de forêt dense humide, avec mangrove sur le littoral.
\item \textbf{Ngaoundéré : proposition B.} Chef-lieu de la région de l'Adamaoua, la ville repose sur le plateau du même nom (altitude moyenne \SI{1100}{m}), sous climat soudanien humide à une seule saison de pluies, avec une végétation de savane arborée.
\item \textbf{Maroua : proposition C.} Maroua s'étend dans la plaine du Diamaré (Extrême-Nord), sous climat soudano-sahélien sec (environ \SI{800}{mm} concentrés sur 3 à 4 mois) ; la végétation est une savane herbeuse à steppe épineuse (acacias, balanites).
\item \textbf{Bafoussam : proposition A.} Chef-lieu de la région de l'Ouest, Bafoussam occupe les hautes terres de l'Ouest (environ \SI{1500}{m}) ; le climat tropical est rafraîchi par l'altitude (températures moyennes \SI{20}{\celsius}) et la forêt a laissé place à des savanes herbeuses d'altitude très anthropisées (cultures, pâturages).
\end{itemize}
\textbf{Point de vigilance :} ne pas confondre les hautes terres de l'Ouest (Bafoussam, volcanisme récent, sols ferralitiques lessivés, climat de montagne) avec le plateau de l'Adamaoua (Ngaoundéré, socle ancien, sols ferrugineux tropicaux, climat soudanien) : les deux sont élevées, mais le climat, la végétation et les sols diffèrent.
\end{corrige}

\begin{corrige}[Exercice 2 — Vrai ou faux ?]
\begin{enumerate}
\item \textbf{Vrai.} Avec \SI{4095}{m} (dernières mesures GPS), le mont Cameroun, volcan encore actif (dernière éruption 2000), est le sommet le plus élevé d'Afrique de l'Ouest et d'Afrique centrale (le Kilimandjaro étant en Afrique de l'Est).
\item \textbf{Faux.} Le climat est équatorial seulement dans le Sud (domaine équatorial guinéen et camerounais) ; vers le Nord on passe au climat soudanien (Adamaoua, Nord) puis au climat soudano-sahélien sec (Extrême-Nord), sans oublier les climats de montagne des hautes terres de l'Ouest (altitude > \SI{1500}{m}).
\item \textbf{Vrai.} Par son altitude moyenne (environ \SI{1000}{m}) et sa position centrale, le plateau de l'Adamaoua alimente des cours d'eau qui rayonnent vers le Nord (Bénoué et affluents : Vina, Mbéré), le Sud (Sanaga et affluents : Lom, Djérem) et l'Est (Kadeï, Boumba, affluents du bassin du Congo via la Sangha).
\item \textbf{Faux.} L'Extrême-Nord et le Nord sont des domaines de savane (savane arborée au sud, savane herbeuse, puis steppe épineuse au nord) ; la forêt dense humide occupe le Sud et le littoral du pays (régions du Sud, du Centre, de l'Est, du Littoral et du Sud-Ouest).
\item \textbf{Vrai.} La Sanaga (\SI{918}{km}), plus long fleuve \emph{entièrement} camerounais, se jette dans le golfe de Guinée (océan Atlantique) au sud d'Édéa, après avoir traversé la retenue d'Édéa et les chutes de la Nachtigal.
\end{enumerate}
\end{corrige}

\begin{corrige}[Exercice 3 — Définitions]
\begin{enumerate}
\item \textbf{Domaine climatique :} vaste portion de territoire caractérisée par un même type de climat, défini par la combinaison des températures (moyennes, amplitudes) et des précipitations (totaux annuels, régime des pluies, durée de la saison sèche). Au Cameroun, on distingue du Sud au Nord les domaines équatorial (guinéen et camerounais), soudanien et soudano-sahélien, ainsi que le domaine de montagne.
\item \textbf{Bassin versant :} surface drainée par un cours d'eau principal et ses affluents ; toutes les eaux de pluie tombées sur cette surface convergent vers le même exutoire (fleuve, lac ou océan). Il est délimité par des lignes de partage des eaux.
\item \textbf{Savane arborée :} formation végétale tropicale faite d'une strate herbeuse continue (graminées hautes) parsemée d'arbres et d'arbustes assez nombreux (karités \emph{Vitellaria paradoxa}, nérés \emph{Parkia biglobosa}, baobabs \emph{Adansonia digitata} au Nord camerounais) ; elle correspond à un climat soudanien à saison sèche marquée (5 à 7 mois).
\item \textbf{Ligne de partage des eaux :} ligne de crête qui sépare deux bassins versants voisins ; de part et d'autre, les eaux s'écoulent vers des exutoires différents. Le plateau de l'Adamaoua porte une ligne de partage majeure entre les bassins de l'Atlantique (direct et via Niger), du Congo et du lac Tchad.
\item \textbf{Milieu biogéographique :} ensemble formé par les conditions physiques d'un lieu (relief, climat, sols, eaux) et par les êtres vivants qui y sont adaptés (végétation, faune), en interaction permanente. La forêt de Korup (réserve de biosphère, forte endémicité) et la savane de Waza (parc national, grande faune sahélienne) sont deux milieux biogéographiques contrastés du Cameroun.
\end{enumerate}
\end{corrige}

\begin{corrige}[Exercice 4 — Calculs climatiques directs]
\begin{enumerate}
\item Amplitude thermique annuelle de Maroua (mois le plus chaud : avril \SI{33}{\celsius} ; mois le plus frais : janvier \SI{25}{\celsius}) :
\[ 33 - 25 = 8 \]
L'amplitude thermique annuelle de Maroua est de \SI{8}{\celsius}.
\item Rapport des précipitations annuelles (Yaoundé \SI{1600}{mm} / Maroua \SI{800}{mm}) :
\[ \frac{1600}{800} = 2 \]
Yaoundé reçoit environ \textbf{deux fois} plus de pluie que Maroua.
\item L'amplitude de Maroua (\SI{8}{\celsius}) est près de trois fois plus forte que celle de Yaoundé (\SI{3}{\celsius}). Yaoundé, proche de l'équateur (3°52' N) et à environ \SI{750}{m} d'altitude, connaît des températures régulières toute l'année (influence maritime, nébulosité). Maroua, à plus haute latitude (10°35' N), plus éloignée de la mer (continentalité) et soumise à l'harmattan sec et chaud en saison sèche, subit de forts écarts thermiques entre janvier (frais) et avril (très chaud).
\end{enumerate}
\textbf{Point de vigilance :} l'amplitude thermique se calcule toujours entre le mois le plus chaud et le mois le plus frais (moyennes mensuelles), et non entre deux mois quelconques ; le rapport de deux grandeurs de même unité n'a pas d'unité.
\end{corrige}

\courssub{Je m'entraîne}

\begin{corrige}[Exercice 5 — Construction de diagrammes climatiques (méthode de Gaussen)]
\begin{enumerate}
\item Le diagramme ombrothermique se construit avec les mois en abscisse, les températures moyennes mensuelles en ordonnée de gauche (courbe rouge, trait plein) et les précipitations totales mensuelles en ordonnée de droite (barres bleues), en respectant l'échelle \textbf{\SI{1}{\celsius} = \SI{2}{mm}} afin de comparer directement courbe et barres (période sèche = courbe au-dessus des barres).
\item \textbf{Yaoundé :} total annuel $= 20+50+140+170+200+160+70+90+250+290+130+30 = \SI{1600}{mm}$ ; amplitude thermique $= 25,0 - 22,0 = \SI{3,0}{\celsius}$ (valeurs moyennes arrondies). \textbf{Maroua :} total annuel $= 0+0+2+25+70+130+220+230+110+25+2+0 = \SI{814}{mm}$, soit environ \SI{800}{mm} ; amplitude thermique $= 33,0 - 25,0 = \SI{8,0}{\celsius}$.
\item \textbf{Yaoundé :} climat équatorial camerounais à quatre saisons (deux saisons de pluies : mars--juin et septembre--novembre ; deux saisons sèches : décembre--février et juillet--août) ; il pleut au moins un peu tous les mois. \textbf{Maroua :} climat soudano-sahélien à une seule saison de pluies (juin--septembre, maximum en août) et une longue saison sèche de 7 à 8 mois quasi sans pluie (octobre--mai).
\item \textbf{Latitude :} Yaoundé, proche de l'équateur, reste sous l'influence de la mousson humide (alizés de sud-ouest) presque toute l'année ; Maroua, vers 10° N, n'est atteinte par la mousson qu'en été boréal (juin--sept). \textbf{Distance à la mer et alizés :} Maroua, très éloignée de l'océan (\SI{>1000}{km}), subit en saison sèche l'harmattan (alizé continental sec et chaud venu du Sahara), d'où des pluies rares, une forte insolation et une forte amplitude thermique.
\end{enumerate}
\textbf{Point de vigilance :} vérifier la somme des douze valeurs mensuelles avant de conclure ; une erreur de report d'un seul mois fausse le total annuel et donc l'identification du domaine climatique. Respecter l'échelle 1/2 (température double des précipitations) pour lire la saison sèche.
\end{corrige}

\begin{corrige}[Exercice 6 — Localisation sur fond de carte muet : les grandes unités physiques]
La carte attendue doit présenter :
\begin{itemize}
\item \textbf{Relief :} la plaine côtière en bande étroite le long du golfe de Guinée ; les plateaux du Sud (Yaoundé, Ebolowa, Sangmélima) au centre-Sud (alt. 600--\SI{800}{m}) ; le plateau de l'Adamaoua en position centrale (alt. \SI{1000}{--1100}{m}) ; les hautes terres de l'Ouest en alignement SW-NE du mont Cameroun vers les Grassfields (mont Oku, \SI{3011}{m}) ; les plaines du Nord (plaine de la Bénoué autour de Garoua, plaine du Diamaré et bassin du lac Tchad autour de Maroua et Kousséri). Les monts Mandara se placent à l'extrême nord-ouest de l'Extrême-Nord, à la frontière nigériane.
\item \textbf{Hydrographie :} les fleuves côtiers (Wouri, Sanaga, Nyong) coulent vers l'Atlantique ; le Dja et la Kadeï (affluents de la Sangha) s'écoulent vers l'est (bassin du Congo) ; la Bénoué et ses affluents (Vina, Mbéré) drainent le Nord vers le Niger (Atlantique via Nigeria) ; le Logone et le Chari drainent l'Extrême-Nord vers le lac Tchad.
\item \textbf{Légende :} ordonnée (relief, hydrographie, sommets, villes repères), avec un titre explicite : \og Les grandes unités physiques du Cameroun \fg, une flèche du nord et une échelle indicative.
\end{itemize}
\textbf{Point de vigilance :} un symbole de la légende ne doit figurer qu'une seule fois et correspondre exactement à ce qui est reporté sur la carte ; la flèche du nord, le titre et l'échelle sont obligatoires, leur absence coûte des points au Bac.
\end{corrige}

\begin{corrige}[Exercice 7 — Classement et interprétation d'un tableau hydrographique]
\begin{enumerate}
\item \textbf{Classification par exutoire :}
\begin{itemize}
\item Bassin de l'Atlantique (direct) : Sanaga, Nyong, Wouri.
\item Bassin du Niger (Atlantique via Nigeria) : Bénoué.
\item Bassin du Congo (via Sangha) : Dja, Kadeï, Boumba.
\item Bassin du Lac Tchad (endoréique) : Logone, Chari.
\end{itemize}
\item La \textbf{Sanaga} (\SI{918}{km}) est le plus long fleuve \emph{entièrement} camerounais : elle naît sur le plateau de l'Adamaoua (confluence Lom/Djérem) et se jette dans l'Atlantique sans quitter le territoire national. La Bénoué est plus longue au total (\SI{1400}{km}) mais coule pour moitié au Nigeria.
\item La Bénoué s'écoule vers l'ouest, franchit la frontière nigériane à Garoua et rejoint le fleuve Niger à Lokoja : ses eaux atteignent l'Atlantique par le delta du Niger, au Nigeria, et non par le littoral camerounais. Elle n'atteint pas le lac Tchad, car elle appartient au sous-bassin Niger, séparé du bassin du Chari-Logone (lac Tchad) par des lignes de partage des eaux sur l'Adamaoua.
\item Le tableau montre \textbf{quatre exutoires distincts} : océan Atlantique direct (Sanaga, Nyong), fleuve Niger (Bénoué), fleuve Congo (via Sangha), lac Tchad (via Logone/Chari). Cette répartition offre de l'eau dans toutes les régions : elle permet l'irrigation des plaines du Nord (riziculture de la SEMRY à Yagoua sur le Logone, riziculture de la SODECOTON sur la Bénoué) et la production hydroélectrique sur les chutes et rapides des fleuves du Sud (barrages d'Édéa et de Song Loulou sur la Sanaga, de Lagdo sur la Bénoué, de Memve'ele sur le Ntem).
\end{enumerate}
\end{corrige}

\begin{corrige}[Exercice 8 — Croquis : le plateau de l'Adamaoua, « château d'eau »]
Le croquis attendu montre le plateau de l'Adamaoua en position centrale, colorié ou hachuré (alt. \SI{1000}{--1100}{m}), d'où partent des flèches bleues en étoile (épaisseur proportionnelle au débit) :
\begin{itemize}
\item Vers le nord : la Bénoué et ses affluents (Vina, Mbéré) avec le symbole du barrage hydroélectrique de Lagdo (\SI{72}{MW}).
\item Vers le sud : la Sanaga (confluence Lom + Djérem) avec le symbole du barrage de Mbakaou (régulation) et les barrages aval (Song Loulou, Nachtigal, Édéa).
\item Vers l'est : la Kadeï (affluent de la Sangha/Congo).
\end{itemize}
La légende comporte trois rubriques ordonnées (unité de relief, cours d'eau et sens d'écoulement, ouvrages hydroélectriques), le titre \og Le plateau de l'Adamaoua, château d'eau du Cameroun \fg, la flèche du nord et une échelle indicative.

\textbf{Paragraphe explicatif attendu :} le plateau de l'Adamaoua mérite le nom de \og château d'eau \fg{} parce que son altitude moyenne (environ \SI{1000}{m}) et sa position centrale en font une ligne de partage des eaux majeure. Les pluies soudaniennes (\SI{1500}{--1800}{mm/an}) qui s'y déversent alimentent des cours d'eau qui rayonnent dans toutes les directions : la Bénoué vers le Nord (bassin Niger), la Sanaga et ses affluents vers le Sud (Atlantique direct), la Kadeï vers l'Est (bassin Congo). Ce plateau approvisionne ainsi les trois grands bassins versants du pays. Ses dénivelés (chutes sur les rebords) ont permis la construction de barrages hydroélectriques (Lagdo, Mbakaou, Song Loulou) qui fournissent l'essentiel de l'électricité nationale et régulent les crues. Il conditionne donc l'eau, l'énergie et l'agriculture d'une grande partie du Cameroun.

\textbf{Point de vigilance :} sur un croquis, les flèches doivent indiquer le sens d'écoulement (de la source vers l'exutoire), jamais l'inverse ; chaque flèche doit être nommée ; la légende doit être structurée par thèmes.
\end{corrige}

\courssub{Je me prépare au Bac}

\begin{corrige}[Exercice 9 — Analyse de documents : deux milieux contrastés sous menace (Korup / Waza -- Lac Tchad)]
\begin{enumerate}
\item \textbf{Document 1 (Forêt de Korup) :} forêt dense sempervirente à plusieurs étages (strates), arbres de plus de \SI{40}{m}, lianes abondantes, sous-bois humide ; climat équatorial guinéen très pluvieux (plus de \SI{3000}{mm/an}, pas de mois sec) ; faune forestière riche (primates : drills, chimpanzés ; éléphants de forêt). Il s'agit du \textbf{domaine forestier équatorial guinéen} (Sud-Ouest). \textbf{Document 2 (Parc de Waza) :} savane herbeuse à acacias épineux (\emph{Acacia seyal}, \emph{Balanites aegyptiaca}) ; climat soudano-sahélien sec (moins de \SI{700}{mm} sur 3--4 mois, juin--sept) ; faune de savane (éléphants, girafes, lions, kob, damalisques). Il s'agit du \textbf{domaine soudanien à sahélien} (Extrême-Nord).
\item Perte en valeur absolue : $25000 - 2000 = \SI{23000}{\square\km}$. Perte en pourcentage (par rapport à l'état initial) :
\[ \frac{23000}{25000} \times 100 = 92 \]
Le lac Tchad a perdu environ \SI{23000}{\square\km}, soit \textbf{92 \%} de sa superficie des années 1960 (variabilité selon les années, mais effondrement majeur).
\item \textbf{Menaces naturelles :} la baisse des pluies et la sécheresse récurrente depuis les années 1970 (régression du lac Tchad, document 3) ; la forte saison sèche qui fragilise la savane de Waza (incendies naturels, stress hydrique). \textbf{Menaces humaines :} les prélèvements d'eau excessifs pour l'irrigation (riz, coton) et la pression démographique sur les rives (document 3) ; la déforestation (exploitation forestière, agriculture sur brûlis) et le braconnage qui pèsent sur la forêt de Korup (document 1).
\item Trois mesures possibles localisées : (1) la création et le renforcement d'aires protégées intégrales (parc national de Korup pour la forêt, parc national de Waza pour la savane, avec corridors écologiques) ; (2) la gestion concertée et durable de l'eau du bassin du lac Tchad (irrigation maîtrisée, restauration des polders, coopération dans le cadre de la Commission du bassin du lac Tchad - CBLT) ; (3) le reboisement et l'agroforesterie (parcs à karités, ceintures vertes) dans les zones de savane de l'Extrême-Nord pour lutter contre la désertification et l'ensablement.
\end{enumerate}
\textbf{Barème indicatif (sur 20) :} description méthodique des deux milieux (vocabulaire précis, chiffres) : 6 pts ; calculs exacts et rédigés (formule, résultat, unité) : 4 pts ; quatre menaces distinguées (naturelles / humaines, localisées) : 5 pts ; trois mesures localisées et réalistes : 5 pts.

\textbf{Points de vigilance :} le pourcentage de perte se calcule par rapport à la superficie \emph{initiale} (\SI{25000}{\square\km}), non à la superficie actuelle ; chaque affirmation doit être appuyée sur un document précis (\og d'après le document 3, la superficie passe de...\fg) ; citer les noms des parcs et organisations (CBLT).
\end{corrige}

\begin{corrige}[Exercice 10 — Étude de cas : l'axe Douala--Ngaoundéré (corridor routier et ferroviaire Transcamerounais)]
\begin{enumerate}
\item \textbf{Présentation :} il s'agit d'un tableau statistique et descriptif présentant les caractéristiques physiques (relief, climat, végétation) de trois étapes du corridor Douala--Ngaoundéré ; le phénomène est étudié à l'échelle nationale, le long d'un axe sud-nord d'environ \SI{700}{km} (route nationale N1 et voie ferrée Transcamerounais).
\item \textbf{Description :} de Douala à Ngaoundéré, le relief s'élève progressivement : plaine côtière basse (alt. < \SI{200}{m}), puis plateaux du Sud (alt. 700--\SI{800}{m} autour de Yaoundé), puis plateau de l'Adamaoua (alt. \SI{1100}{m} à Ngaoundéré). Le climat passe de l'équatorial camerounais (deux saisons de pluies, \SI{>2500}{mm}) à l'équatorial de transition / guinéen (quatre saisons, \SI{1600}{mm} à Yaoundé), puis au soudanien humide (une seule saison de pluies, \SI{1500}{mm} à Ngaoundéré). La végétation évolue de la forêt dense humide et de la mangrove (Douala) à la forêt dégradée en mosaïque cultures/forêt (plateaux du Sud), puis à la savane arborée (Adamaoua).
\item \textbf{Explication :} la \emph{latitude} croissante (3° N $\to$ 7° N) éloigne progressivement de l'équateur : la mousson humide atlantique influence le Nord pendant une période de plus en plus courte, d'où le passage de deux saisons de pluies (Douala) à une seule (Ngaoundéré). L'\emph{altitude} croissante (0 $\to$ \SI{1100}{m}) rafraîchit les températures (gradient adiabatique) et accentue les contrastes saisonniers sur l'Adamaoua (amplitude thermique plus forte).
\item \textbf{Complémentarité :} le Sud forestier et humide fournit bois (grumes, sciages), cacao (région du Centre, Sud), café arabica (Ouest), bananes (Littoral, Njombé), exportés par le port autonome de Douala ; le Nord pastoral et soudanien envoie vers le Sud bétail (Adamaoua, Nord), coton (Extrême-Nord, filière SODECOTON), oignons (plaine du Diamaré, Maroua). Le corridor (route bitumée N1 et Transcamerounais) permet ces échanges entre milieux aux productions différentes : la diversité physique devient ainsi un facteur d'intégration économique nationale.
\item \textbf{Contraintes :} les fortes pluies équatoriales dégradent les routes (ornières, glissements de terrain sur les pentes des plateaux) et provoquent des coupures de pistes au Sud ; le relief accidenté (montées vers les plateaux du Sud et de l'Adamaoua) rend la construction et l'entretien coûteux (ouvrages d'art, terrassements) ; au Nord, la saison sèche et les vents de sable (harmattan) ensablent les voies et imposent un entretien permanent (désensablement).
\end{enumerate}
\textbf{Barème indicatif (sur 20) :} présentation du document (nature, objet, échelle, source) : 2 pts ; description ordonnée (relief, climat, végétation) : 5 pts ; explication par deux facteurs nommés (latitude, altitude) : 4 pts ; complémentarité illustrée par des produits précis et des infrastructures : 5 pts ; deux contraintes localisées et expliquées : 4 pts.

\textbf{Point de vigilance :} la présentation d'un document (nature, objet, échelle, date, source) est une étape obligatoire et notée ; ne jamais la confondre avec la description du contenu.
\end{corrige}

\begin{corrige}[Exercice 11 — Question de synthèse : dissertation géographique]
\textbf{Sujet :} « La diversité physique du Cameroun : atout ou obstacle à l'intégration nationale ? »

\textbf{Introduction.} La diversité physique désigne la variété des reliefs, climats, végétations et réseaux hydrographiques d'un territoire ; l'intégration nationale est le processus qui unit économiquement, politiquement et culturellement les différentes parties d'un État. Le Cameroun, pays charnière entre l'Afrique centrale et l'Afrique de l'Ouest, s'étend du golfe de Guinée (2° N) au lac Tchad (13° N) et présente une mosaïque de milieux qui lui vaut le surnom d'\og Afrique en miniature \fg. On peut alors se demander : cette diversité physique favorise-t-elle ou freine-t-elle l'unité du pays ?

\textbf{Première partie : un atout majeur pour la complémentarité et le développement.}
La variété des milieux fonde des complémentarités agricoles et économiques fortes : le Sud forestier et humide produit cacao (régions du Centre, Sud, Est), café arabica (hautes terres de l'Ouest, Dschang, Bafoussam) et robusta (Littoral, Sud-Ouest), bananes d'exportation (plantations CDC/PAMOL à Tiko, Njombé, Mungo), bois d'œuvre (forêts du Sud, Est) ; le Nord soudanien et sahélien fournit coton (Extrême-Nord, filière SODECOTON, Maroua, Kaele), bétail (Adamaoua, Nord, cheptel national \textasciitilde 6 millions de têtes), oignons (plaine du Diamaré, Maroua, export vers l'Afrique centrale) ; l'Adamaoua offre des potentialités laitères et céréalières (maïs, sorgho). Ces échanges internes (flux Nord-Sud) tissent la solidarité économique nationale. La diversité des paysages nourrit un tourisme varié : plages de Kribi et de Limbé, ascension du mont Cameroun (\SI{4095}{m}), safaris à Waza, Bénoué, Bouba Ndjida, forêt primaire de Korup (biodiversité). Les fleuves aux régimes variés offrent un potentiel hydroélectrique considérable (barrages d'Édéa \SI{263}{MW}, Song Loulou \SI{384}{MW}, Nachtigal \SI{420}{MW} sur la Sanaga ; Lagdo \SI{72}{MW} sur la Bénoué ; Memve'ele sur le Ntem), source d'énergie pour l'industrie (aluminium Edea, régie de distribution ENEO) et les villes.

\textbf{Deuxième partie : des contraintes qui freinent l'aménagement et creusent les inégalités.}
Le relief accidenté des hautes terres de l'Ouest (pentes fortes, sols instables) et les pluies diluviennes du littoral (\SI{>3000}{mm}, érosion, lessivage) rendent la construction et l'entretien des routes et voies ferrées très coûteux (coût au km multiplié par 3 à 5 en zone de montagne) ; certaines pistes du Nord sont coupées en saison des pluies (inondations des plaines alluviales du Logone, de la Bénoué), d'autres ensablées en saison sèche (harmattan sur la N1 Extrême-Nord). Les milieux inégalement favorables ont produit un peuplement et une mise en valeur inégaux : le Sud, mieux desservi, plus arrosé, moins malade (paludisme moins endémique qu'au Nord ? Non, paludisme partout, mais sommeil...), concentre villes (Douala, Yaoundé), industries (70\% du PIB), administrations et infrastructures ; l'Extrême-Nord souffre de sécheresses récurrentes, d'insécurité alimentaire, de pression foncière et d'insécurité (Boko Haram), le Nord et l'Adamaoua d'enclavement relatif. La diversité peut ainsi creuser les inégalités régionales (indicateur de développement humain inégal) et alimenter des tensions (revendications autonomistes, conflits agro-pastoraux).

\textbf{Conclusion.} La diversité physique du Cameroun est ambivalente : contrainte lorsqu'elle isole et oppose les régions (enclavement, coûts d'infrastructures, risques climatiques), elle devient richesse lorsque l'aménagement la met en valeur (complémentarités, énergie, tourisme). Les grands travaux récents (bitumage des axes structurants N1, N2, N10, pont sur le Wouri à Douala, barrages hydroélectriques de Nachtigal et Memve'ele, extension du Transcamerounais vers Ngaoundéré et projet vers le Tchad) montrent que la volonté politique, appuyée sur la planification spatiale (SNAT, PDN), peut transformer les contraintes naturelles en fondements durables de l'intégration nationale.

\textbf{Barème indicatif (sur 20) :} introduction (définitions, cadre spatial, problématique) : 4 pts ; première partie argumentée et illustrée (exemples nommés : régions, villes, fleuves, entreprises) : 5 pts ; deuxième partie argumentée et illustrée : 5 pts ; conclusion équilibrée avec ouverture (aménagement, prospective) : 3 pts ; cohérence, langue, exemples localisés, structure binaire : 3 pts.

\textbf{Points de vigilance :} chaque argument doit être accompagné d'un exemple nommé (région, ville, fleuve, infrastructure, entreprise) ; éviter le plan catalogue (simple liste de milieux : relief, climat, végétation...) ; la problématique doit être une véritable question, annoncée dès l'introduction et reprise en conclusion ; utiliser le vocabulaire géographique (exutoire, ligne de partage, complémentarité, enclavement, aménagement du territoire).
\end{corrige}

\newpage

\coursec{Fiche bilan}

\begin{popfigure}
\begin{tikzpicture}[mindmap, grow cyclic, every node/.style={concept, text=white, font=\small\bfseries},
root concept/.append style={concept color=popink, font=\bfseries},
level 1/.append style={level distance=3.6cm, sibling angle=60},
level 2/.append style={level distance=2.2cm, sibling angle=40, font=\scriptsize}]
\node[root concept]{Diversité physique du Cameroun}
  child[concept color=poporange]{ node {Relief}
    child { node {Dorsale camerounaise} }
    child { node {Hauts plateaux} }
    child { node {Plaines et basses terres} }
  }
  child[concept color=popblue]{ node {Climats}
    child { node {Équatorial (4 sous-types)} }
    child { node {Tropical soudanien} }
    child { node {Tropical sahélien} }
  }
  child[concept color=popgreen]{ node {Végétation}
    child { node {Forêt dense (S)} }
    child { node {Savanes (centre)} }
    child { node {Steppes (N)} }
  }
  child[concept color=popgold]{ node {Hydrographie}
    child { node {Bassin du Congo} }
    child { node {Bassin de l'Atlantique} }
    child { node {Bassin du Niger--Tchad} }
  }
  child[concept color=poppurple]{ node {Milieu biogéographique}
    child { node {Sols : ferrallitiques, ferrugineux} }
    child { node {Faune et aires protégées} }
  }
  child[concept color=poppink]{ node {Intégration nationale}
    child { node {Complémentarités} }
    child { node {Contraintes} }
  };
\end{tikzpicture}
\legende{Carte mentale de la leçon : les six grandes entrées de la diversité physique du Cameroun et leurs contenus essentiels à retenir.}
\end{popfigure}

\begin{bilan}
\textbf{Définitions clés}
\begin{itemize}
  \item \cle{Dorsale camerounaise} : alignement volcanique et horstéen orienté SO--NE, du mont Cameroun (\SI{4095}{m}, point culminant) aux monts Mandara ; ligne de partage des eaux majeure.
  \item \cle{Climat équatorial} : climat chaud ($\approx$ \SI{25}{\celsius}) et humide toute l'année (4 sous-types : guinéen, camerounien, congolais, submontagnard).
  \item \cle{Climat tropical} : climat chaud à saison sèche marquée ; soudanien au centre, sahélien à l'Extrême-Nord.
  \item \cle{Bassin versant} : surface drainée par un fleuve et ses affluents ; le Cameroun en compte trois (Congo, Atlantique, Niger--Tchad).
  \item \cle{Milieu biogéographique} : ensemble formé par le sol, la végétation et la faune, en interaction avec le climat.
\end{itemize}

\textbf{Repères chiffrés à connaître par c\oe ur}
\begin{center}
\begin{tabularx}{\linewidth}{|l|X|}\hline
\rowcolor{popblueL} \textbf{Élément} & \textbf{Donnée essentielle} \\ \hline
Point culminant & Mont Cameroun, \SI{4095}{m} \\ \hline
Précipitations record & Debundscha : plus de \SI{10000}{mm/an} (parmi les lieux les plus arrosés du monde) \\ \hline
Grand fleuve & Sanaga (\SI{918}{km} environ, entièrement camerounais) \\ \hline
Frontière bioclimatique clé & Adamaoua : charnière entre Sud forestier et Nord soudano-sahélien \\ \hline
Grands ensembles végétaux & Forêt dense (sud), savanes (centre), steppes (Extrême-Nord) \\ \hline
\end{tabularx}
\end{center}

\textbf{Méthodes express}
\begin{itemize}
  \item \textbf{Décrire un relief} : localiser $\rightarrow$ citer les unités $\rightarrow$ donner altitudes et formes $\rightarrow$ expliquer (tectonique, volcanisme, érosion).
  \item \textbf{Lire un climat} : températures (moyennes, amplitude) $+$ pluies (total, régime : 1 ou 2 saisons sèches) $\rightarrow$ nommer le type.
  \item \textbf{Construire un croquis} : fond de carte $\rightarrow$ tracés $\rightarrow$ légende numérotée et ordonnée $\rightarrow$ titre $+$ orientation $+$ échelle indicative.
\end{itemize}
\end{bilan}

\begin{aretenir}[Checklist avant le Bac]
\begin{itemize}
  \item $\square$ Je sais citer et localiser les grandes unités de relief (dorsale camerounaise, hauts plateaux de l'Ouest et de l'Adamaoua, plaines du Nord, bas plateau et plaine littorale du Sud).
  \item $\square$ Je connais le point culminant du pays et l'origine volcanique de la dorsale.
  \item $\square$ Je distingue les trois grands types de climats et les quatre sous-types du climat équatorial.
  \item $\square$ Je sais expliquer le rôle de la dorsale et de la mousson dans la répartition des pluies.
  \item $\square$ Je sais associer chaque domaine climatique à sa végétation (forêt, savane, steppe).
  \item $\square$ Je connais les trois bassins versants et les principaux fleuves (Sanaga, Wouri, Nyong, Bénoué, Logone, Dja).
  \item $\square$ Je sais définir le milieu biogéographique et citer les grands types de sols.
  \item $\square$ Je sais relier diversité physique et intégration nationale (complémentarités agricoles, contraintes de communication).
  \item $\square$ Je maîtrise la construction d'un croquis légendé (TP 1) : titre, légende ordonnée, nord, échelle indicative.
  \item $\square$ Je donne les chiffres avec prudence : ordre de grandeur signalé quand la donnée exacte est incertaine.
\end{itemize}
\end{aretenir}

\findecours

\end{document}
